% ===========================================================================
%  17 — Registro de decisoes arquiteturais
% ===========================================================================
\chapter{Registro de decisões arquiteturais}
\label{cap:adr}

Cada decisão abaixo segue a forma de um ADR~\cite{nygard2011adr}: contexto,
decisão, alternativas descartadas, consequências --- inclusive as ruins --- e o
gatilho que faria reavaliá-la. Decisão sem alternativa descartada e sem gatilho de
reavaliação não é decisão registrada: é justificativa.

São \textbf{73 decisões}, numeradas até ADR-73, \textbf{sem lacunas}. A ADR-65
ficou vaga por um tempo: o número se perdeu numa colisão entre duas frentes
paralelas que reservaram o mesmo identificador --- o risco que a \adrr{ADR-61}
registra --- e a lacuna foi ocupada pelo \textbf{guia da aula} (\adrr{ADR-65}),
em vez de a numeração seguir de 71 para 72.

A colisão voltou a acontecer no mesmo dia, e em dobro: \textbf{três} frentes
paralelas fecharam decisão em 2026-09-06, e duas escreveram ADR-65 sem se ver.
Prevaleceu a mais antiga; as outras seguiram para \adrr{ADR-72} e
\adrr{ADR-73}. É a razão de a \adrr{ADR-61} mandar conferir \cd{origin/dev} na
hora de escrever o número --- e a prova de que conferir só a branch integrada
não basta: o número precisa ser reservado contra as frentes ABERTAS.

Renumerar as seis ADRs entre 66 e 71 para fechar a lacuna, na época, foi
considerado e recusado, e a razão não é o trabalho: dentro deste repositório a
troca é mecânica, porque os \cd{\textbackslash label} são nominais e o
\cd{\textbackslash adrr} é formatação, não referência resolvida. O problema é o
que sobra depois. Renumerar não deixa referência QUEBRADA --- deixa referência
\textbf{errada}: \enquote{ADR-69} continuaria resolvendo, para outra decisão, em
silêncio, em onze mensagens de \emph{commit} imutáveis, nos PRs já mergeados e em
qualquer cópia deste PDF que já tenha circulado. É a forma exata do
\cd{expires\_at} que cinco ADRs deste capítulo registram como armadilha. Um
identificador vale por nunca mudar de significado; a ordem cronológica é o que se
abre mão --- e é por isso que a \adrr{ADR-65} não é, em data, a mais antiga das
que a cercam.

\section{Índice}

\begin{tabularx}{\linewidth}{@{}L{16mm}X L{18mm} L{20mm}@{}}
\toprule
\headrow \thd{ID} & \thd{Decisão} & \thd{Data} & \thd{Estado}\\
\midrule
\adrr{ADR-01} & Experiências como \emph{bundles} estáticos, não serviços & 2026-08 & Aceita\\
\zebra \adrr{ADR-02} & \emph{Monorepo} único (pnpm + Turborepo) & 2026-08 & Aceita\\
\adrr{ADR-03} & Autenticação por OIDC genérico & 2026-08 & Aceita\\
\zebra \adrr{ADR-04} & Uma \emph{stack} por contrato; escolas por RLS & 2026-08 & Aceita\\
\adrr{ADR-05} & Cloud SQL HA, não AlloyDB & 2026-08 & Aceita\\
\zebra \adrr{ADR-06} & Pub/Sub para telemetria desde o início & 2026-08 & Aceita\\
\adrr{ADR-07} & IA atrás de um \emph{gateway} na própria API & 2026-08 & Aceita\\
\zebra \adrr{ADR-08} & RLS forçada e uma porta única de banco & 2026-08-24 & Aceita\\
\adrr{ADR-09} & Domínio por experiência na API; gabarito \emph{server-side} & 2026-08-24 & Aceita\\
\zebra \adrr{ADR-10} & Bem-estar: DASS-21 fixo, IA só no entorno, sem visão de gestor & 2026-08-25 & Aceita\\
\adrr{ADR-11} & Telemetria sem PII e vocabulário fechado no \emph{payload} & 2026-08-26 & Aceita\\
\zebra \adrr{ADR-12} & \cd{grau} como invariante de API, não regra de tela & 2026-08-26 & Aceita\\
\adrr{ADR-13} & Retenção sem prazo declarado, em troca de reprocessamento & 2026-08-26 & Aceita\\
\zebra \adrr{ADR-14} & Telemetria no PostgreSQL particionado; sem BigQuery, sem Redis & 2026-08-27 & Aceita\\
\adrr{ADR-15} & O hub assina os \emph{tokens} enquanto o IdP central não entra & 2026-08 & Aceita, transitória\\
\zebra \adrr{ADR-16} & Credencial por ADC; modelo \cd{3.6-flash} em \cd{global} & 2026-08 & Aceita, com custo\\
\adrr{ADR-17} & \emph{Claim} atômico em vez de fila externa & 2026-08 & Aceita\\
\zebra \adrr{ADR-18} & Exceções registradas: jogo em JS clássico, exp-70 fora do DS & 2026-08-26 & Aceita\\
\adrr{ADR-19} & \emph{Design system} externo como fonte única de interface & 2026-08 & Aceita\\
\zebra \adrr{ADR-20} & Um banco, um \emph{schema}; prefixo \cd{exp\_NN\_} & 2026-08-27 & Aceita\\
\adrr{ADR-21} & \emph{Push subscription} para a própria API, sem \emph{worker} & 2026-08 & Aceita\\
\zebra \adrr{ADR-22} & exp-02 1.1.0: \emph{briefing} na voz da escola, conclusão por botão & 2026-08-27 & Aceita, com custo\\
\adrr{ADR-23} & \emph{Function calling} entra no \cd{AiProvider}, não na experiência & 2026-08-27 & Aceita\\
\zebra \adrr{ADR-24} & Teto de custo de IA por SESSÃO; a quota diária é fusível & 2026-08-27 & Aceita, com número a medir\\
\adrr{ADR-25} & O Relógio da Aula é tabela da experiência, não coluna em \cd{turmas} & 2026-08-27 & Aceita\\
\zebra \adrr{ADR-26} & Sem prazo de retenção na exp-03; exclusão por \enquote{zerar a turma} & 2026-08-27 & Aceita\\
\adrr{ADR-27} & Transcrição do estudante é dado do instrumento: professor lê, gestor não & 2026-08-27 & Aceita\\
\bottomrule
\end{tabularx}

% ---------------------------------------------------------------------------
\section{ADR-01 --- Experiências como \emph{bundles} estáticos}

\begin{especificacao}{Contexto}
O catálogo alvo é de 48 experiências, cada uma uma aplicação React independente.
O pico de carga é uma parede às oito da manhã, e o grosso dos \emph{bytes} é o
código e a mídia dessas aplicações. Uma pessoa opera tudo.
\end{especificacao}

\textbf{Decisão.} Cada experiência é um \emph{bundle} Vite 100\% estático,
publicado em \arq{gs://\{contrato\}-exp/\{expId\}/\{versao\}/} e servido pela
CDN. Não há processo, porta, nem \emph{scaling} por experiência. O hub monta a URL
do \cd{iframe} com \cd{versaoAtiva} do manifesto.

\textbf{Alternativas descartadas.}
\begin{itemize}
  \item \emph{Microsserviço por experiência}: 48 serviços com \emph{deploy}, log,
    escala e \emph{alerting} próprios --- para servir arquivos estáticos.
  \item \emph{Rotas do hub monolítico}: cada experiência entraria no \emph{bundle}
    do hub, e o hub cresceria sem limite. Uma mudança em qualquer experiência
    exigiria \emph{deploy} do hub inteiro.
\end{itemize}

\textbf{Consequências.}
\begin{itemize}
  \item Custo marginal de publicação $\approx$ armazenamento + tráfego de CDN.
  \item O caminho carrega a versão \implica{} o \emph{bundle} é imutável e pode
    ficar um ano na borda.
  \item \textbf{\emph{Rollback} é um \cd{UPDATE}}: trocar \cd{versaoAtiva}. Sem
    \emph{build}, sem imagem, sem \emph{deploy}. É a operação de emergência mais
    rápida da plataforma.
  \item \emph{Ruim}: a experiência não tem \emph{backend} próprio, então toda
    lógica de domínio dela vive na API --- o que concentra responsabilidade num
    serviço só (mitigado por \adrr{ADR-09}).
  \item \emph{Ruim}: um \emph{bundle} publicado não é reconstruído quando a API
    muda, o que torna os contratos do Capítulo~\ref{cap:contratos} rígidos.
\end{itemize}

\textbf{Gatilho de reavaliação.} Uma experiência que precise de execução no
servidor por natureza --- multiplayer em tempo real com autoridade de servidor, ou
processamento pesado por sessão.

% ---------------------------------------------------------------------------
\section{ADR-02 --- \emph{Monorepo} único}

\textbf{Contexto.} A SDK é consumida por 48 aplicações. O \emph{design system} é
externo, mas a ponte é interna.

\textbf{Decisão.} Um \emph{monorepo} com pnpm \emph{workspaces} e
Turborepo~\cite{turborepo2026}. \cd{apps/*}, \cd{apps/experiencias/*},
\cd{packages/*}.

\textbf{Alternativa descartada.} Repositório por experiência, com a SDK publicada
em npm. Uma mudança na SDK só seria testada contra as 48 aplicações depois de
publicada --- e a incompatibilidade apareceria num \emph{bundle} servido a estudantes
semanas depois.

\textbf{Consequências.} Mudança na SDK testa tudo no mesmo \emph{pull request};
o CI builda só o que mudou (\cd{--filter='...[HEAD\^{}]'}). \emph{Ruim}: o
repositório cresce, e \cd{pnpm install} completo é lento.

\textbf{Gatilho.} Estúdio externo produzindo experiência --- aí a experiência sai
para repositório próprio e passa a consumir a SDK publicada.

% ---------------------------------------------------------------------------
\section{ADR-03 --- Autenticação por OIDC genérico}

\textbf{Contexto.} O WSO2 Identity Server é exigência do cliente e da estratégia
multi-contrato: é único para várias redes.

\textbf{Decisão.} A aplicação \textbf{não conhece o WSO2}: conhece OIDC. A API
valida \emph{token} por \env{OIDC\_ISSUER}, \env{OIDC\_JWKS\_URI} e
\env{OIDC\_AUDIENCE}. Trocar de provedor de identidade é trocar essas
variáveis.

\textbf{Alternativa descartada.} SDK do WSO2 na aplicação. Amarraria o código a um
produto que pode mudar por decisão comercial, e tornaria o teste local dependente
de um servidor Java em \emph{bootstrap}.

\textbf{Consequências.} \env{OIDC\_ISSUER} e \env{OIDC\_JWKS\_URI} vão como
\emph{env} em texto plano (não são segredo: são URLs públicas), derivadas de
\cd{hub\_domain} --- e é isso que garante que o \cd{iss} do \emph{token} e o
esperado pela API \textbf{não podem} divergir.

\textbf{Gatilho.} Nenhum previsto. É a decisão mais estável do conjunto.

% ---------------------------------------------------------------------------
\section{ADR-04 --- Uma \emph{stack} por contrato; escolas por RLS}

\textbf{Contexto.} Cada contrato é uma relação jurídica com um ente federativo;
cada contrato tem 100+ escolas.

\textbf{Decisão.} Isolamento \textbf{físico} entre contratos (\emph{stack} GCP
inteira, por \emph{module} Terraform) e \textbf{lógico} entre escolas
(\cd{escola\_id} + RLS).

\textbf{Alternativas descartadas.}
\begin{itemize}
  \item \emph{Uma \emph{stack} por escola}: cem infraestruturas por contrato,
    operadas por uma pessoa. E o custo fixo de uma escola pequena passaria a ser o
    de uma rede.
  \item \emph{Uma \emph{stack} para todos os contratos}: um incidente num
    município seria um incidente em todos, e o desenho não sobreviveria à primeira
    análise jurídica de um edital.
\end{itemize}

\textbf{Consequências.} \cd{escola\_id} em toda tabela de dado de estudante, RLS
forçada, e \cd{admin-contrato} como única exceção cross-escola. \emph{Ruim}: cada
contrato novo é uma \emph{stack} nova para provisionar, monitorar e faturar.

\textbf{Gatilho.} Exigência contratual de isolamento físico \emph{por escola} ---
o que mudaria o modelo comercial, não só a arquitetura.

% ---------------------------------------------------------------------------
\section{ADR-05 --- Cloud SQL HA, não AlloyDB}

\textbf{Decisão.} PostgreSQL gerenciado (Cloud SQL) com HA, PITR e \emph{pooler}
de conexões obrigatório.

\textbf{Alternativa descartada.} AlloyDB. \textbf{Motivo:} a robustez que este
sistema precisa vem de HA + PITR + \emph{pooler}, não do \emph{engine}. Ambos são
PostgreSQL, então migrar depois é trivial se a carga justificar.

\textbf{Consequências.} O \emph{pooler} é \textbf{obrigatório}: o
\emph{autoscaling} do Cloud Run multiplica \emph{pools} do Prisma e estouraria
\cd{max\_connections}. E o banco só tem IP privado --- o que obriga a migração a
rodar como Cloud Run Job dentro da VPC.

\textbf{Gatilho.} Carga medida que sature uma instância de Cloud SQL com
\emph{pooler} bem dimensionado.

% ---------------------------------------------------------------------------
\section{ADR-06 --- Pub/Sub para telemetria desde o início}

\textbf{Contexto.} No pico, quatro mil estudantes simultâneos gerando evento a cada
poucos segundos.

\textbf{Decisão.} O canal \emph{analítico} de telemetria vai para o Pub/Sub.
\textbf{O resto da aplicação fala PostgreSQL direto} --- inclusive o canal
\emph{operacional} do motor.

\textbf{Alternativa descartada.} Tudo síncrono no PostgreSQL. Colocaria um
\cd{INSERT} sincronizado por interação de estudante dentro da mesma transação que
serve a tela.

\textbf{Consequências.} Uma peça de infraestrutura a mais --- e é a única exceção
declarada à regra \enquote{nada de fila sem gargalo medido}, junto com
\adrr{ADR-07}. Em troca: o pico é amortecido por um \emph{buffer} \textbf{durável}
(ver \adrr{ADR-14}).

\textbf{Gatilho.} Nenhum. É decisão tomada de antemão e mantida.

% ---------------------------------------------------------------------------
\section{ADR-07 --- IA atrás de um \emph{gateway} na própria API}

\textbf{Decisão.} Toda chamada de IA passa por \arq{apps/api/src/ai}, com
interface neutra de \emph{vendor}, \emph{prompt} de sistema no servidor, quota
diária persistida por perfil e \emph{rate limit}. Pré-gerar em lote sempre que a
pedagogia permitir.

\textbf{Alternativas descartadas.}
\begin{itemize}
  \item \emph{Chamada direta do \emph{bundle}}: exporia chave e \emph{prompt} de
    sistema, e tornaria a quota incontrolável.
  \item \emph{Serviço de IA separado}: um Cloud Run a mais para operar, sem
    ganho --- a API já tem CPU sempre alocada.
\end{itemize}

\textbf{Consequências.} Nenhum módulo fora de \arq{src/ai} conhece o
\emph{vendor}; trocar de provedor é criar um arquivo. \emph{Ruim}: a API carrega o
custo de latência das gerações longas, o que obriga o padrão 202 +
\emph{polling}.

\textbf{Gatilho.} Volume de IA que justifique isolar o consumo, ou requisito de
GPU própria.

% ---------------------------------------------------------------------------
\section{ADR-08 --- RLS forçada e uma porta única de banco}

\textbf{Decisão.} \cd{ENABLE} \emph{e} \cd{FORCE ROW LEVEL SECURITY} em toda
tabela de dado de estudante; acesso ao banco \textbf{exclusivamente} por
\cd{PrismaService.comTenancy}, que abre transação e executa
\cd{set\_config('app.escola\_id', ..., true)} antes do trabalho.

\textbf{Alternativas descartadas.}
\begin{itemize}
  \item \emph{Apenas filtro na aplicação}: uma \emph{query} nova que esqueça o
    filtro devolve dado de outra escola, e ninguém percebe.
  \item \emph{Apenas RLS}: um erro de migração desliga a proteção inteira.
  \item \emph{Duas \emph{roles}} (dona para migração, sem \emph{ownership} para a
    aplicação): desnecessário com \cd{FORCE}, que sujeita o próprio dono às
    políticas.
\end{itemize}

\textbf{Consequências.} Duas trancas independentes. O terceiro argumento
\cd{true} do \cd{set\_config} é \textbf{crítico}: local à transação, então a
conexão volta limpa ao \emph{pool}. \cd{comTenancy} sem escola \textbf{lança} ---
não consulta sem filtro. \emph{Ruim}: toda leitura, inclusive trivial, abre
transação.

\textbf{Gatilho.} Nenhum. É a decisão que sustenta o produto.

% ---------------------------------------------------------------------------
\section{ADR-09 --- Domínio por experiência na API; gabarito \emph{server-side}}

\textbf{Contexto.} A exp-02 (2026-08-24) foi a primeira experiência com gabarito.

\textbf{Decisão.} \textbf{Não há módulo genérico de avaliação.} Cada experiência
com lógica de domínio ganha um módulo próprio na API (\arq{src/<slug>/}), com
\arq{catalogo.ts} + \arq{regras.ts} + \emph{service} + \emph{controllers} por
papel. O gabarito nunca vai ao \emph{bundle}; a API corrige.

\textbf{Alternativa descartada.} Um módulo genérico de avaliação com
\enquote{itens} e \enquote{respostas corretas}. Uma API genérica de correção
convida a mandar a resposta certa junto --- e cada experiência tem uma mecânica
que não cabe num formato comum.

\textbf{Consequências.} Decisões concretas herdadas desta ADR: ranking com
colegas \textbf{fictícios} (sem pontuação de menor real exposta a outro menor);
URLs das fontes como \textbf{texto} (o \emph{sandbox} bloqueia \emph{popups}, e o
estudante não sai da plataforma); salvos \textbf{congelados} depois da conclusão
(nota e estado não podem divergir); estudante roteado pela \textbf{série}.
\emph{Ruim}: um módulo novo por experiência --- 48 módulos no limite.

\textbf{Gatilho.} Se três experiências consecutivas tiverem mecânica idêntica,
extrair o comum. Nunca antes.

% ---------------------------------------------------------------------------
\section{ADR-10 --- Bem-estar: instrumento fixo, IA só no entorno, sem gestor}

\textbf{Contexto.} 2026-08-25. Um \emph{check-in} de bem-estar \emph{do
professor}, com o DASS-21.

\textbf{Decisão.} O instrumento entra \textbf{\emph{verbatim}} (versão brasileira
de Vignola \& Tucci~\cite{vignola2014dass21}) em microdoses de três itens por dia.
A IA gera abertura, pergunta aberta e acolhimento, atrás de
\cd{regras.validarTextoDaIa}, com \emph{fallback} fixo, e \textbf{nunca vê}
pontuação, nome ou escola. Correção \emph{server-side} ao sétimo \emph{check-in}.
Item sentinela (17 ou 21 $\geq 2$) ou termo de risco \implica{} mensagem fixa de
apoio, sem IA.

\textbf{Por ser dado de saúde (LGPD art.~11): não há visão de gestor nem agregado
por escola.} RLS + posse por \cd{professor\_id}, e o próprio professor apaga tudo.

\textbf{Alternativas descartadas.}
\begin{itemize}
  \item \emph{IA reescrevendo os itens} para \enquote{soar mais natural}:
    quebra a validação psicométrica e invalida as faixas de corte.
  \item \emph{Questionário completo de uma vez}: engajamento vem de custar menos
    de um minuto, não de ser curto uma vez.
  \item \emph{Painel agregado para a rede}: sai de um \cd{GROUP BY} e é
    exatamente o que não se deve fazer com dado de saúde de uma equipe pequena,
    onde \emph{n} baixo reidentifica.
\end{itemize}

\textbf{Consequências.} O \emph{check-in} nunca deixa de abrir. \emph{Ruim}: a
rede não tem nenhuma visão de bem-estar dos professores --- o que é uma perda
real de informação de gestão, aceita.

\textbf{Gatilho.} Uma visão agregada e anonimizada é decisão futura, com DPA e
\emph{k}-anonimato explícito.

% ---------------------------------------------------------------------------
\section{ADR-11 --- Telemetria sem PII e vocabulário fechado}

\textbf{Decisão.} A telemetria carrega \cd{estudante\_id} \textbf{opaco}, e a
identidade é derivada do \emph{token} no servidor --- o que o cliente enviou nesses
campos é descartado. Em \cd{exp\_70\_sessoes.payload}, \cd{eventos[].dados} aceita
apenas número, booleano ou \emph{string} de um \textbf{vocabulário fechado}.

\textbf{Alternativa descartada.} Campo \cd{texto} livre no \emph{payload} (que
existia). Era o único canal de texto arbitrário --- e o caminho por onde o nome de
uma criança chegaria a um JSONB de longa duração.

\textbf{Consequências.} Trava \textbf{estrutural}, não lista de proibidos: não há
como enviar texto arbitrário, então não há como revisar mal. \emph{Ruim}: um
evento novo que precise de rótulo textual exige acrescentar o rótulo ao
vocabulário --- ou seja, exige uma decisão explícita.

\textbf{Gatilho.} Nenhum.

% ---------------------------------------------------------------------------
\section{ADR-12 --- \cd{grau} como invariante de API}

\textbf{Contexto.} O mapeamento comportamento \seta{} valor de Schwartz é
\textbf{hipótese de \emph{design}, não instrumento validado}: Schwartz se valida
por questionário (PVQ para adultos, PBVS-C para
crianças~\cite{schwartz2012overview,doring2010pbvsc}), não por telemetria de jogo.

\textbf{Decisão.} A API devolve uma \textbf{união discriminada} em que só a
variante \cd{scored} tem \cd{numero: number}; as outras têm \cd{null}
\textbf{por tipo}. E a distinção entre \cd{zero} (\enquote{teve a chance e não
foi}, entra na média) e \cd{none} (\enquote{a oportunidade não chegou}, sai da
conta) é preservada em todo o caminho.

\textbf{Alternativa descartada.} Manter a decisão na camada de interface, com um
comentário avisando. Foi assim no aplicativo de origem, e qualquer tela nova podia
imprimir \enquote{Poder 100} a partir de um sinal só. \textbf{Prosa não é
controle; um tipo é.}

\textbf{Consequências.} O hub fica estruturalmente incapaz de desenhar o número
errado --- ele nunca o recebe, e o compilador nem deixa procurar. Decisões
herdadas: o relatório abre pelo \textbf{nível} do instrumento antes do desempenho;
a listagem de turma \textbf{não} mostra valor psicométrico (listagem é ordenável,
e uma coluna \enquote{Poder} viraria ranking de crianças);
\cd{cartoesServidos} $\neq$ \enquote{apareceram}.

\textbf{Gatilho.} Se um instrumento validado por questionário for adotado, a
união ganha uma variante --- não perde a distinção.

% ---------------------------------------------------------------------------
\section{ADR-13 --- Retenção sem prazo declarado}

\textbf{Contexto.} 2026-08-26. Havia \cd{expires\_at} de 90 dias no \emph{schema}
da exp-70 e \textbf{nenhuma rotina que o cumprisse}.

\textbf{Decisão.} A coluna \textbf{saiu}
(\arq{20260826170000\_vila\_sem\_prazo}). Não há prazo de retenção declarado para
\cd{exp\_70\_sessoes} nem para \cd{telemetria\_eventos}.

\textbf{Alternativa descartada.} Manter a coluna. \textbf{Coluna de prazo que
ninguém lê é pior que prazo nenhum}, porque quem lê o \emph{schema} acredita nela
--- é a mesma armadilha de um campo \cd{consentimento\_em} que existia no
aplicativo de origem: a regra saiu do código, a coluna ficou, e todo mundo
continuou acreditando.

\textbf{Consequências.} Ganha-se o \textbf{reprocessamento}, que é requisito real
do instrumento: \cd{payload} é a única fonte para recomputar o perfil quando
\cd{VERSAO\_MAPEAMENTO} subir, e 90 dias matavam exatamente a comparação entre
duas gerações do mapeamento. \emph{Ruim, e é explícito}: dado comportamental de
menor fica armazenado indefinidamente, e isso é uma exposição de conformidade em
aberto.

\textbf{Gatilho.} Exigência contratual de prazo. Quando vier, ele volta
\textbf{junto com a rotina que o cumpre} --- e, com a tabela particionada
(\adrr{ADR-14}), essa rotina é \cd{DROP TABLE particoes.<mês>}. A pergunta é da
plataforma inteira, não de um módulo.

% ---------------------------------------------------------------------------
\section{ADR-14 --- Telemetria no PostgreSQL particionado}

\textbf{Contexto.} 2026-08-27. Avaliou-se mandar o canal analítico para o BigQuery
por \emph{subscription} nativa do Pub/Sub. Mediu-se também o custo real do
caminho de gravação.

\textbf{Decisão.} O canal analítico grava no \textbf{mesmo PostgreSQL}, com
\cd{INSERT} multi-linha e \cd{telemetria\_eventos} particionada por mês.
\textbf{Sem BigQuery. Sem Redis.}

\textbf{Alternativas descartadas.}
\begin{itemize}
  \item \emph{BigQuery}: os relatórios do professor \textbf{não leem}
    \cd{telemetria\_eventos} (leem \cd{exp\_01\_events}), então não resolvia
    problema nenhum hoje --- e criava uma segunda plataforma de dados para
    operar.
  \item \emph{Redis como \emph{buffer}}: quem amortece o pico já é o Pub/Sub, e
    ele é \textbf{durável}. Um \emph{buffer} em memória perde os eventos do
    processo que morre antes do \emph{flush}, ou seja, reintroduz a perda que o
    Pub/Sub eliminou. \textbf{O gargalo medido era o laço de \cd{INSERT}} --- um
    lote de 50 custava 50 \emph{round-trips}, o que no pico dava
    $\sim$800 \cd{INSERT}/s; em lote, $\sim$16/s.
\end{itemize}

\textbf{Consequências.} A chave primária é composta \cd{(event\_id, ts)}, porque o
PostgreSQL exige a chave de partição em todo índice único --- e a \emph{dedupe}
segue correta, porque a reentrega traz o mesmo \cd{ts}. As partições vivem no
\emph{schema} \cd{particoes}, fora do alcance do Prisma. Há partição
\cd{DEFAULT} como rede de segurança, porque evento recusado é telemetria perdida
em silêncio. \emph{Ruim}: as faixas vão até 2028-12 e \textbf{não há \emph{job}
agendado} --- criar partição é item de operação.

\textbf{Gatilho.} Volume que sature o PostgreSQL, ou um relatório que
efetivamente precise de análise sobre \cd{telemetria\_eventos} em escala.

% ---------------------------------------------------------------------------
\section{ADR-15 --- O hub assina os \emph{tokens} de experiência}

\textbf{Contexto.} O IdP central (WSO2) tem ciclo de vida próprio e não estava
pronto quando as experiências precisaram de \emph{token}.

\textbf{Decisão.} O \textbf{hub} assina os \emph{tokens} de experiência (RS256,
TTL 15 min, \emph{claims} \cd{\{sub, escola\_id, turma\_id?, perfil\}}, \cd{aud}
\cd{plataforma-api}) e serve o JWKS em \rota{/.well-known/jwks.json}. A API valida
por \env{OIDC\_ISSUER} e \env{OIDC\_JWKS\_URI}.

\textbf{Alternativa descartada.} Esperar o WSO2. Faria toda a plataforma depender
de um servidor Java em \emph{bootstrap} para que qualquer experiência pudesse ser
testada.

\textbf{Consequências.} Trocar para o WSO2 é trocar duas variáveis --- o
\cd{OidcGuard} não muda uma linha. \textbf{Efeito colateral útil:} o CI usa a
mesma chave para assinar o \emph{token} de cinco minutos do \emph{upsert} de
manifesto, o que é o motivo de as experiências serem publicadas \emph{depois} do
hub. \emph{Ruim}: o hub acumula responsabilidade de emissor, e
\env{TOKEN\_SIGNING\_KEY\_PEM} passa a ser obrigatória em produção.

\textbf{Gatilho.} WSO2 pronto e configurado por contrato. A ADR é
\textbf{transitória por desenho}.

% ---------------------------------------------------------------------------
\section{ADR-16 --- ADC como credencial; modelo \emph{versus} região}

\textbf{Decisão, parte 1 --- credencial.} O \emph{provider} de IA obtém o
\emph{token} com \cd{google-auth-library}, que resolve
\env{GOOGLE\_APPLICATION\_CREDENTIALS} \seta{} arquivo de
\cd{gcloud auth application-default login} \seta{} \emph{metadata server}.
\textbf{Nunca chave de API} --- não existe variável para ela no projeto.

\textbf{Contexto.} Antes havia \emph{só} o \emph{metadata server}, e por isso
\emph{nenhuma} chamada de IA funcionava fora do Cloud Run. Em desenvolvimento a
API ficava presa em \cd{mock} e a exp-90 devolvia um plano \textbf{falso} --- o
que é pior que falhar, porque parecia funcionar.

\textbf{Decisão, parte 2 --- modelo e região.} Modelo \cd{gemini-3.6-flash},
localização \cd{global}, para texto \emph{e} para plano de aula.

\textbf{Medição.} O \cd{3.6-flash} dá \textbf{404} em \cd{southamerica-east1},
\cd{us-east1} e \cd{us-central1}. Só o endpoint \cd{global} o serve.

\textbf{Consequência aceita, e ela é séria.} \textbf{O tráfego de IA sai de São
Paulo, inclusive o do estudante}, e isso precisa constar do contrato e do DPA da rede.

\textbf{Gatilho.} Exigência de residência de dados. Voltar à região é trocar
\env{VERTEX\_LOCATION} \textbf{e} o modelo, juntos --- a região publica a família
2.5, não a 3.x.

\textbf{Decisões relacionadas.} \env{VERTEX\_THINKING\_BUDGET} = \cd{0}: nos
modelos 2.5 e 3.x os \emph{tokens} de pensamento saem do \cd{maxOutputTokens}, e
com pensamento livre o plano anual voltava \textbf{vazio} com
\cd{finishReason: MAX\_TOKENS}. Teto de saída do plano em \textbf{64k}, não 32k.
Modelo de imagem \cd{gemini-3.1-flash-image}, com proporção \cd{4:5}
\textbf{pedida} ao modelo --- o \emph{default} anterior não existe mais, e sem
pedir a proporção o \cd{cover} descarta metade da largura.

% ---------------------------------------------------------------------------
\section{ADR-17 --- \emph{Claim} atômico em vez de fila externa}

\textbf{Contexto.} Duas rotinas longas e assíncronas: geração de \emph{feed}
(exp-01) e geração de plano de aula (exp-90). Duas ou três instâncias do Cloud Run
rodando ao mesmo tempo.

\textbf{Decisão.} \emph{Claim} atômico na própria linha
(\cd{UPDATE ... WHERE status = 'pendente'}, com contagem de linhas afetadas) e
\emph{worker} \emph{in-process}. Para o motor de \emph{feed}, \emph{lock}
consultivo por \cd{run} (\cd{pg\_advisory\_xact\_lock}).

\textbf{Alternativas descartadas.}
\begin{itemize}
  \item \emph{Fila externa} (Pub/Sub, Cloud Tasks): peça a mais para operar, para
    um problema que um \cd{UPDATE} condicionado resolve.
  \item \emph{Sem \emph{lock} no motor}: dois pedidos simultâneos de bloco (duas
    abas, \emph{reload} no meio do \emph{fetch}, \emph{retry}) se sobrescreviam e
    um bloco inteiro desaparecia do estado --- posts entregues voltavam a ser
    servidos e a série por bloco do dossiê ficava corrompida.
\end{itemize}

\textbf{Consequências.} A API precisa de \textbf{CPU sempre alocada} e
\cd{min-instances} $\geq$ 1, porque há trabalho fora do ciclo de requisição.
Retomada de órfã é disparada pelo \emph{polling} do painel, não por
\emph{scheduler}. \emph{Ruim}: se todas as instâncias morrerem e ninguém abrir o
painel, uma geração órfã fica parada até que alguém abra.

\textbf{Gatilho.} Primeira reclamação real de geração parada.

% ---------------------------------------------------------------------------
\section{ADR-18 --- Exceções registradas}

\textbf{Contexto.} Duas regras da plataforma têm exceção, e exceção não
registrada vira precedente.

\textbf{Exceção 1 --- \cd{any} e ESLint.}
\arq{apps/experiencias/exp-70-vila-do-lume/public/jogo/} --- cerca de dez mil
linhas de JavaScript clássico com globais compartilhadas e \arq{three.min.js}
vendorizado --- fica como está, e está nos \cd{ignores} do ESLint.
\textbf{Motivo:} é um instrumento com distâncias e tempos calibrados, e o
compilador não tem como saber que uma refatoração mudou a \emph{medida}. O
\emph{shim} (\arq{src/}, TS \emph{strict}) é a única porta: \cd{window.VILA}. A
rede que substitui o \cd{tsc} ali é o \emph{smoke} visual.

\textbf{Exceção 2 --- \emph{design system}.} A exp-70 não usa o
\cd{@unooria/design-system}. \textbf{Motivo:} é jogo em \emph{canvas} 3D, sem DOM
de interface.

\textbf{Consequências.} As duas exceções são \textbf{nominais} --- valem para uma
pasta e uma experiência, não para uma categoria. Qualquer outra experiência que
peça a mesma exceção precisa de uma ADR nova.

% ---------------------------------------------------------------------------
\section{ADR-19 --- \emph{Design system} externo como fonte única de interface}

\textbf{Decisão.} Todo elemento de interface vem do \cd{@unooria/design-system},
consumido pela ponte \arq{packages/ui}. Componente que falta \textbf{nasce no
repositório do \emph{design system}} e chega por versão nova --- não como arquivo
local.

\textbf{Alternativa descartada.} Componentes locais em \arq{packages/ui}. Com 48
experiências, um componente local viraria 48 divergências --- e a fidelidade ao
clone é o que a experiência \emph{mede}.

\textbf{Consequências.} Duas famílias no mesmo pacote que \textbf{nunca se
cruzam}; a linha entre elas é a \textbf{revelação}, não o arquivo. Compensação
local só como exceção marcada no CSS, com pedido correspondente ao repositório do
\emph{design system}. \emph{Ruim}: uma lacuna do pacote bloqueia a experiência até
a próxima publicação --- e \cd{\^{}0.1.x} não sobe \emph{minor}, então atualizar é
ato deliberado.

\textbf{Gatilho.} Nenhum. É decisão de produto e de marca, não de engenharia.

% ---------------------------------------------------------------------------
\section{ADR-20 --- Um banco, um \emph{schema}; prefixo \cd{exp\_NN\_}}

\textbf{Decisão.} Um banco, um \emph{schema} lógico. Tabela de experiência leva
\cd{exp\_NN\_}; tabela de plataforma fica sem prefixo; tabela que é contrato entre
duas experiências carrega \textbf{os dois números}
(\cd{exp\_01\_exp\_70\_perfis\_importados}).

\textbf{Alternativas descartadas.}
\begin{itemize}
  \item \emph{Banco por experiência}: elas \textbf{não} são isoladas --- o perfil
    da Vila alimenta a largada do Instanooria, e o PostgreSQL não faz \cd{JOIN}
    entre bancos. E um \emph{pool} do Prisma por experiência estouraria
    \cd{max\_connections}.
  \item \emph{\emph{Schema} por experiência}: a única coisa que daria além do
    prefixo é fronteira de \cd{GRANT}, e existe \textbf{uma} \emph{role} de
    aplicação.
\end{itemize}

\textbf{Consequências.} Migração de referência
\arq{20260827090000\_tabelas\_por\_experiencia}: \cd{RENAME}, nunca
\cd{DROP}/\cd{CREATE}, renomeando \emph{constraints}, índices e políticas junto.
Duas armadilhas registradas lá: \cd{FOR ... IN SELECT} sobre o catálogo
\textbf{relê o que o próprio laço alterou}, e \cd{RENAME CONSTRAINT} renomeia
\textbf{também} o índice que a sustenta --- as duas produzem
\cd{exp\_01\_exp\_01\_...\_pkey}.

\textbf{Gatilho.} Necessidade de fronteira de \cd{GRANT} por experiência --- o
que só faria sentido com mais de uma \emph{role} de aplicação.

% ---------------------------------------------------------------------------
\section{ADR-21 --- \emph{Push subscription} para a própria API}

\textbf{Decisão.} O canal analítico é consumido por \emph{push subscription}
direto num \emph{endpoint} da própria API (\rota{POST /telemetry/ingest}), e não
por um \cd{\{contrato\}-telemetry-worker} separado.

\textbf{Alternativa descartada.} \emph{Worker} dedicado, como previa o desenho
original de infraestrutura. Seria um Cloud Run a mais com \emph{deploy}, log,
escala e \emph{alerting} próprios.

\textbf{Consequências.} A rota grava dado de estudante e \textbf{não} passa pelo
\cd{OidcGuard}, então precisa de tranca própria: \cd{PubSubPushGuard} verifica
emissor Google, \cd{aud} e \cd{email} da conta de serviço de \emph{push}, e
\textbf{falha fechada} sem configuração. O Terraform casa os dois lados. Gravação
idempotente por \cd{ON CONFLICT}, porque a entrega é \emph{at least once}.

\textbf{Gatilho.} Volume que justifique isolar o consumo do tráfego de estudante.

% ---------------------------------------------------------------------------

\section{ADR-22 --- exp-02 1.1.0: \emph{briefing} antes do disfarce}

\textbf{Contexto.} Até a 1.0.0 o estudante da exp-02 caía direto no \emph{feed} de
\emph{Shorts} sem nenhum enquadramento: a primeira vez que a palavra
\enquote{missão} aparecia para ele era no diálogo que abria sozinho depois do
terceiro salvo --- isto é, \emph{depois} de a medida já ter sido tomada. Não havia
contador visível de favoritos, não havia bandeja, e o teto de três se manifestava
como um \emph{toast} de 5,2\,s. O Figma de redesenho (arquivo
\cd{oX1NGR8sPyduAS52hDaJFX}, seis telas de 1440$\times$900, público de nove anos)
acrescentou o enquadramento antes da ação e o retorno durante ela.

\textbf{Decisão.} Três mudanças, tomadas juntas:

\begin{enumerate}
  \item \textbf{Tela de entrada na voz da escola.} O \emph{briefing} usa a seção
    da \textbf{plataforma} (\cd{@edu/ui/ds}) --- a primeira vez que ela aparece
    \emph{antes} da revelação. O critério da linha da revelação
    (Capítulo~\ref{cap:ds}) continua sendo \enquote{quem está falando?}: no
    \emph{briefing} quem fala é a escola, e não há disfarce a proteger, porque a
    criança ainda não viu um \emph{Short}.
  \item \textbf{Conclusão pelo botão.} O terceiro salvo apenas habilita
    \enquote{Concluir missão} na bandeja; o diálogo abre no clique.
  \item \textbf{Ganho de inscritos visível.} Um \emph{badge} \cd{+N} ao lado do
    chip do canal, por 1,4\,s, quando o número sobe.
\end{enumerate}

\textbf{Alternativas descartadas.} (a) Reproduzir o \emph{wireframe} ao pé da
letra, com o azul e o tema claro fixo dele --- forçaria a paleta da plataforma
para dentro da ficção e quebraria o seletor claro/escuro do estudante; o Figma é
\emph{wireframe}, e vale como especificação de estrutura, texto e estados, não de
cor. (b) Reduzir o catálogo a oito vídeos, como o \emph{wireframe} ilustra ---
mexeria no instrumento (o balanço 10~\emph{fake} / 11 factual) e na
comparabilidade com dados já coletados. (c) Manter a abertura automática do
diálogo no terceiro salvo, com o botão como caminho alternativo --- duas portas
para a mesma coisa, e a primeira empurra a criança para a virada no instante do
toque. (d) Animar só o total, sem o \emph{badge} de delta.

\textbf{Consequências ruins, aceitas.} O \emph{badge} torna legível
\textbf{durante} a missão que um vídeo vale $+$1.000 inscritos e outro $+$50, ou
seja, expõe o padrão que o \emph{ranking} existe para revelar. Foi aceito porque
o total já se movia --- o delta sempre foi derivável de dois quadros seguidos ---
e porque a nota da etapa 3 do Figma decide isso explicitamente: o ganho varia, e
o motivo \textbf{não} é explicado ali. A tela de entrada também gasta um toque a
mais antes do \emph{feed}, e um estudante apressado a atravessa sem ler.

\textbf{O que não mudou.} As curtidas continuam fora do cálculo; o gabarito
(\cd{factual}) continua no servidor e fora do \emph{wire} --- há teste que
enumera os campos públicos do catálogo; o \emph{score} continua saindo de
\rota{POST /unootube/credibilidade}. Nenhuma rota nova, nenhuma migração.

\textbf{Efeito colateral corrigido no hub.} O \emph{feed} passou a montar depois
de um clique \emph{dentro} do \cd{iframe}, e isso revelou um defeito que já
existia: clicar dentro do \cd{iframe} tira o foco da janela do hub, o
\arq{ExpFrame.tsx} traduzia esse \cd{blur} em \cd{hub:pause}, e o \cd{focus}
correspondente nunca voltava enquanto o estudante continuasse ali --- o vídeo parava
no primeiro toque em \enquote{Salvar} e não voltava mais, em qualquer versão. A
correção é uma linha: quando \cd{document.activeElement} é o próprio
\cd{iframe}, o foco foi para a experiência e não há perda de visibilidade
nenhuma.

\textbf{Gatilho de reavaliação.} Distribuição de \cd{fontesConfiaveis} da turma
subindo de forma anômala --- o \emph{badge} é o primeiro suspeito. Ou o professor
relatando que a turma atravessa o \emph{briefing} sem ler, o que pediria que ele
fosse do professor e não da tela.

% ---------------------------------------------------------------------------
\section{ADR-23 --- \emph{Function calling} entra no \cd{AiProvider}}
\label{adr:ferramentas}

\begin{especificacao}{Contexto}
A exp-03 (Missão Tempestade) tem um agente conversacional em que \textbf{nenhum
número pode sair do modelo de linguagem}: as estimativas vêm de um motor
probabilístico determinístico, e o agente apenas narra o que a ferramenta
devolveu. O projeto de origem implementava isso com um laço de \emph{function
calling} do SDK do vendor --- até seis idas e voltas por turno, com três
tentativas que reaproveitam o motivo da recusa anterior. A interface
\cd{AiProvider} desta plataforma tinha \cd{generate}, \cd{generateJson} e
\cd{generateImage}: nenhuma delas conversa.
\end{especificacao}

\textbf{Decisão.} A interface ganha um quarto método, \cd{conversar}, com tipos
neutros de vendor (\cd{AiFerramenta} reusa o \cd{AiJsonSchema} que já existia). O
laço de ferramentas, a preservação da assinatura de raciocínio e o \emph{retry}
que aprende ficam dentro de \arq{vertex.provider.ts}; as regras de recusa ficam
no domínio, passadas como \emph{callback} \cd{validar}. \cd{tokensUsados} é a
soma de todas as voltas e de todas as tentativas, \textbf{inclusive as
rejeitadas} --- é o que foi pago.

\textbf{Alternativas descartadas.}
\begin{itemize}
  \item \emph{A experiência falar com o vendor}: violaria a regra de que só
    \arq{src/ai} conhece o provedor, e poria a chave e o protocolo do vendor no
    caminho de uma experiência.
  \item \emph{Módulo da experiência chamar o SDK do vendor diretamente}: mesmo
    problema um nível acima --- o dia de trocar de provedor passaria a exigir
    mexer em cada experiência que conversa.
  \item \emph{Resposta estruturada em vez de ferramentas} (\cd{generateJson}
    devolvendo uma lista de ações): funcionaria para um turno simples, mas o
    agente encadeia \emph{carregar fonte} \seta{} \emph{cruzar dados} \seta{}
    \emph{narrar} decidindo cada passo com o resultado do anterior. Uma lista de
    ações declarada de antemão não tem essa realimentação.
\end{itemize}

\textbf{Consequências.}
\begin{itemize}
  \item Qualquer experiência futura ganha agente com ferramentas de graça, e a
    troca de provedor continua sendo um arquivo novo em \arq{src/ai}.
  \item \emph{Ruim}: \cd{conversar} é obrigatório na interface, então as três
    implementações existentes tiveram de ganhá-lo --- inclusive os \emph{fakes}
    dos testes de outros módulos.
  \item \emph{Ruim}: o \cd{MockProvider} não conhece os \emph{slugs} do domínio,
    então chama a primeira ferramenta declarada com argumento vazio. Em
    desenvolvimento o mapa não enche pela conversa; enche pelo painel.
  \item Três defeitos do provedor ficaram cobertos por teste com \cd{fetch}
    dublado: a assinatura de raciocínio reenviada \emph{verbatim} (sem ela o
    Vertex recusa a continuação com 400 e a narrativa nunca chega), o descarte
    das partes de pensamento (senão o raciocínio do modelo aparece na tela do
    estudante) e o teto de voltas (modelo preso pedindo ferramenta vira recusa, não
    bolha vazia).
\end{itemize}

\textbf{Gatilho de reavaliação.} Um provedor cujo protocolo de ferramentas não
caiba na assinatura de \cd{conversar} --- por exemplo, ferramentas paralelas na
mesma volta. Ou latência de turno acima do que a sala tolera, que empurraria para
\cd{streamGenerateContent} e mudaria a forma do retorno.

% ---------------------------------------------------------------------------
\section{ADR-24 --- Teto de custo de IA por SESSÃO}
\label{adr:teto-sessao}

\begin{especificacao}{Contexto}
O \cd{QuotaService} limita tokens por usuário e por dia, com teto por perfil
(\env{AI\_DAILY\_TOKEN\_QUOTA\_PER\_ESTUDANTE} valia 20.000). Esse número foi
dimensionado para o tutor da exp-01: uma conversa curta. O agente da exp-03 tem
\emph{prompt} de sistema de cerca de 11.500 caracteres, reenviado inteiro a cada
volta de ferramenta e a cada tentativa; uma partida de sete turnos fica na ordem
de 100 a 200 mil tokens --- duas ordens de grandeza acima.
\end{especificacao}

\textbf{Decisão.} O controle que importa passa a ser \textbf{por sessão}, no
módulo da experiência: \cd{tokens\_gastos} e \cd{chamadas\_ia} na linha da
partida, com tetos em \arq{regras/defaults.ts}. A quota diária continua como
\textbf{fusível da plataforma} e sobe para acomodar duas partidas. A checagem
acontece antes de chamar o provedor, e o consumo é registrado no \cd{finally} ---
inclusive quando o turno falha, porque o token foi gasto de todo modo.

\textbf{Alternativas descartadas.}
\begin{itemize}
  \item \emph{Quota por experiência no \cd{QuotaService}}: um limite por
    \cd{expId} em variável de ambiente. Não mexeria no teto do tutor, mas
    complicaria um serviço que hoje é simples e por perfil --- e continuaria
    sendo por dia, quando o que tem começo e fim é a sessão.
  \item \emph{Encolher o instrumento para caber em 20.000}: reduzir o
    \emph{prompt}, as voltas de ferramenta e o histórico. Cabe na quota atual e
    mexe justamente no que foi calibrado em \emph{playtest}.
  \item \emph{Teto em centavos}, como no projeto de origem: preço depende de
    tabela do vendor, que muda sem aviso e não é conhecida pela aplicação. Token
    é o que o provedor devolve.
\end{itemize}

\begin{especificacao}{Medição --- 2026-08-27, \cd{gemini-3.6-flash} no endpoint \cd{global}}
Uma \textbf{partida completa} (8 turnos, 7 chamadas de IA, do briefing ao
desfecho) custou \textbf{\num{71397} tokens} em 40 segundos.

Por turno: 7,5k a 8k na maioria, com dois picos (20,4k e 16,4k) nos turnos de
mais voltas. O turno de conversa pura --- \enquote{não entendi, explica de
novo} --- custou \num{3846} e \textbf{não debitou pergunta}, que é a regra
funcionando.

E o que domina o custo \textbf{não é o tamanho do contexto: é o número de
voltas} de ferramenta, porque cada volta reenvia o histórico inteiro. Um turno de
duas voltas custou \num{11481}; o mesmo turno com quatro fontes já carregadas e
seis falas de histórico, mas uma volta, custou \num{8463} --- \emph{menos}, com
mais contexto.
\end{especificacao}

\textbf{Os números, com a aritmética.} \cd{TETO\_TOKENS\_SESSAO} $=$ \num{150000}:
o dobro de uma partida completa medida, o que cabe a missão inteira com folga e
ainda corta antes de o custo virar problema. \cd{TETO\_CHAMADAS\_SESSAO} $=$ 24:
ao custo medido o teto de \emph{tokens} morde primeiro, por volta do 19º turno ---
este é a rede contra o laço, não o controle de custo.

\begin{atencao}
A medição encontrou um defeito que a estimativa não teria encontrado: com
\env{AI\_DAILY\_TOKEN\_QUOTA\_PER\_ESTUDANTE} em \num{20000} --- o valor dimensionado
para o tutor da exp-01, e que estava \textbf{fixo no \emph{compose} de
produção} --- a missão recusava o \textbf{quarto turno} com \enquote{a sala de
operações fechou por hoje}. A quota subiu para \num{300000} (cerca de quatro
partidas, porque \enquote{Recomeçar} abre partida nova) em quatro lugares: os dois
\emph{composes}, o \emph{default} da validação de ambiente e o
\arq{.env.example}.
\end{atencao}

\textbf{Consequências.} O custo fica limitado onde ele acontece, e a mensagem de
estouro é uma frase de jogo (\enquote{a sala de operações fechou por hoje}) em vez
de vocabulário de infraestrutura --- o estudante não sabe o que é quota. \emph{Ruim}:
são dois tetos no mesmo caminho, e um estouro exige olhar os dois para saber qual
disparou.

\textbf{Gatilho de reavaliação.} A primeira medição real. E a terceira
experiência com agente conversacional, que tornaria o teto por sessão um padrão a
subir para a plataforma.

% ---------------------------------------------------------------------------
\section{ADR-25 --- O Relógio da Aula é tabela da experiência}
\label{adr:relogio-turma}

\begin{especificacao}{Contexto}
A Missão Tempestade tem uma restrição de tempo que é parte do instrumento: cinco
minutos de relógio de parede valem a hora do cenário. No projeto de origem esse
cronômetro armava sozinho quando o estudante entrava na investigação --- e numa sala
isso é uma corrida com largadas diferentes: quem chegou primeiro terminava antes
de o professor acabar de explicar. As colunas do relógio moravam em
\cd{turmas}.
\end{especificacao}

\textbf{Decisão.} O relógio é da \textbf{turma}, e vive em
\cd{exp\_03\_relogio\_turma} --- tabela \emph{da experiência}, com
\cd{escola\_id} e RLS como as outras. Quem dá a largada é o professor, no painel,
e ela vale para a sala inteira ao mesmo tempo. Enquanto \cd{iniciado\_em} é nulo o
estudante vê o cronômetro parado e o campo de fala desabilitado.

\textbf{Alternativa descartada.} \emph{Colunas em \cd{turmas}}, como no origem.
\cd{turmas} serve a plataforma inteira: uma coluna de cronômetro de uma missão
faria toda experiência futura herdar o vocabulário dela, e a próxima que
precisasse de tempo teria de escolher entre reusar colunas com nome errado ou
adicionar as suas. É o mesmo raciocínio do prefixo \cd{exp\_NN\_}
(\adrr{ADR-20}).

\textbf{Consequências.}
\begin{itemize}
  \item A sala começa junto, e a espera passa a ser a aula --- não uma tela
    travada. A ficha do professor abre com essa faixa, porque \enquote{não abre}
    é o suporte que ela evita.
  \item A extensão de tempo usa um \textbf{pivô}: sem ele, esticar de cinco para
    sete minutos jogaria a cidade do minuto 48 para o 34, e a chuva andaria para
    trás na cara da turma.
  \item A largada é \textbf{idempotente} --- é exatamente o clique que um
    professor nervoso dá duas vezes.
  \item Estudante sem turma vê \enquote{aguardando}, que é a verdade. O
    \emph{fallback} do origem armava um cronômetro individual, e com ele um estudante
    solto jogava sem largada.
  \item \emph{Ruim}: o professor passa a ter uma ação sem a qual a aula não
    acontece. É dependência operacional nova, e o suporte precisa saber dela.
\end{itemize}

\textbf{Gatilho de reavaliação.} Uma segunda experiência com tempo de sala ---
aí o relógio da turma vira serviço da plataforma, não tabela de uma missão.

% ---------------------------------------------------------------------------
\section{ADR-26 --- Sem prazo de retenção na exp-03}
\label{adr:retencao-exp03}

\begin{especificacao}{Contexto}
O projeto de origem tinha \cd{expires\_at} de 30 dias em \cd{messages} e
\cd{marcos}, com uma rota de \emph{cron} que o cumpria. Esta plataforma não tem
\emph{job} agendado, por decisão --- e já registrou duas vezes o custo de uma
coluna de prazo sem rotina: o \cd{expires\_at} da exp-70 e o
\cd{turmas.consentimento\_em} do aplicativo de origem, onde a regra saiu do
código, a coluna ficou, e quem leu o \emph{schema} acreditou nela.
\end{especificacao}

\textbf{Decisão.} Nenhuma das cinco tabelas da exp-03 tem coluna de prazo. No
lugar entra \textbf{um} caminho de exclusão que existe de fato: \enquote{zerar a
turma} (\rota{DELETE /unoogpt/professor/turmas/:turmaId}), acionado por quem tem
a posse pedagógica, apagando as partidas \emph{daquela turma} e devolvendo o
relógio ao início.

\textbf{Alternativas descartadas.}
\begin{itemize}
  \item \emph{Portar o \cd{expires\_at} de 30 dias}: prazo que nenhuma rotina
    cumpre é pior que prazo nenhum, porque cria a crença de que há retenção.
  \item \emph{Trazer o prazo E o \emph{job}}: seria o primeiro agendador da
    plataforma, com o custo operacional que a decisão de não ter agendador
    evitou. Se um contrato exigir prazo, ele volta \textbf{junto com a rotina},
    nunca só a coluna.
  \item \emph{\cd{DELETE /unoogpt/dados} no estilo do bem-estar}: lá o professor
    é dono das linhas; aqui o dado é do estudante, e um botão \enquote{apagar tudo} no
    professor seria posse que ele não tem.
\end{itemize}

\textbf{Consequências.} O relógio é apagado junto, e o diálogo diz por quê:
rodar de novo com a aula \enquote{já começada} faria a segunda rodada nascer no
meio da tempestade. \emph{Ruim}: sem prazo, a transcrição de um menor de idade
fica no banco até alguém apagar --- e a pergunta de retenção continua aberta para
\cd{telemetria\_eventos} e para a plataforma inteira, não só para este módulo.

\textbf{Gatilho de reavaliação.} Exigência contratual de prazo, ou a primeira
rotina agendada da plataforma --- a partir dela o custo marginal de cumprir um
prazo cai muito.

% ---------------------------------------------------------------------------
\section{ADR-27 --- Transcrição é dado do instrumento}
\label{adr:transcricao}

\begin{especificacao}{Contexto}
A missão inteira acontece numa conversa em texto livre, escrita por menor de
idade. O professor precisa ler essa conversa: é nela que aparece \emph{como} a
criança perguntou, e o relatório sem ela é uma lista de contagens. Ao mesmo tempo,
o \cd{AGENTS.md} registra a trava oposta na exp-70 --- o campo de texto livre foi
removido do \emph{payload} porque era o caminho por onde o nome de uma criança
chegaria a um JSONB de 90 dias.
\end{especificacao}

\textbf{Decisão.} A transcrição vive em coluna dedicada
(\cd{exp\_03\_mensagens.texto}), com RLS por escola, e é servida
\textbf{apenas ao professor}: \cd{UnoogptTurmasController} tem
\cd{@Roles("professor", "admin-contrato")} --- e \cd{admin-contrato} é a
exceção cross-escola de sempre, não o gestor ---, enquanto o agregado do gestor
está em
\cd{UnoogptAgregadoController}, que \textbf{não tem nenhuma rota que aceite
\cd{estudanteId}}. Ela nunca vai a log nem à telemetria.

\textbf{Alternativa descartada.} Um controller só, com \cd{@Roles} incluindo
gestor e uma exceção dentro do método que devolve a criança. É a forma que vaza:
basta um método novo esquecer a exceção. Com dois controllers, nenhum descuido
futuro leva o gestor à conversa de uma criança, porque não existe rota para isso
no controller dele.

\textbf{Consequências.}
\begin{itemize}
  \item O relatório individual abre pelo \textbf{nível} --- graus, com
    \cd{zero} (\enquote{teve a chance e não foi}) distinto de \cd{none} (\enquote{a
    oportunidade não chegou}) --- antes do que a criança fez. É a regra da
    exp-70, e existe para o adulto não concluir sobre a criança o que o
    instrumento não mediu.
  \item A \textbf{listagem} da turma não tem coluna de resgatados: listagem é
    ordenável, e ordenar crianças por acerto viraria \emph{ranking}. O
    \emph{smoke} visual falha se uma aparecer.
  \item A telemetria analítica carrega contagem, nunca o texto.
  \item \emph{Ruim}: o dado mais sensível da experiência não tem prazo
    (\adrr{ADR-26}), e a única forma de apagá-lo é uma ação do professor.
\end{itemize}

\textbf{Gatilho de reavaliação.} Pedido de visão de rede sobre conduta ---
aí o caminho é agregado sem identidade no controller do gestor, nunca abrir a
linha da criança.

% ---------------------------------------------------------------------------
\section{ADR-28 --- Identidade por pessoa; OneRoster na fronteira}
\label{adr:identidade-enturmamento}

\begin{especificacao}{Contexto}
As experiências foram portadas de projetos onde \textbf{cada atividade tinha o
próprio cadastro de usuário}. No monorepo a unificação técnica já existia --- toda
experiência recebe \cd{estudante.id} pelo \cd{hub:init} e grava \cd{estudante\_id} com
chave estrangeira para \cd{users} ---, mas três resíduos impediam que
\enquote{acompanhar por estudante} fosse verdade:

\begin{enumerate}
  \item A conta era de \textbf{assento}, não de pessoa. Está escrito no
    \emph{schema} desde o porte: \cd{exp\_01\_runs.nome\_no\_app} existe
    \enquote{porque a conta é do APARELHO: outra turma reusa o mesmo login}. O nome
    real da criança vivia dentro de UMA experiência, e nenhuma das outras o
    enxergava.
  \item Não existia \textbf{enturmamento}. O vínculo era a coluna
    \cd{users.turma\_id}: um estudante cabia em uma turma só, o PROFESSOR não tinha
    turma nenhuma --- e por isso todo professor enxergava todas as turmas da
    escola --- e não havia ano letivo. Em fevereiro o 8A vira 9A e todo o
    acumulado do ano anterior passaria a ser lido pela turma nova, em silêncio.
  \item A entrada de pessoas fazia \emph{upsert} por \cd{usuario} e exigia
    \textbf{senha em texto claro no arquivo}. Corrigir um login criava um segundo
    estudante, com o histórico preso no primeiro.
\end{enumerate}
\end{especificacao}

\textbf{Decisão.} O modelo interno passa a ser \textbf{isomórfico ao OneRoster
1.2} (1EdTech): entram \cd{periodos\_letivos} (\cd{academicSessions}) e
\cd{matriculas} (\cd{enrollments}, a tripla usuário $\times$ turma $\times$ escola
com papel e datas), mais as colunas de identidade em \cd{users}, \cd{turmas} e
\cd{escolas} (\cd{matricula\_externa}, \cd{fonte\_id}, \cd{codigo\_inep}).

O padrão vive \textbf{só na fronteira} \arq{apps/api/src/rostering/} --- a mesma
disciplina de \arq{src/ai}, onde apenas o \emph{provider} conhece o fornecedor.
Nada do domínio conhece \cd{sourcedId}, o contrato do \cd{hub:init} não muda, e os
nomes seguem em pt-BR.

\textbf{Alternativas descartadas.}
\begin{itemize}
  \item \emph{Modelar do zero, sem padrão}: chegaríamos a quase a mesma tabela.
    O padrão não é o motivo --- é a conferência de que o desenho está certo, o
    vocabulário de graça e um argumento de edital.
  \item \emph{LTI 1.3 / NRPS}: responde \enquote{quem chegou agora}, em tempo real,
    e pressupõe estar embutido num LMS. Não provisiona quem ainda não entrou, que
    é exatamente o caso da matrícula da rede antes do ano letivo.
  \item \emph{Ed-Fi}: padrão de \emph{data warehouse} de rede; escopo muito maior
    que o problema, para uma equipe de uma pessoa.
  \item \emph{Esperar o padrão nacional}: o MEC publicou a Portaria 269, de
    24/03/2026, criando a Infraestrutura Nacional de Dados da Educação
    (EducaDados) e o Padrão Nacional de Interoperabilidade, com implantação
    gradual e especificação técnica ainda não publicada. Adotar o \emph{modelo}
    agora e manter o \emph{formato} atrás de adaptadores é o que serve aos dois
    futuros.
  \item \emph{Adotar também os \emph{bindings} do padrão}: o CSV do OneRoster são
    sete arquivos com identificadores globais; nenhum secretário de escola os
    preenche. A tela oferece uma \textbf{planilha simples}, de uma linha por
    pessoa, e o mesmo núcleo aplica o lote.
  \item \emph{\cd{demographics} do OneRoster}: raça e etnia no vocabulário do
    censo dos Estados Unidos, além de nascimento e sexo. Dado sensível
    (LGPD art.~11) que a plataforma não usa para nada --- não entra nem como
    coluna.
\end{itemize}

\textbf{Consequências.}
\begin{itemize}
  \item \cd{users.turma\_id} \textbf{sobrevive}, com novo significado: cache da
    matrícula principal. O \emph{login} e o \emph{claim} \cd{turma\_id} do
    \emph{token} (contrato congelado) o leem antes de haver escola no contexto,
    onde \cd{matriculas} --- que tem RLS --- não é alcançável.
  \item \enquote{Coordenação} continua sendo \cd{gestor}: \cd{PAPEIS} é contrato
    congelado entre o hub e a API, e um quinto papel atravessaria
    \emph{middleware}, \emph{guards} e RLS. A planilha aceita sinônimos
    (\cd{coordenacao}, \cd{diretor}, \cd{principal}) e normaliza.
  \item \emph{Ruim}: quatro tabelas e um módulo a mais, e uma migração nas duas
    tabelas mais centrais três meses antes de produção.
  \item \emph{Armadilha registrada}: consulta que leia \cd{matriculas} fora de
    \cd{comTenancy} volta \textbf{vazia}, sem erro. O sintoma é lista de pessoas
    \enquote{sem turma}.
\end{itemize}

\textbf{Nota de implementação --- o binding CSV, arquivo por arquivo.} O padrão
empacota os sete arquivos num zip com \cd{manifest.csv}. Aqui cada um é enviado
em uma chamada (\rota{POST /rostering/oneroster/:arquivo}), na ordem de
dependência: empacotamento é embalagem, não semântica, e aplicar um por vez dá o
mesmo resultado sem dependência de descompactação e sem lote pela metade guardado
em lugar nenhum. Fora de ordem também não corrompe: a linha que aponta para um
\cd{sourcedId} inexistente volta no relatório \textbf{dizendo qual arquivo
falta}, em vez de criar meia entidade. Quatro decisões dentro dessa leitura:

\begin{itemize}
  \item \cd{orgs.csv} só \textbf{aprende} o código INEP da escola desta
    importação. Escola é \emph{tenant}: criar um a partir de arquivo enviado por
    \emph{upload} seria abrir um inquilino com \emph{stack}, RLS e dados de
    menores por trás de um clique.
  \item Na 1.2 o papel saiu de \cd{users.csv} e foi para \cd{roles.csv}. A
    coluna \cd{role} é aceita quando vem (provedores seguem emitindo a da 1.1);
    sem ela, a pessoa nasce como \textbf{estudante} --- o menor privilégio --- e o
    relatório diz quantas linhas dependem do \cd{roles.csv}. Falhar fechado: um
    professor que entra como estudante perde acesso e reclama; um estudante que entrasse
    como professor leria a turma inteira.
  \item A coluna \cd{password}, que o padrão permite, é \textbf{ignorada}
    (\adrr{ADR-29}) --- e a contagem de linhas que a traziam vai ao relatório,
    para não sumir em silêncio.
  \item Em \cd{bulk}, ausência de linha \textbf{não} desativa ninguém: é
    contada e listada. Lote completo que chegou pela metade não pode apagar uma
    rede por omissão.
\end{itemize}

\textbf{Gatilho de reavaliação.} Publicação da especificação técnica do
EducaDados --- em especial se ela adotar o CPF como identificador do estudante,
que entra em coluna de reconciliação de acesso restrito e \textbf{nunca} em
telemetria, JSONB ou \emph{claim} de \emph{token}.

% ---------------------------------------------------------------------------
\section{ADR-29 --- Senha provisória com troca obrigatória}
\label{adr:senha-provisoria}

\begin{especificacao}{Contexto}
Antes da integração com o WSO2, alguém precisa criar as contas e entregar uma
credencial a 10 mil pessoas que nunca acessaram a plataforma. O formato anterior
resolvia isso com uma coluna \cd{senha} na planilha --- arquivo que circula por
e-mail, WhatsApp e \emph{pen drive}, e que fica no computador da secretaria.
\end{especificacao}

\textbf{Decisão.} A planilha \textbf{não tem coluna de senha}. A senha inicial é
escolhida na tela de importação, em dois modos: uma para o lote inteiro, ou uma
por pessoa (recomendado para estudantes). Quem nasce da importação recebe
\cd{senha\_provisoria = true}, e o primeiro acesso cai em \rota{/trocar-senha}
sem saída. Reimportar \textbf{não} redefine a senha de quem já entrou.

\textbf{Alternativas descartadas.}
\begin{itemize}
  \item \emph{Manter a senha na planilha}: é o vazamento por construção, e a
    reimportação passaria a redefinir a senha de quem já trocou.
  \item \emph{Delegar tudo ao WSO2 agora}: o piloto precisa de contas antes do
    provedor de identidade existir, e a plataforma continuaria sem caminho para a
    escola sem federação.
  \item \emph{Link mágico por e-mail}: a maioria dos estudantes da rede pública não
    tem e-mail institucional em uso, e a plataforma não envia e-mail hoje.
\end{itemize}

\textbf{Consequências.}
\begin{itemize}
  \item A trava é \textbf{tripla}: \emph{middleware} (páginas), BFF (dados) e
    \rota{/api/exp-token}. Só com o \emph{middleware}, bastaria digitar
    \rota{/exp/...} para pular a troca.
  \item Ao trocar, a sessão é \textbf{refeita} com a senha nova, não remendada:
    o \emph{cookie} antigo ainda diz \enquote{precisa trocar}, e sem refazê-lo o
    \emph{middleware} devolveria a pessoa à mesma tela --- agora com a senha
    anterior já inválida. É o beco sem saída mais fácil de criar neste fluxo.
  \item As senhas geradas são devolvidas \textbf{uma vez}, para impressão. Não
    ficam em coluna nem em \emph{log}; depois disso o caminho é o reajuste de
    senha pelo gestor, que remarca a conta como provisória.
  \item \emph{Ruim, e dito na própria tela}: com uma senha só para o lote,
    qualquer estudante entra na conta de outro até o primeiro acesso dele. Por isso o
    modo \enquote{uma por pessoa} existe e é o recomendado para turma grande.
  \item O \emph{seed} de desenvolvimento nasce com \cd{senha\_provisoria = false}:
    prender os usuários de demonstração na troca pararia os seis
    suítes de ponta a ponta na mesma tela.
\end{itemize}

\textbf{Gatilho de reavaliação.} Entrada do WSO2 em produção --- a partir daí a
senha local vira caminho de exceção (escola sem federação), não o principal.

% ---------------------------------------------------------------------------
\section{ADR-30 --- A rede: lotação por escola e escola ativa no \emph{token}}
\label{adr:rede-lotacoes}

\begin{especificacao}{Contexto}
Vender a uma rede estadual não é vender a uma escola grande. São milhares de
unidades, professores que dão aula em mais de uma, e crianças que se transferem
entre elas. O modelo tinha \cd{users.escola\_id} \textbf{único}, e ele é a chave
de \emph{tenancy}: circula no \emph{claim} \cd{escola\_id} do JWT e no
\cd{set\_config} da RLS. Um professor em duas escolas não caberia nele.

O tamanho do problema, medido: \textbf{38 tabelas} carregam \cd{escola\_id} e
\textbf{23 arquivos} da API derivam a escola do \emph{token}. (Números de
2026-09-05; eram 31 e 19 quando a decisão foi tomada. A diferença é o custo que
a alternativa descartada teria hoje --- ele só cresce, e é por isso que a
medida é refeita em vez de citada.)
\end{especificacao}

\textbf{Decisão.} O \emph{token} continua carregando \textbf{UMA} escola. Nenhuma
das 38 tabelas muda, nenhum dos 23 arquivos muda, a RLS continua sendo
\cd{escola\_id = current\_setting('app.escola\_id')} e a regra inviolável nº 1
continua literalmente verdadeira. O que entra é a camada \emph{acima}:

\begin{itemize}
  \item \cd{lotacoes} --- o vínculo PESSOA $\times$ ESCOLA, datado.
    \enquote{Lotação} é o termo que a rede pública usa para a alocação de um
    servidor a uma unidade. Professor e gestor \textbf{acumulam}: lotação nova não
    encerra a anterior. Estudante tem \textbf{uma ativa}, garantida por índice único
    parcial no banco.
  \item \textbf{Escola ativa na sessão.} Quem tem mais de uma escolhe no
    cabeçalho; a troca passa pelo servidor (\rota{POST /api/escola-ativa}), que
    valida contra \rota{GET /escolas/minhas} --- \textbf{nunca} contra a lista que
    o \emph{cookie} carrega. \emph{Cookie} é dado do cliente, e aceitar a escola
    que ele pedir deixaria qualquer professor ler qualquer unidade da rede.
  \item \cd{users.escola\_id} passa a ser \textbf{cache} da lotação principal ---
    o mesmo papel que \cd{turma\_id} já tinha em relação a \cd{matriculas}.
\end{itemize}

\textbf{Alternativas descartadas.}
\begin{itemize}
  \item \emph{\emph{Token} com várias escolas e RLS em \cd{= ANY(...)}}: tocaria
    as 38 tabelas e trocaria a garantia \enquote{esta conexão vê UMA escola} por
    \enquote{vê um conjunto}. A propriedade que sustenta o isolamento de um
    contrato B2G é justamente a primeira.
  \item \emph{Um usuário por escola (dois cadastros para o mesmo professor)}: é o
    que acontece hoje sem esta camada, e é o problema --- duas senhas, dois
    logins, e nada ligando as duas metades da mesma pessoa.
  \item \emph{\cd{lotacoes} com RLS}: ela é lida no LOGIN, antes de se saber a
    escola. Uma política por \cd{escola\_id} esconderia justamente a resposta
    (\enquote{quais escolas são suas}), e o sintoma seria o professor de duas
    unidades não conseguindo trocar para a segunda. Fica fora, como
    \cd{users}/\cd{turmas}, com filtro explícito nos \emph{services}.
\end{itemize}

\textbf{Consequências.}
\begin{itemize}
  \item Criar escola é \cd{admin-contrato} e ninguém mais --- inclusive no
    \cd{orgs.csv}. Escola é \emph{tenant}: abrir inquilino com RLS, dados de
    menores e \emph{stack} por trás não é operação de quem administra uma unidade.
  \item \emph{Ruim}: a escola ativa é estado de SESSÃO. Duas abas do mesmo
    professor em escolas diferentes brigam pelo mesmo \emph{cookie} --- a última
    troca vence, e a outra aba passa a mostrar a escola nova sem avisar.
  \item \emph{Ruim}: o campo extra que eu pus dentro do item da \cd{SideBar} para
    filtrar o menu virou atributo no DOM (o tipo do DS tem índice aberto e ele
    espalha o que não conhece). A restrição foi para fora do item; fica
    registrado porque é o tipo de erro que o \emph{typecheck} não pega.
\end{itemize}

\textbf{Gatilho de reavaliação.} Um contrato em que a MESMA pessoa precise ver
duas escolas ao mesmo tempo na mesma tela --- aí a conversa volta para o
\emph{token} com conjunto de escolas, e com ela o custo nas 38 tabelas.

% ---------------------------------------------------------------------------
\section{ADR-31 --- CPF como identidade, guardado só em \emph{hash}}
\label{adr:cpf}

\begin{atencao}
\textbf{SUPERADA pela \adrr{ADR-33}.} O \emph{hash} entregava o que a
\emph{plataforma} precisa do CPF e não o que o \emph{gestor} precisa: ver o
número. Fica registrada porque o raciocínio continua valendo para qualquer dado
que a plataforma só precise \emph{comparar} --- e porque a lição de processo é a
parte mais útil dela.
\end{atencao}

\begin{especificacao}{Contexto}
Numa rede, a criança troca de escola e o dado dela precisa continuar sendo dela.
Nome muda (casamento dos pais, correção de registro), \emph{login} muda,
matrícula muda a cada rede --- e a escola de origem e a de destino não
compartilham nenhuma dessas chaves. Sem um identificador imutável, cada
transferência cria uma segunda pessoa e o histórico nasce partido em duas.

O identificador que o Brasil usa é o CPF, e é o que a Portaria MEC 269/2026
adota como chave do estudante no EducaDados. É também o dado pessoal mais
sensível que esta plataforma tocaria, de menores de idade, num banco que já
guarda perfil de valores e \emph{check-in} de saúde.
\end{especificacao}

\textbf{Decisão.} O CPF é a chave de identidade e \textbf{nunca} é guardado. O
banco tem \cd{users.cpf\_hash} = \cd{HMAC-SHA256(cpf, CPF\_PEPPER)} e
\cd{users.cpf\_final}, os três últimos dígitos em claro para confirmação humana
na tela. O \emph{hash} é \textbf{determinístico} de propósito: é o que permite
procurar (basta \emph{hashear} o que o gestor digitou) e reconhecer a mesma
pessoa. Único por \textbf{contrato}, não por escola --- é a identidade da pessoa
na rede inteira.

\textbf{Alternativas descartadas.}
\begin{itemize}
  \item \emph{Guardar o número em claro}: dá tudo o que o \emph{hash} dá e mais
    um acervo de CPF de menores, que a plataforma não precisa para nada --- ela
    não emite documento nem presta contas ao fisco. O SIS da rede já tem o número.
  \item \emph{\emph{Hash} com sal por linha} (\cd{bcrypt}): inutilizaria o campo
    como identidade. Sem determinismo não há busca nem reconhecimento, que é a
    única razão de o CPF estar aqui.
  \item \emph{\emph{Hash} sem pimenta}: CPF tem 11 dígitos e espaço de busca
    trivial --- um SHA-256 puro é reversível por força bruta em minutos. A pimenta
    é o que torna o \emph{hash} inútil fora desta instalação.
  \item \emph{Cifrar em vez de \emph{hashear}}: exigiria guardar a chave capaz de
    devolver o número, e o número é exatamente o que não se quer poder devolver.
\end{itemize}

\textbf{Consequências.}
\begin{itemize}
  \item \env{CPF\_PEPPER} \textbf{não muda}. Trocá-la invalida todos os
    \emph{hashes} e a plataforma deixa de reconhecer quem já cadastrou --- o
    sintoma é toda criança transferida virando cadastro novo. Ausente, a coluna
    \cd{cpf} da importação é \textbf{recusada} com mensagem explícita, nunca
    aceita e guardada de outro jeito.
  \item O dígito verificador é conferido na leitura da planilha. Não é
    preciosismo: CPF digitado errado numa planilha de 10 mil linhas cria uma
    identidade fantasma que nunca vai casar com a da rede, e o erro só aparece
    meses depois, na primeira transferência.
  \item O CPF é \textbf{opcional}. Rede sem o dado continua funcionando com
    \cd{matricula\_externa} como chave --- e perde só o reconhecimento
    automático na transferência.
  \item \emph{Ruim}: três dígitos em claro são identificáveis em combinação com o
    nome. É o preço de o gestor conseguir confirmar que é a pessoa certa antes de
    transferir uma criança.
  \item O teste estrutural \arq{identidade-unica.spec.ts} falha se aparecer
    qualquer coluna com \cd{cpf} no nome fora de \cd{users.cpf\_hash} e
    \cd{users.cpf\_final}.
\end{itemize}

\textbf{Gatilho de reavaliação.} A especificação técnica do EducaDados exigir
\emph{enviar} o CPF de volta à rede --- aí ele volta a entrar, e a decisão passa a
ser onde a chave de cifra mora, não se o número existe.

% ---------------------------------------------------------------------------
\section{ADR-32 --- Transferência: o dado fica, a identidade segue}
\label{adr:transferencia}

\begin{especificacao}{Contexto}
\enquote{Os dados do estudante seguem o estudante, independente da escola} é requisito de
produto. E colide de frente com a regra inviolável nº 1: tudo o que a criança
produziu tem \cd{escola\_id} da escola onde produziu, com RLS por cima. Depois de
uma transferência, a escola nova não alcança nada disso --- e a antiga continua
alcançando tudo.
\end{especificacao}

\textbf{Decisão.} O \textbf{dado fica onde foi produzido}; o que segue o estudante é a
\textbf{identidade}. A transferência encerra a lotação e as matrículas na origem,
abre na destino e move o \emph{cache} \cd{users.escola\_id} --- e mais nada. Para
a leitura atravessada existe \textbf{uma} porta: \rota{GET /estudantes/:id/historico},
que devolve \textbf{participações} (em que escola, em que turma, em que missão,
quando), nunca o conteúdo dos instrumentos.

\textbf{Alternativas descartadas.}
\begin{itemize}
  \item \emph{Reescrever \cd{escola\_id} nas 38 tabelas na transferência}: apaga
    de onde o fato veio. \enquote{Esta criança respondeu isso na escola A, em
    2026} deixaria de ser recuperável, e um \emph{UPDATE} em massa por
    transferência é a operação mais fácil de deixar pela metade.
  \item \emph{Afrouxar a RLS para o \cd{estudante\_id} atravessar}: transformaria a
    exceção numa regra geral. A porta única é auditável; uma política mais larga
    não é.
  \item \emph{O histórico devolver o instrumento completo}: abriria, por uma tela
    de secretaria, o retrato de valores da Vila e a transcrição do unooGPT --- o
    que \adrr{ADR-12}, \adrr{ADR-27} e \adrr{ADR-10} fecharam. Quem lê o
    instrumento é quem tem a posse pedagógica na escola onde ele aconteceu.
\end{itemize}

\textbf{Consequências.}
\begin{itemize}
  \item O \cd{historico} só abre para quem tem \textbf{lotação na escola onde o
    estudante está hoje} --- é a unidade responsável por ele --- ou para
    \cd{admin-contrato}.
  \item A transferência é a única operação que troca o \cd{app.escola\_id}
    \textbf{no meio de uma transação}: ela encerra na origem e abre no destino, e
    sem trocar o contexto o \emph{update} na origem não veria linha nenhuma (a RLS
    o esconderia, sem erro, e a criança ficaria matriculada nas duas escolas). O
    \cd{set\_config(..., true)} é local à transação, então a conexão volta limpa
    ao \emph{pool}.
  \item Quem recebe a criança é a escola de \textbf{destino}: um gestor só
    transfere PARA a própria escola, e o caminho normal é importar a planilha com
    o CPF --- o importador reconhece e transfere.
  \item \emph{Ruim}: o professor da escola nova vê que a criança participou de
    uma missão na escola anterior, mas não o que ela respondeu. Para continuar um
    trabalho pedagógico iniciado em outra unidade, isso é pouco --- e é
    deliberado até haver uma regra de consentimento que diga o contrário.
\end{itemize}

\textbf{Gatilho de reavaliação.} Pedido pedagógico formal de continuidade entre
escolas da mesma rede --- aí a conversa é sobre consentimento e sobre QUAL
instrumento atravessa, um a um, nunca sobre abrir a fronteira inteira.

% ---------------------------------------------------------------------------
\section{ADR-33 --- O CPF é guardado, e o gestor o vê}
\label{adr:cpf-em-claro}

\begin{especificacao}{Contexto}
\adrr{ADR-31} decidiu guardar apenas \cd{HMAC-SHA256(cpf, CPF\_PEPPER)} e os três
últimos dígitos. A decisão pesou o risco de o banco acumular CPF de menores e
\textbf{não perguntou quem usa o dado}.

Quem usa é o gestor da escola, e ele precisa \textbf{ver o número}: confere
identidade contra o sistema de gestão da rede, preenche documento de
transferência e atende a família no balcão com o documento na mão. \emph{Hash} é
de mão única --- não existe configuração que o exiba, porque não existe chave.
Não era uma limitação a ajustar; era a funcionalidade ausente por construção.
\end{especificacao}

\textbf{Decisão.} O CPF é guardado: \cd{users.cpf}, os 11 dígitos sem pontuação
(a formatação é da tela, e um \emph{CHECK} no banco impede que
\enquote{529.982.247-25} e \enquote{52998224725} coexistam como duas pessoas).
Único por \textbf{contrato}, não por escola. Exibido formatado na lista do gestor,
buscável com ou sem pontuação, e corrigível na edição --- o balcão erra, e um
dígito trocado numa planilha de dez mil linhas só aparece meses depois, na
primeira transferência.

\textbf{Sobre a base legal.} CPF é dado pessoal \textbf{comum}, não dado sensível
do art.~11 da LGPD --- que trata de saúde, raça, religião, filiação e biometria.
O tratamento tem base clara num contrato de execução de política pública de
educação, e é o que todo sistema de gestão escolar brasileiro guarda. O cuidado
que a plataforma tem com \cd{exp\_91\_checkins} (saúde) não é o mesmo caso.

\textbf{Alternativas descartadas.}
\begin{itemize}
  \item \emph{Manter o \emph{hash} e exibir os três dígitos}: é o que havia, e não
    atende --- três dígitos não servem para conferir documento nem para preencher
    formulário.
  \item \emph{\emph{Hash} para identidade $+$ coluna cifrada para exibição}: duas
    representações do mesmo fato, uma chave a operar e a rotacionar, e um vazamento
    de \emph{dump} $+$ chave devolve o número de todo modo. O ganho real sobre a
    cifra de disco que o Cloud SQL já aplica é pequeno; a complexidade é
    permanente.
  \item \emph{Exibir mascarado na lista e completo só no cadastro}: foi
    considerado e recusado pelo dono do produto --- o gestor confere pela lista, e
    obrigá-lo a abrir cada pessoa transforma uma conferência de turma em trinta
    cliques.
\end{itemize}

\textbf{Consequências.}
\begin{itemize}
  \item \textbf{O que sustenta a decisão não é a coluna, é a contenção.}
    \cd{users.cpf} é o ÚNICO lugar onde o número existe: nunca em \emph{log} (nem
    em mensagem de erro), nunca em telemetria, JSONB, \emph{claim} de \emph{token}
    ou \emph{bundle} de experiência. O teste estrutural cobra pelos dois lados ---
    que só exista essa coluna, e que nenhum JSONB de \cd{telemetria\_eventos} ou
    \cd{progresso} carregue onze dígitos seguidos, que é o caminho mais fácil de o
    número escapar.
  \item O índice único é da \textbf{rede}, não da escola. Duas importações
    simultâneas com o mesmo CPF colidem no banco --- é o comportamento correto, e
    o efeito colateral apareceu primeiro nos testes, onde duas suítes usavam o
    mesmo número.
  \item \env{CPF\_PEPPER} deixou de existir.
  \item \emph{Ruim}: o banco passa a ser um acervo de CPF de menores de idade.
    A proteção é de acesso (só gestor e admin-contrato alcançam a rota) e de
    contenção, não de criptografia de aplicação.
  \item \emph{Lição de processo, registrada porque custou retrabalho}: uma escolha
    \textbf{irreversível} --- \emph{hash} não volta atrás --- foi tomada sem o
    requisito de quem usa o dado estar escrito. A pergunta \enquote{quem precisa
    LER isso, e como?} devia ter vindo antes da primeira migração.
\end{itemize}

\textbf{Gatilho de reavaliação.} Exigência contratual de cifra em nível de
aplicação, ou um incidente de acesso indevido --- aí a conversa é sobre onde a
chave mora e quem a usa, não sobre se o número existe.

\section{ADR-34 --- A regional agrupa; quem autoriza é a lotação}
\label{adr:nivel-regional}

\begin{especificacao}{Contexto}
Uma rede estadual não chega organizada em escolas soltas. Em São Paulo são
\textbf{Diretorias de Ensino}; em outras redes, polos ou secretarias municipais.
Entre a escola e o contrato existe um nível intermediário real, com gente que
trabalha nele: o supervisor que acompanha as cinquenta unidades da sua diretoria.

A \adrr{ADR-30} deixou a plataforma pronta para a pessoa em várias escolas --- a
lotação existe, o \emph{token} carrega uma escola por vez e o cabeçalho troca a
escola ativa. O que faltava era o \emph{agrupamento}: sem ele, pôr um supervisor
nas cinquenta unidades é abrir cinquenta telas.

As colunas \cd{escolas.tipo} e \cd{escolas.pai\_id} tinham sido acrescentadas
antes, \enquote{porque custam zero}, e \textbf{nenhuma consulta as lia}. Era
exatamente a armadilha que este projeto já documentou duas vezes --- o
\cd{expires\_at} sem rotina e o \cd{turmas.consentimento\_em} do aplicativo de
origem: a regra sai do código, a coluna fica, e quem lê o \emph{schema} acredita
nela.
\end{especificacao}

\textbf{Decisão.} O nível regional passa a ser \textbf{real e lido}, com um
vocabulário fechado de três valores em \cd{escolas.tipo} --- \cd{escola} (a
unidade, onde há turma, estudante e dado), \cd{regional} (o que agrupa unidades) e
\cd{rede} (a raiz). \cd{pai\_id} liga unidade a regional, e um \emph{trigger} no
banco recusa \cd{pai\_id} apontando para uma \cd{escola} ou para si mesma.

E a decisão que dá sentido a tudo: \textbf{a hierarquia agrupa e nunca autoriza}.
Estar numa diretoria não dá acesso a nada. Quem dá acesso é a lotação, uma por
unidade. A regional é o \emph{insumo} da lotação em lote
(\rota{POST /escolas/lotacoes-em-lote}) --- ela monta a lista de escolas que o
admin do contrato confirma antes de conceder.

\textbf{Alternativas descartadas.}
\begin{itemize}
  \item \emph{Um papel novo, \cd{supervisor}}: \cd{PAPEIS} é contrato congelado
    entre o \emph{middleware} do \emph{hub}, os \emph{guards} da api e a RLS; um
    quinto papel toca os três. E não resolveria nada que a lotação já não
    resolva --- o supervisor faz o que um gestor faz, em mais de uma escola.
  \item \emph{Autorizar por hierarquia} (\enquote{quem está na regional lê as
    unidades dela}): é uma segunda fonte de verdade sobre acesso, ao lado da
    lotação, e as duas divergem no dia em que alguém corrige um \cd{pai\_id}.
    Pior: transformaria uma coluna \emph{organizacional} em coluna de segurança,
    e a RLS deixaria de ser legível numa linha.
  \item \emph{\cd{regional\_id} no \emph{token}, com a api somando as escolas
    da diretoria}: quebraria a regra 1 em todas as 38 tabelas com \cd{escola\_id}
    e daria ao supervisor uma tela que mistura duas escolas --- que é
    precisamente o que o isolamento lógico existe para impedir.
  \item \emph{Aceitar o critério no lugar da lista} (\enquote{lote todos da
    regional X}): o servidor expandiria o critério e o humano confirmaria um
    número que não viu. A rota exige a \textbf{lista} de \cd{escolaIds}, no
    máximo mil.
\end{itemize}

\textbf{Consequências.}
\begin{itemize}
  \item \textbf{Nada mudou na tenancy.} O \emph{token} continua carregando uma
    escola, a política de RLS continua sendo a mesma linha, e as 38 tabelas com
    \cd{escola\_id} não foram tocadas. O nível regional inteiro vive acima delas.
  \item A regional \textbf{não recebe lotação}: \cd{lotarEmLote} ignora o que não
    for \cd{tipo = 'escola'} e devolve a lista de ignoradas. Lotar alguém numa
    regional daria uma escola sem turma nem estudante no seletor do cabeçalho.
  \item O panorama (\rota{GET /escolas/minhas/resumo}, tela
    \cd{/admin/minhas-escolas}) \textbf{lista, não soma}: uma linha por unidade,
    com turmas, estudantes e quantos ainda não fizeram o primeiro acesso. Somar as
    escolas numa métrica só seria a tela que a decisão acima recusa.
  \item O \cd{orgs.csv} do OneRoster passa a criar regionais a partir de
    \cd{district}/\cd{local} e a ligar unidades por \cd{parentSourcedId}. Pai
    desconhecido vira \emph{aviso} e a escola entra sem regional --- perder a
    escola por causa do agrupamento seria trocar o essencial pelo acessório.
  \item \emph{Ruim}: a lotação em lote é uma operação de poucos cliques que
    concede acesso a muitas escolas. A contenção é ser exclusiva de
    \cd{admin-contrato}, exigir a lista explícita e mostrar os nomes antes de
    confirmar.
  \item \emph{Lição de processo}: coluna acrescentada \enquote{porque custa zero}
    e nunca lida é dívida com aparência de recurso. Ou nasce com o consumidor,
    ou não nasce.
\end{itemize}

\textbf{Gatilho de reavaliação.} Uma rede que exija leitura agregada de fato
--- \enquote{o desempenho da diretoria inteira} --- muda a pergunta: aí é
relatório com escopo próprio (como \cd{/vila/agregado/*} faz para o gestor),
nunca alargamento do \emph{token}.

\section{ADR-35 --- Perfil portátil: um documento assinado, não uma integração}
\label{adr:perfil-portatil}

\begin{especificacao}{Contexto}
Um contrato é uma rede e uma \emph{stack}. A \adrr{ADR-32} resolveu a criança
que troca de escola DENTRO da rede: o dado fica onde foi produzido e a leitura
atravessa por uma porta só. Fica de fora o caso que o dono do produto levantou:
a criança que sai da rede estadual para a municipal --- outro contrato, outro
banco, nenhum \emph{JOIN} possível.

E há a restrição que decide o desenho: \textbf{não há garantia nenhuma de que a
rede de destino comprou o produto}. Qualquer solução que exija duas instalações
desta plataforma conversando resolve só o caso raro em que as duas existem.
\end{especificacao}

\textbf{Decisão.} A saída é um \textbf{documento assinado}, e não uma integração.
Um arquivo JSON \cd{\{ perfil, assinatura \}} que a escola de origem emite a
pedido do titular (LGPD art.~18, V --- portabilidade) e que a escola de destino
importa como um ato dela. Para uma rede que nunca vai instalar esta plataforma, o
mesmo arquivo continua sendo um documento legível.

Quatro decisões dentro dessa:

\begin{itemize}
  \item \textbf{O hub assina; a api verifica.} Assinar exige a chave privada, que
    só o hub tem --- é a MESMA que assina os \emph{tokens} de experiência e é
    publicada em \rota{/.well-known/jwks.json}, então nenhum material de chave
    novo entrou. \cd{aud: "perfil-portatil"} separa os dois usos. Verificar
    acontece na api porque é o passo que \textbf{não pode confiar em quem chama}:
    se a api recebesse o veredito pronto do hub, bastaria forjar o campo.
  \item \textbf{O que vem de fora nunca se mistura.} O documento vive em
    \cd{perfis\_portateis} e em lugar nenhum mais: não vira matrícula, progresso
    nem perfil da Vila, e não entra em média. É a afirmação de OUTRA instalação
    sobre a criança, possivelmente de outra \cd{versao\_mapeamento}. Separado, o
    professor lê os dois e sabe qual é qual.
  \item \textbf{Assinatura inválida é recusa; emissor inalcançável é rótulo.} São
    coisas diferentes e a tela diz qual foi. Não conseguir falar com o emissor é
    limitação de rede; assinatura que não confere é adulteração, e essa não entra.
  \item \textbf{O documento se liga a alguém que já está aqui.} Importar não cria
    estudante: a matrícula é ato da escola (planilha ou OneRoster), e só depois o
    perfil se anexa --- pelo CPF, que é a única chave que significa alguma coisa
    numa base onde nenhum id da origem existe.
\end{itemize}

\textbf{Alternativas descartadas.}
\begin{itemize}
  \item \emph{API entre as duas instalações}: exige que as duas existam, que se
    conheçam e que uma libere acesso à outra. Não cobre o caso comum, que é a
    rede de destino não ser cliente.
  \item \emph{Um banco só para várias redes} (\cd{rede\_id} em tudo): foi medido
    e descartado antes --- e não resolveria este caso, porque a rede de destino
    continuaria fora da instalação.
  \item \emph{Exportar CSV sem assinatura}: quem recebe não tem como saber se o
    arquivo foi editado no caminho, e um histórico escolar sem procedência não
    serve para decidir nada.
  \item \emph{Assinar o documento como \emph{payload} do JWS}: o arquivo viraria
    uma \emph{string} opaca. Ele é entregue à FAMÍLIA, que tem direito de lê-lo;
    por isso o JWS cobre o \emph{hash} canônico e o documento fica legível ao
    lado.
  \item \emph{Buscar o JWKS de qualquer emissor}: o \cd{emissor.hub} vem de
    dentro do arquivo, ou seja, de um desconhecido --- seria SSRF. A allowlist
    \env{PORTABILIDADE\_EMISSORES} é a decisão de confiança: o documento
    \textbf{seleciona} entre origens que o operador escreveu, e nunca introduz
    uma URL.
\end{itemize}

\textbf{Consequências.}
\begin{itemize}
  \item \textbf{O CPF passa a existir num segundo lugar}, e é a única exceção à
    contenção da \adrr{ADR-33}: \cd{perfis\_portateis.documento}. O CPF é o
    CONTEÚDO do documento --- é o que reconhece a criança lá fora --- e é
    guardado \emph{verbatim} porque alterá-lo destruiria a assinatura. A exceção
    está numa lista explícita dentro do teste estrutural, para que a próxima seja
    uma decisão e não um vazamento que ninguém viu.
  \item \textbf{Bem-estar (exp-91) nunca entra no documento.} DASS-21 é dado de
    saúde, a posse é do professor e não existe rota de gestor para ele em lugar
    nenhum. Um documento que o gestor emite não pode carregar o que nem o gestor
    lê.
  \item \env{PORTABILIDADE\_EMISSORES} vazia (o default, e o caso de dev) faz
    TODO documento entrar como \cd{nao-verificada}. Degradar visivelmente é
    melhor que recusar o direito do titular --- e a tela sai em âmbar, não em
    verde.
  \item \textbf{A entrada da allowlist pode apontar o JWKS}
    (\cd{origem=urlDoJwks}), e é isso que torna o caminho criptográfico
    exercitável fora de produção. O \cd{emissor.hub} do documento é o endereço do
    \textbf{navegador} --- a mesma distinção de \env{EXP\_INTERNAL\_HOST}, e pela
    mesma razão: de dentro do contêiner da api, \cd{localhost:3000} é a própria
    api. A regra de canal passou a valer sobre a URL do JWKS: \cd{https} sempre,
    exceto \textbf{fora de produção}, onde \cd{http} é aceito e avisado no
    \emph{log} (quem escreveu a entrada foi o operador, no servidor, e o tráfego
    não sai da máquina). Em produção, \cd{http} é recusado --- a chave pública
    viria por um canal que qualquer um reescreve.
  \item Sem isso, a garantia mais importante do recurso --- documento adulterado
    é RECUSADO --- não teria prova no ambiente onde se trabalha todo dia, e o
    \emph{smoke} visual registrava o caminho degradado como se fosse o normal.
  \item O caminho criptográfico é coberto por teste com as chaves injetadas, e
    foi esse teste que pegou o pior erro possível deste recurso: a classificação
    do erro olhava o TEXTO da exceção, e a palavra \enquote{failed} de
    \enquote{signature verification failed} caía na regra de erro de rede.
    Documento FORJADO entraria marcado \enquote{não deu para conferir} em vez de
    recusado. Hoje a classificação é pelo \cd{code} da biblioteca.
  \item \emph{Ruim}: reemitir gera outro documento, e a escola de destino pode
    acabar com dois do mesmo estudante. A lista mostra data \emph{e hora} e permite
    apagar; unificar por conteúdo esconderia que foram duas emissões.
\end{itemize}

\textbf{Gatilho de reavaliação.} A especificação técnica do EducaDados (Portaria
MEC 269/2026) definir um formato nacional de transferência de histórico: aí este
documento vira um \emph{binding} daquele formato, e o que sobrevive é a decisão
de ser documento assinado em vez de integração.

\section{ADR-36 --- O OpenAPI é derivado do Zod, não escrito ao lado dele}
\label{adr:openapi-derivado}

\begin{especificacao}{Contexto}
A api tinha 146 rotas em 27 \emph{controllers} e nenhuma documentação executável
(hoje são 170 operações em 32 \emph{controllers} --- o inventário cresce, a
decisão não muda):
quem quisesse saber o que \rota{POST /rostering/importar} aceita lia o
\emph{controller}. A única descrição existente é o Apêndice~\ref{ap:endpoints}
deste documento --- prosa escrita à mão, que nada obriga a concordar com o
código.

A complicação está em como a api valida. Não há DTO de classe: são 42
\emph{schemas} Zod declarados como \cd{const} dentro dos \emph{controllers} e
aplicados por \cd{.parse()}/\cd{.safeParse()} no corpo do método, sem
\emph{pipe} compartilhado. O \cd{@nestjs/swagger} infere esquema de classe
anotada com \cd{class-validator}; sobre Zod ele não infere nada.
\end{especificacao}

\textbf{Decisão.} O documento é \textbf{derivado} do \emph{schema} que a rota já
usa para validar. Um decorator único, \cd{@Contrato}, declara resumo, prosa,
corpo, \emph{query} e \textbf{resposta}; \cd{zod-to-json-schema} traduz para o
dialeto do OpenAPI 3.0; e \cd{openapi/documento.ts} monta cada operação.

Três decisões dentro dessa:

\begin{itemize}
  \item \textbf{O \cd{@Contrato} não chama \cd{@ApiOperation}/\cd{@ApiResponse}.}
    Ele só grava metadata. Um decorator não enxerga o verbo da rota --- o
    \cd{@Get}/\cd{@Post} pode ser aplicado depois dele, e a ordem de avaliação é
    de baixo para cima --- então teria de ADIVINHAR o código de sucesso, e
    também não enxerga o \cd{@Roles}. Quem monta a operação é o construtor do
    documento, que percorre as rotas já registradas e tem o verbo, o
    \cd{@HttpCode} e o \cd{@Roles} à mão.
  \item \textbf{Os papéis vêm do \cd{@Roles} real}, nunca de um campo repetido no
    \cd{@Contrato}: repetir abriria espaço para o documento dizer
    \cd{gestor} enquanto o \emph{guard} exige \cd{admin-contrato}.
  \item \textbf{O esqueleto de caminhos vem do roteador do Nest}
    (\cd{SwaggerModule.createDocument}), que é a única fonte incapaz de inventar
    uma rota que não existe. As duas travessias --- a do roteador e a do
    \cd{DiscoveryService} --- precisam concordar, e divergência de um caminho
    sequer é erro alto na subida.
\end{itemize}

\textbf{Alternativas descartadas.}
\begin{itemize}
  \item \emph{Anotar \cd{@ApiBody}/\cd{@ApiResponse} à mão nas 146 rotas}: criaria
    uma segunda descrição do contrato ao lado da primeira, sem nada que as
    obrigue a concordar. É literalmente a armadilha do \cd{expires\_at} que ficou
    no \emph{schema} depois que a rotina saiu do código.
  \item \emph{Migrar a validação para DTOs com \cd{class-validator}}: reescreveria
    a validação das 146 rotas para ganhar inferência automática, e trocaria o Zod
    --- que a convenção da casa usa em TODO input externo, \emph{env} incluído ---
    por um segundo validador só nas rotas.
  \item \emph{\cd{nestjs-zod}}: exige DTO de classe (\cd{createZodDto}) e
    \emph{pipe} global. É a mesma reescrita, com uma dependência a mais.
  \item \emph{\cd{@asteasolutions/zod-to-openapi}}: monta o documento a partir de
    um registro próprio, fora do roteador do Nest --- ou seja, pode listar rota
    que não existe e omitir rota que existe, que é exatamente o defeito que se
    quer evitar.
\end{itemize}

\textbf{Consequências.}
\begin{itemize}
  \item \textbf{\cd{.refine()} some da conversão, em silêncio.} JSON Schema não
    expressa função JavaScript, e a biblioteca descarta. O
    \cd{AjustarMatriculaSchema} exige \enquote{\cd{turmaDestinoId} OU
    \cd{principal}} e isso não aparece como esquema. Por isso o \cd{@Contrato}
    tem o campo \cd{notas}: restrição que o Zod checa e o JSON Schema não sabe
    dizer vira texto na descrição. A validação não muda --- o que as notas
    consertam é o documento não mentir por omissão.
  \item \textbf{Esquemas saem embutidos, sem \cd{\$ref}}
    (\cd{\$refStrategy: "none"}). É seguro porque não há \cd{z.lazy} em lugar
    nenhum da api. Custa repetição no JSON e ganha um documento que não depende
    de um registro de nomes estáveis.
  \item \emph{Ruim}: a assinatura publicada do \cd{zodToJsonSchema} estoura o
    compilador com \textbf{TS2589} (\enquote{type instantiation is excessively
    deep}) --- o tipo de retorno instancia \cd{JsonSchema7Type}, uma união
    recursiva grande, mesmo quando \cd{target: "openApi3"} resolve para
    \cd{object}. O módulo fixa localmente a assinatura que usa. Some quando a
    biblioteca simplificar o retorno.
  \item \textbf{A entrega foi em etapas, com lista que só encolhe.} Documentar
    146 rotas num commit seria um \emph{diff} que ninguém revisa, e uma resposta
    mal descrita passaria batida. \cd{tests/openapi-cobertura.spec.ts} percorre o
    roteador e cobra \cd{@Contrato} de toda rota, salvo as nomeadas numa lista
    explícita --- o mesmo molde do \cd{SEM\_CHAVE\_POR\_DECISAO} de
    \cd{identidade-unica.spec.ts}, com a diferença de que aquelas exceções são
    permanentes e estas eram dívida. \textbf{A lista está vazia}: todas as rotas
    declaram contrato (170 operações em 2026-09-05), e rota nova sem ele é teste
    vermelho.
  \item \textbf{O Apêndice~\ref{ap:endpoints} deixou de ser escrito à mão.} Ele
    é gerado do \arq{docs/openapi.json}, que é gerado do \cd{@Contrato}. Eram
    547 linhas de prosa descrevendo 146 rotas sem que nada obrigasse essa prosa a
    concordar com o código --- exatamente o defeito que esta ADR existe para
    fechar, e que teria sobrevivido dentro do próprio documento que a registra.
\end{itemize}

\textbf{Gatilho de reavaliação.} Subir para Zod 4, que traz
\cd{z.toJSONSchema} nativo: aí a dependência de conversão sai, e o TS2589
provavelmente com ela.

\section{ADR-37 --- O \emph{schema} de resposta é conferido, não só publicado}
\label{adr:contrato-de-resposta}

\begin{especificacao}{Contexto}
O \emph{schema} de ENTRADA não tem como envelhecer: a api o usa para validar em
toda requisição, então errar nele quebra a rota na hora. O de SAÍDA não tinha
nenhum --- os \emph{services} devolvem entidade do Prisma direto, \emph{spreads}
e formas ad-hoc --- e passa a existir só para o documento.

Afirmação que nada exercita é a mesma coisa que a coluna de prazo sem rotina que
a cumpra: quem lê acredita, e ninguém percebe quando deixa de ser verdade.
\end{especificacao}

\textbf{Decisão.} Três camadas, em ordem de força.

\begin{enumerate}
  \item \textbf{Teste}, com \cd{.parse()} estrito, nas suítes que já exercitam
    rota de verdade. É a conferência que dói.
  \item \textbf{Runtime fora de produção}: o \cd{ContratoInterceptor} roda
    \cd{safeParse} na resposta e avisa no \emph{log}. Nunca lança.
  \item \textbf{Produção}: o interceptor é \emph{no-op}.
\end{enumerate}

\textbf{Alternativas descartadas.}
\begin{itemize}
  \item \emph{Validar a resposta também em produção}: um \emph{schema} estrito
    demais transformaria rota que funciona em 500. O documento não vale uma
    indisponibilidade.
  \item \emph{Só publicar, sem conferir nada}: é o \cd{expires\_at} de novo, agora
    em forma de \emph{schema}.
\end{itemize}

\textbf{Consequências.}
\begin{itemize}
  \item \textbf{O interceptor serializa antes de conferir}
    (\cd{JSON.parse(JSON.stringify(...))}). Ele roda ANTES da serialização, onde
    ainda há \cd{Date} do Prisma; sem essa volta, todo campo de data acusaria
    divergência falsa contra um \cd{z.string().datetime()} correto. É também
    parte do motivo de ele não rodar em produção.
  \item \textbf{O aviso carrega só caminho e código do problema, nunca o valor.}
    A mensagem do Zod embute o recebido em alguns casos
    (\cd{invalid\_enum\_value}), e a resposta desta api carrega nome de criança e
    CPF --- regra de não registrar PII em \emph{log}, nem dentro de mensagem de
    erro.
  \item Funcionou na primeira rota documentada: o \emph{schema} do
    \rota{POST /auth/login} declarava \cd{escolaId} como \cd{uuid}, e
    \cd{Escola.id} é \cd{String @id} sem \emph{default} --- o id vem da rede. O
    aviso apareceu no primeiro \emph{login} real.
\end{itemize}

\textbf{Gatilho de reavaliação.} A conferência aparecer no perfil de latência de
desenvolvimento a ponto de atrapalhar --- aí ela passa a valer só em teste.

\section{ADR-38 --- O \emph{deploy} paraleliza os \emph{jobs}, não o portão}
\label{adr:deploy-paralelo}

\begin{especificacao}{Contexto}
Medido nos \emph{runs} de 2026-08-28, um \emph{push} comum levava \textasciitilde
9,6\,min até a escola e um \cd{forcar}, \textasciitilde 26\,min --- e 80\% disso
eram as duas imagens Docker, por duas causas que se somavam: \cd{COPY . .} antes
do \cd{pnpm install} invalidava a camada de dependências em todo \emph{commit}
(o \emph{cache} GHA só acertava apt e corepack), e api e hub buildavam em série
no mesmo \emph{job} sem depender uma da outra.
\end{especificacao}

\textbf{Decisão.} Duas mudanças, nenhuma nova peça de infra: (1) os
\emph{Dockerfiles} abrem com \cd{COPY pnpm-lock.yaml pnpm-workspace.yaml} +
\cd{pnpm fetch} --- que baixa o \emph{store} lendo só o \emph{lockfile} --- e o
\cd{pnpm install --offline} roda depois do \cd{COPY . .}; (2) o \emph{workflow}
vira quatro \emph{jobs} --- \cd{mudancas} (o \emph{diff}, calculado uma vez),
\cd{deploy-api} e \cd{deploy-hub} em paralelo atrás do portão, e
\cd{publicar-exps} depois dos dois, porque o \cd{PUT} do manifesto é validado
contra o JWKS que o hub serve.

\textbf{Alternativas descartadas.}
\begin{itemize}
  \item \emph{Buildar as imagens em paralelo com o portão}, deixando só
    \cd{gcloud run deploy} e a migração esperarem o verde: ganharia 2--6\,min de
    caminho crítico, mas gastaria \emph{runner} em todo \emph{commit} vermelho, e
    o pipeline passaria a ter dois regimes de \emph{gating} para raciocinar. O
    portão continua na frente de tudo.
  \item \emph{Cache remoto do Turborepo} ligando \cd{verificar}, o \emph{build}
    do \emph{runner} e os \emph{builds} dentro das imagens: é a otimização certa
    quando o \cd{tsc}/\cd{next build} virarem o gargalo --- hoje o gargalo era o
    \cd{pnpm install} repetido, e um serviço de \emph{cache} é uma peça a mais
    para operar.
  \item \emph{Copiar o \arq{dist/} do runner para dentro da imagem} em vez de
    buildar lá dentro: acopla o \emph{Dockerfile} ao \emph{workflow} (um
    \cd{docker build} local deixaria de funcionar sozinho) para economizar um
    passo que, com a camada de dependências cacheada, custa segundos.
\end{itemize}

\textbf{Consequências.}
\begin{itemize}
  \item O hub não espera mais o \emph{deploy} da api. A ordem era acidental ---
    os dois conversam por HTTP e o \emph{rollout} do Cloud Run nunca garantiu
    revisão casada --- mas fica registrado: se um dia um contrato exigir
    hub-depois-da-api, é um \cd{needs} de uma linha.
  \item \cd{publicar-exps} usa \cd{!cancelled()} com \cd{result != 'failure'}
    nos \cd{needs}: os \emph{deploys} podem ter sido PULADOS (um \emph{push} que
    só toca experiência), e pulado não pode bloquear a publicação.
  \item Cada \emph{job} que autentica no GCP repete o \emph{setup} de
    \textasciitilde 1\,min (WIF, gcloud, buildx). É o preço do paralelismo, pago
    em paralelo.
  \item O \cd{pnpm fetch} depende só do \emph{lockfile}: mudou dependência, a
    camada refaz inteira. É o comportamento certo --- e é raro na frequência de
    \emph{commit} deste repo.
\end{itemize}

\textbf{Gatilho de reavaliação.} O \cd{nest build}/\cd{next build} dentro das
imagens passar de \textasciitilde 2\,min com a camada de dependências já em
\emph{cache} --- aí o cache remoto do Turborepo entra na conversa.

\section{ADR-39 --- Vulnerabilidade de dependência trava o portão, com prazo para os \emph{majors}}
\label{adr:snyk-no-portao}

\begin{atencao}
\textbf{O portão descrito acima NÃO está ativo (situação em 2026-09-05).} O
passo \cd{Snyk} do \arq{ci.yml} está \textbf{comentado} desde o \emph{commit}
\cd{07d9ec0} (\enquote{\emph{chore: disable Snyk vulnerability scan in CI
workflow}}). Enquanto ele estiver assim, \cd{high} e \cd{critical} não deixam o
CI vermelho, e nada impede que uma dependência vulnerável suba.

A consequência que dói mais é a de segunda ordem: o prazo de 2026-10-31 do
\arq{apps/api/.snyk} era cobrado \emph{pelo próprio CI} --- vencido sem o
\emph{upgrade}, o passo voltaria a falhar. Com o passo comentado, o prazo vence
\textbf{em silêncio}, e o \arq{.snyk} passa a ser exatamente o que a alternativa
descartada abaixo chama de armadilha: um registro em que todo mundo acredita e
que nada cumpre --- a mesma forma do \cd{expires\_at} sem rotina.

Isto está registrado aqui porque a ADR descrevia uma proteção que deixou de
existir, e ADR que descreve o que não está no código é pior que ADR nenhuma.
Reativar o passo é decisão de quem cuida do CI; o \arq{AGENTS.md} da raiz
proíbe que um agente o faça por conta própria, e proíbe também rodar a varredura
na sessão (\adrr{ADR-55}). O que muda com esta nota é que a ausência do portão
deixou de ser invisível.
\end{atencao}

\begin{especificacao}{Contexto}
A plataforma guarda dado de menor e roda em rede pública; dependência com
vulnerabilidade conhecida é passivo que cresce em silêncio. Na primeira
varredura (Snyk, 2026-08-29) havia \emph{criticals} reais: o Next~15 do hub e o
\cd{next@16.3.0} que o \emph{design system} arrasta para dentro de todo
consumidor, mais \emph{highs} transitivas na api (\cd{multer}, \cd{js-yaml},
\cd{lodash}) e três cujo \emph{fix} exige \emph{major} de \emph{framework}
(Nest~11) ou \emph{bump} de terceiro (Prisma).
\end{especificacao}

\textbf{Decisão.} O portão (\cd{ci.yml}) ganha
\cd{snyk test --all-projects --severity-threshold=high}: \cd{high} ou
\cd{critical} = CI vermelho = nada sobe. O estoque existente foi tratado em
três movimentos --- o hub subiu para o \textbf{Next~16} (a linha 15.x não tinha
\emph{patch} para as \emph{criticals}), o transitivo consertável entrou em
\cd{pnpm.overrides} na raiz, e o que exige Nest~11/Prisma está em
\arq{apps/api/.snyk} com \textbf{motivo e prazo} (2026-10-31): vencido o prazo
sem o \emph{upgrade}, o próprio passo volta a falhar. \cd{SNYK\_TOKEN} ausente
= passo vermelho, nunca pulado (a regra do \emph{guard} do Pub/Sub).

\textbf{Alternativas descartadas.}
\begin{itemize}
  \item \emph{Só relatório, sem travar} (\cd{continue-on-error}): vira o
    \cd{expires\_at} --- todo mundo acredita que há uma verificação, e ninguém
    lê o relatório de um passo que nunca fica vermelho.
  \item \emph{Bloquear até consertar TUDO, sem \arq{.snyk}}: com \emph{fix} que
    é \emph{major} de \emph{framework}, o pipeline ficaria vermelho por semanas
    e o \emph{deploy} pararia --- e um portão que todo mundo precisa contornar
    ensina a contornar portões.
  \item \emph{Rodar só no push para \cd{main}}: o PR mergearia vulnerável e o
    autor descobriria depois do merge, no pior momento para reverter.
\end{itemize}

\textbf{Consequências.}
\begin{itemize}
  \item Cada execução do portão testa os 12 projetos do \emph{workspace} contra
    a API do Snyk --- em plano gratuito há cota mensal de testes. Se a cota
    virar o problema, a válvula é mover a trava para o \emph{push} em
    \cd{main} + um agendado diário, nunca desligar.
  \item O \arq{.snyk} expira em 2026-10-31: o \emph{upgrade} do Nest 10\seta 11
    tem prazo de verdade, cumprido pelo próprio CI.
  \item \cd{pnpm.overrides} força versão por baixo de quem não atualizou
    (\cd{@unooria/design-system>next}, \cd{multer}, \cd{js-yaml}, \cd{lodash});
    cada entrada deve SUMIR quando o pai atualizar --- \emph{override} esquecido
    é \emph{pin} que ninguém lembra por quê.
  \item O \emph{fix} definitivo do \cd{next} transitivo é do repo
    \arq{storybook-unooria}: o DS não deveria arrastar \cd{next} para dentro de
    experiência Vite que nunca o executa.
\end{itemize}

\textbf{Gatilho de reavaliação.} A cota do Snyk barrar o portão, ou o
\arq{.snyk} vencer: os dois têm a mesma resposta --- resolver a causa (plano ou
\emph{upgrade}), não afrouxar a trava.

\section{ADR-40 --- A raiz é a \emph{landing}; a marca vem do DS, não do modelo}
\label{adr:landing}

\begin{especificacao}{Contexto}
A raiz do hub (\rota{/}) só roteava: sem sessão, \rota{/login}; com sessão, a
\emph{home} do papel. Quem chegava ao endereço da plataforma via um formulário
de usuário e senha e nada sobre o produto --- não havia \emph{landing}. A tela
foi desenhada no Stitch (projeto \cd{16689192725057015920}, ``Unooria ---
Educação Digital Infinita''), e o pedido era implementá-la fiel ao modelo.

O HTML exportado pelo Stitch traz duas coisas que não são da marca: o tema
Material que a ferramenta gera sozinha (primária \cd{\#396093}, Nunito Sans,
Inter, Material Symbols \emph{Outlined}, Tailwind por CDN) e o logo como PNG
hospedado no Google. E o \cd{login-form} e a troca de senha fazem
\cd{router.replace("/")} contando com o roteamento da raiz --- é o que leva a
pessoa à área certa sem o \emph{client} conhecer o papel.
\end{especificacao}

\textbf{Decisão.} Três partes.

\begin{enumerate}
  \item \textbf{\rota{/} sem sessão é a \emph{landing}; com sessão, continua
    roteando} para a \emph{home} do papel (ou para \rota{/trocar-senha}). Nada
    muda no \cd{login-form}, na troca de senha nem no \emph{middleware} (a raiz
    segue fora do \cd{matcher}). Quem já está dentro e quer ver a \emph{landing}
    sai da conta.
  \item \textbf{Fiel ao modelo em estrutura, copy, seções e movimento} ---
    navegação fixa em vidro, \emph{hero} sobre o \emph{shader} WebGL do modelo
    (portado do \cd{<script>} inline para um \emph{client component} com ciclo
    de vida), cards flutuantes, revelação ao rolar, rodapé --- \textbf{e do DS em
    cor, tipografia, ícones, logo e botões}: tokens \cd{--unooria-*} (a primária
    do modelo vira \cd{brandnavy-900}; os roxos, \cd{brandpurple-600/700}; o
    vermelho, \cd{danger-600}), Open Sans, \cd{Icon} (Material Symbols
    \emph{Rounded}, os mesmos nomes de ícone), \cd{Logo}/\cd{LogoSymbol}
    vetoriais no lugar dos PNG e \cd{Button} nas duas chamadas. O painel
    ``conceito'' de Nossa Resposta, que no modelo era um PNG do símbolo sob um
    gradiente, é o \cd{--unooria-color-gradient-1} com o \cd{LogoSymbol} ---
    mesma imagem, sem \emph{raster}.
  \item \textbf{O CSS da \emph{landing} vive ao lado do componente}
    (\arq{components/landing/landing.css}), não em \arq{plataforma.css}: é
    composição de página, não compensação de lacuna do DS. Toda cor ali é
    \emph{token}; translucidez sai de \cd{color-mix()} sobre o \emph{token}, não
    de literal \cd{rgba}.
\end{enumerate}

\textbf{Alternativas descartadas.}
\begin{itemize}
  \item \emph{Landing sempre visível, com uma rota nova (\rota{/inicio}) para o
    roteamento por sessão}: exigiria trocar o destino do \cd{login-form} e da
    troca de senha e criar uma segunda porta que só redireciona. Mais peças
    para o mesmo resultado --- e quem está logado raramente quer a página de
    venda.
  \item \emph{Landing fora do hub, como HTML estático no \emph{bucket}}: a
    página teria de duplicar o CSS e as fontes do DS, e o botão ``Começar
    Agora'' atravessaria origens. O hub já serve a casca do DS
    (\cd{CascaPlataforma}) --- a \emph{landing} é a quinta superfície a usá-la.
  \item \emph{Reproduzir a paleta e as fontes do Stitch}: seria a primeira
    superfície da plataforma fora da marca, e por definição a mais vista. A
    regra do DS (Capítulo~\ref{cap:ds}) existe exatamente para isso.
  \item \emph{Trazer o Tailwind por CDN e o tema do modelo}: o hub já tem
    Tailwind só para utilitárias de \emph{layout}, com o tema VAZIO de
    propósito --- cor por nome de classe é o que a decisão de
    \arq{plataforma.css} proíbe.
\end{itemize}

\textbf{Consequências.}
\begin{itemize}
  \item \textbf{O botão primário sai azul-céu, não navy} (\cd{\#0d0d35} no
    modelo): é o \cd{Button} do DS. Se a marca quiser o CTA em navy, o lugar é o
    pacote, não uma classe local.
  \item \textbf{O cleanup do \emph{shader} NÃO chama \cd{loseContext()}.} O
    Strict Mode do React (dev) monta, desmonta e remonta o efeito no MESMO
    canvas, e um canvas só tem um contexto WebGL --- perdido de propósito na
    primeira desmontagem, volta perdido na segunda montagem e o fundo fica em
    branco. Foi o primeiro sintoma do \emph{smoke}
    (\cd{contextLost: true}). Com \cd{prefers-reduced-motion} desenha um quadro
    e para; na aba oculta, para até voltar; sem WebGL, o gradiente do DS
    declarado no CSS ocupa o lugar.
  \item \textbf{A revelação ao rolar só esconde depois que o JS montou}
    (\cd{.landing--animada}). Esconder no CSS e revelar no JS deixaria a página
    em branco para quem o \emph{bundle} não alcança.
  \item \textbf{Cor de ícone vai pela \emph{prop} \cd{color} do \cd{Icon}}, não
    por classe: a folha do pacote vem depois no grafo de importação e, com peso
    igual, vence --- os ícones saíam pretos.
  \item Os \emph{links} de Privacidade, Termos e Contato do rodapé apontavam
    para \cd{\#}, como no modelo: as páginas não existiam. \textbf{Isso mudou
    em 2026-09-05} --- Privacidade e Termos apontam para \rota{/privacidade} e
    \rota{/termos} (\adrr{ADR-70}); só o Contato segue em \cd{\#}, e continua
    sendo pendência de produto.
  \item Suíte própria (\arq{scripts/e2e/publico/01-landing.mjs}): fontes do
    DS carregadas (Open Sans e a de LIGADURA dos ícones), canvas com contexto,
    cor de fundo COMPUTADA do card de vidro, todo \cd{.landing-revelar} visível
    ao fim, nada transbordando em celular, e o CTA levando ao \rota{/login}. Só
    precisa do hub no ar.
\end{itemize}

\textbf{Gatilho de reavaliação.} A plataforma ganhar uma área pública maior
(institucional, ajuda, termos): aí a \emph{landing} vira um grupo de rotas
\cd{(publico)} com \emph{layout} próprio, e os \emph{links} do rodapé deixam de
ser \cd{\#}. \textbf{Disparou pela metade em 2026-09-05}: as duas páginas legais
nasceram (\adrr{ADR-70}) e os \emph{links} delas deixaram de ser \cd{\#}, mas
como rotas próprias --- o grupo \cd{(publico)} com \emph{layout} compartilhado
ainda não existe, e nasce quando a terceira página pública aparecer.

\section{ADR-41 --- O relatório do estudante sai do aparelho: o iframe expande, não remonta}
\label{adr:relatorio-expandido}

\begin{especificacao}{Contexto}
A explicabilidade da exp-01 é a peça mais densa do produto e rodava inteira
dentro do aparelho de $375\times812$: seis blocos, um gráfico só e três dos cinco
posts de maior \emph{dwell} --- quando o \emph{payload} já carregava mais. As
barras por bloco e a linha da escalada existiam em \arq{graficos.tsx} e não
tinham consumidor de tela; a ``escalada'' virara uma frase por falta de espaço.

Ela acontece DEPOIS da revelação, e nesse ponto não há mais disfarce a proteger:
a regra da ponte (\arq{packages/ui/src/ds.ts}) já autorizava o DS da plataforma
nas ``telas pós-revelação do estudante''.
\end{especificacao}

\textbf{Decisão.} No passo do relatório, a experiência publica
\cd{exp:estado \{ fase: "explicabilidade", passo: "relatorio" \}} e o
\cd{ExpHost} liga \cd{expandido} no \cd{DispositivoFrame}, que troca APENAS o
CSS do próprio \cd{<div>} interno. O relatório vira claro, responsivo, com os
três gráficos e o \emph{top} 5. A revelação continua no aparelho, escura: ela é
a fronteira, e não estava comprimida.

\textbf{Alternativas descartadas.}
\begin{itemize}
  \item \emph{Um \cd{Modal size="full"} do hub buscando \rota{/runs/:id/dossie}}:
    o BFF já permitiria (a rota autoriza o estudante a ler o próprio) e havia
    precedente em \arq{modal-estudante.tsx}. Descartado porque duplicaria a
    renderização do relatório em dois apps --- ou obrigaria a experiência a
    perder os \emph{exports} em texto e PNG, que são dela.
  \item \emph{Uma rota própria do hub} (\rota{/estudante/relatorio/:runId}): tira o
    estudante do aparelho no meio de uma aula conduzida por botão.
  \item \emph{\cd{requestFullscreen()} de dentro do \emph{iframe}}: o
    \cd{allow="\ldots fullscreen"} já está delegado a todas as experiências e
    ninguém usa. O hub perderia o controle da moldura, e a largura lógica
    continuaria a mesma --- ficaria maior, não menos comprimido.
\end{itemize}

\textbf{Consequências.}
\begin{itemize}
  \item \textbf{Expandir é ESTILO, nunca árvore.} Tirar o \cd{children} do
    \cd{DispositivoFrame} para renderizá-lo fora desmontaria o \cd{<iframe>} e
    o estudante voltaria à revelação com a sessão recarregada. É por isso que o
    modo é uma \emph{prop}, e não um \cd{if} trocando o \emph{wrapper}.
  \item \textbf{A condição é POSITIVA, e o \emph{default} é contraído.}
    Qualquer outro estado --- inclusive mensagem perdida --- cai no aparelho de
    375px. O \cd{IntersectionObserver} do \emph{dwell} se refere ao
    \emph{iframe}: um FEED numa área larga inflaria o sinal central do motor.
  \item \textbf{O layout largo é \cd{md:}, o estreito é o \emph{default}.} Se a
    expansão não acontecer, o estudante vê exatamente o que via.
  \item \textbf{Nenhuma mensagem hub$\rightarrow$exp para o layout}: o
    \emph{iframe} muda de largura de verdade, e as \emph{media queries} de
    dentro respondem sozinhas. Os dois gráficos a mais são renderizados e
    escondidos por CSS --- perguntar a largura em JS custaria um
    \cd{matchMedia} e um render a mais na hidratação.
  \item A volta usa o caminho que já existia: o professor avança, o
    \emph{polling} de 5s vê a fase nova, a experiência publica
    \cd{\{ fase: "perdidos" \}} sem \cd{passo}, e o aparelho volta. Até 5s de
    atraso, que é a latência que a troca de tela já tinha.
  \item \textbf{Paridade}: o relatório do professor (\arq{modal-estudante.tsx})
    ganhou a seção \emph{``por que o feed dele começou desse lado''}, que ele
    não lia --- e é a frase que a criança acabou de ler.
\end{itemize}

\textbf{Gatilho de reavaliação.} Uma segunda experiência precisar do mesmo
gesto: aí \cd{expandido} deixa de ser lido de dois campos conhecidos no
\cd{ExpHost} e vira mensagem própria no \cd{ExpParaHubSchema}, com \emph{bump}
de \cd{protocolVersion}.


\section{ADR-42 --- ``O que o algoritmo escondeu'' é a EXPOSIÇÃO, não o domínio}
\label{adr:exposicao}

\begin{especificacao}{Contexto}
Havia duas contas para a mesma frase. A tela do estudante (\rota{/runs/:id/perdidos})
decidia qual era ``o lado dele'' CONTANDO os posts servidos
(\cd{consumo.A >= consumo.B}); o relatório do professor
(\arq{dossie/computar.ts}) decidia pelo \cd{dominioDe()}, o \emph{score} do
motor --- \emph{dwell}, sinais, escada de intensidade.

Enquanto as duas apontam o mesmo lado, os números batem. Elas divergem
exatamente quando o estudante VIRA O JOGO, porque virar o jogo \emph{é} a definição
de a atenção ter terminado do lado oposto ao que o feed empurrou. Ou seja: as
duas telas se contradiziam justamente sobre os estudantes que o professor mais quer
discutir.

O caso está registrado numa auditoria de UX: sobre o MESMO estudante, o painel dizia
``terminou dominante em Processo Lento e Gradual, virou o jogo'' e a tela dele
dizia ``32 posts que nunca chegaram até você, todos de Processo Lento e
Gradual''. Duas frases sobre a mesma criança, nomeando lados opostos. O
relatório do professor (``escondeu 30'') contava os posts do lado que o
algoritmo estava EMPURRANDO nele.
\end{especificacao}

\textbf{Decisão.} Uma conta só, em \arq{apps/api/src/dossie/exposicao.ts},
consumida pelas duas rotas: ``o lado dele'' é o que o feed mais SERVIU, e o
número é quantos aprovados do outro lado nunca foram servidos.
\cd{ladoDominante} continua existindo e continua no relatório --- ele responde
``para onde a atenção dele foi'', que é outra pergunta, e é a que sustenta o
``virou o jogo''.

\textbf{Alternativa descartada.} \emph{Padronizar no domínio.} É a leitura mais
fiel ao que a atenção do estudante fez, mas ``o algoritmo escondeu'' é uma
afirmação sobre o que o feed ENTREGOU: medir pelo domínio credita à máquina o
resultado do comportamento da criança. E com \cd{dominio $\le$ 0{,}5} o
resultado era \cd{indefinido}, que virava ``escondeu 0 posts''.

\textbf{Consequências.}
\begin{itemize}
  \item \textbf{O número do professor MUDA} em atividades já rodadas ---
    \cd{VERSAO\_DO\_RELATORIO} sobe para 4 e os dossiês congelados são
    recomputados na leitura. É o preço de ter contado errado.
  \item \textbf{Sem post de tese servido, não se nomeia lado nenhum.} A rota
    devolve \cd{ladoPerdido: null} e a tela diz que a sessão terminou antes de o
    feed abrir. ``0 posts que nunca chegaram até você'' afirmaria o contrário do
    que houve --- é a mesma distinção entre \cd{zero} e \cd{none} da Vila do
    Lume.
  \item \textbf{Post desaprovado não conta como escondido}: nunca teve chance de
    aparecer para ninguém, e culpar o algoritmo por uma decisão editorial do
    professor seria outra mentira pequena.
  \item A trava é \arq{apps/api/tests/exposicao-lado-oposto.spec.ts}, e o teste
    do meio monta exatamente o estudante que virou o jogo.
\end{itemize}

\textbf{Gatilho de reavaliação.} Alguém precisar de ``o oposto do lado
dominante'' para outra tela: é uma terceira grandeza, e ela nasce com nome
próprio --- nunca reaproveitando esta.


\section{ADR-43 --- A resposta do estudante guarda a primeira e a última}
\label{adr:conclusao-primeira}

\begin{especificacao}{Contexto}
A tela da resposta não tem botão de confirmar: tocar na pílula já registra, e o
quadro diz ``pode trocar de ideia enquanto espera --- vale a última''. O estudante
oscila de propósito enquanto os colegas terminam.

O \cd{upsert} sobrescrevia \cd{inclinacao} e a primeira escolha desaparecia. O
\cd{created\_at} ficava congelado no primeiro envio enquanto a inclinação
mudava, ou seja, passava a mentir sobre a própria linha. E a correlação ``lado
em que o feed nasceu $\times$ resposta dada'' --- o dado mais forte que esta
atividade produz --- não tinha como dizer se estava lendo a reação imediata ao
feed ou o que sobrou depois de dois minutos de espera e de olhar a tela do
colega.
\end{especificacao}

\textbf{Decisão.} Migração aditiva em \cd{exp\_01\_conclusoes}:
\cd{inclinacao\_inicial} (gravada uma vez, nunca mais tocada), \cd{trocas} e
\cd{atualizado\_em}. A correlação continua usando a ÚLTIMA --- é a que o estudante
assumiu e é a que ele vai defender no debate --- e a análise da turma passa a
DIZER isso, com quantos mexeram na resposta e quantos terminaram noutro lado.

\textbf{Alternativas descartadas.}
\begin{itemize}
  \item \emph{Emitir evento de telemetria a cada toque}: o canal analítico tem
    retenção indefinida e \cd{estudante\_id}; guardar a oscilação de opinião de uma
    criança num JSONB de eventos é mais dado do que a pergunta pede.
  \item \emph{Trocar a correlação para a primeira escolha}: seria decidir pelo
    professor qual das duas conta. A decisão é dele, e agora ele tem os dois
    números.
\end{itemize}

\textbf{Consequências.} \cd{inclinacao\_inicial} é NULL nas linhas anteriores à
migração, e é assim que se lê ``esta resposta é anterior à mudança'':
preenchê-la com a inclinação atual inventaria uma primeira escolha que ninguém
observou. Reenviar a mesma opção (toque duplo, rede lenta) não conta como troca.
As duas colunas novas vão ao CSV da análise.

\textbf{Gatilho de reavaliação.} Aparecer demanda por ver a oscilação completa
(não só a primeira e a última): aí é tabela de histórico, não mais colunas.


\section{ADR-44 --- A janela do envolvimento termina com o feed, e o total sai das listagens}
\label{adr:envolvimento}

\begin{especificacao}{Contexto}
Duas coisas, medidas na mesma auditoria.

A primeira: o painel ao vivo mostrava a imersão da turma DERRETENDO conforme a
aula avançava ($52{,}7 \rightarrow 42{,}8 \rightarrow 38{,}7 \rightarrow 35{,}1$)
e isso lia como adolescente desengajando. Não era. O \emph{buffer} de eventos da
experiência vive enquanto a \emph{run} vive: durante o debate, a explicação e a
reflexão --- mais de metade do roteiro de 100 minutos --- cada \cd{hub:pause}
ainda registra \cd{perdeu\_foco}. Com a conta fechando no último evento da
SESSÃO, \cd{feedMs} crescia depois de o feed fechar, cada vão virava
\cd{ociosoMs}, e \cd{ativoMs / feedMs} caía sozinha.

A segunda: o total 0--100 aparecia como número numa COLUNA de tabela, por
criança, no monitor e na análise --- ordenável, comparável. O próprio código
dizia, em comentário, que ``o total existe para ordenar uma lista, não para ser
exibido sozinho'', e a tela o exibia sozinho. E ele é somável justamente onde
menos se sustenta: \emph{Gestos} e \emph{Exploração} ficam em 0/20 na maioria
das sessões, então o teto real não é 100.
\end{especificacao}

\textbf{Decisão.} \cd{fimDoFeedDe(eventos)} --- o último \cd{visto}, que é o fim
do feed por definição --- corta a janela em \cd{calcularTempo} E no desconto por
perda de foco de \cd{calcularImersao}. E o total sai das duas listagens por
estudante, ficando as cinco barras (com o eixo no \cd{title}); o número segue no
popup do estudante, que é uma criança por vez, e nos agregados de turma. Um nome só:
\emph{Envolvimento}, que estava ``Imersão'' em quatro telas e ``Envolvimento''
numa.

\textbf{Alternativas descartadas.}
\begin{itemize}
  \item \emph{Cortar por \cd{run.finishedAt}}: \cd{clientTs} é o relógio do
    APARELHO e \cd{finishedAt} é o do servidor. Comparar os dois numa escola com
    relógio desregulado traria o mesmo erro de volta por outro caminho.
  \item \emph{Parar de registrar eventos fora do feed, no cliente}: corrigiria
    daqui para a frente e não tocaria em nenhum dado já coletado --- e depende de
    a experiência estar atualizada.
  \item \emph{Recalibrar os tetos de Gestos e Exploração}: exige medição em
    turmas reais. Fica em aberto; mexer nos tetos por intuição trocaria um número
    errado por outro.
\end{itemize}

\textbf{Consequências.} \cd{VERSAO\_DO\_RELATORIO} sobe para 4 e recomputa o que
está congelado. Um teste do \arq{dossie-relatorio.spec.ts} precisou de um
\cd{visto} para delimitar o feed: ele montava trinta \cd{perdeu\_foco} e nenhum
post visto, ou seja, descrevia perdas de foco fora de qualquer feed --- uma
distinção que antes não existia. O professor perde a leitura de um número por
estudante na tabela; ganha o padrão dos cinco eixos, que é o que a conversa com ele
usa.

\textbf{Gatilho de reavaliação.} Dados de turmas reais mostrarem que
\emph{Gestos} e \emph{Exploração} pontuam: aí os tetos são recalibrados e a
pergunta de exibir o total volta a fazer sentido.

\section{ADR-45 --- \emph{Agent} e \cd{fetch} vêm do MESMO pacote undici}
\label{adr:undici-fetch}

\begin{especificacao}{Contexto}
O provider do Vertex usa um \emph{dispatcher} próprio: \cd{generateContent}
não-\emph{streaming} só devolve os \emph{headers} quando a geração termina, e o
teto de 300\,s do \cd{fetch} do Node estourava no meio de um plano anual. A
solução era um \cd{Agent} do pacote \cd{undici} com 15 minutos, passado como
opção \cd{dispatcher} ao \cd{fetch} GLOBAL.

Isso acopla duas coisas que ninguém obriga a concordar. O \cd{fetch} global usa
o undici EMBUTIDO na \emph{runtime} (\cd{process.versions.undici}), que é outro
pacote, com outra versão. Com os \arq{Dockerfile} fixados em \textbf{Node 26.7}
(undici \textbf{8.x}) e a API declarando \cd{"undici": "\^{}6"}, o \cd{Agent} da
6 deixou de satisfazer a interface de \emph{handler} da 8.

O sintoma não aponta para nada: \cd{fetch failed}. A razão real
(\cd{invalid onError method}) vive só em \cd{error.cause}, e o \cd{catch} da
geração logava apenas \cd{error.message}. \textbf{TODA} chamada de IA morria ---
plano de aula, \emph{unooGPT}, arte --- e por dias isso pareceu falha de
credencial, porque havia de fato um ADC vencido por cima.
\end{especificacao}

\textbf{Decisão.} Importar o \cd{fetch} do PRÓPRIO undici
(\cd{import \{ Agent, fetch \} from "undici"}). Vindos do mesmo pacote, os dois
concordam por construção, e a versão embutida na \emph{runtime} deixa de ser
uma dependência invisível. O \emph{log} da geração passa a imprimir o
\cd{cause} junto --- só no \emph{log}: a mensagem que vai ao banco é a que o
professor lê.

\textbf{Alternativas descartadas.}
\begin{itemize}
  \item \emph{Subir \cd{undici} para \cd{\^{}8}}: funciona hoje e quebra no
    próximo \emph{bump} de Node, porque a condição continua sendo ``duas versões
    coincidirem'' --- e nada no CI verifica isso. Seria a mesma armadilha,
    adiada.
  \item \emph{Largar o \emph{dispatcher} e usar o \cd{fetch} global puro}:
    devolveria o teto de 300\,s que motivou o \emph{dispatcher}, e o plano anual
    voltaria a morrer no meio.
  \item \emph{Fixar o \cd{undici} pela versão da \emph{runtime}}: exigiria ler
    \cd{process.versions.undici} e casar no \arq{package.json}. Duas fontes de
    verdade sobre a mesma versão, uma delas fora do controle do repo.
\end{itemize}

\textbf{Consequências.} O teste \arq{ai-conversar.spec.ts} dublava
\cd{global.fetch} e passou a não interceptar nada --- voltava \emph{array}
vazio. Ele agora mocka o módulo \cd{undici} encaminhando para \cd{global.fetch},
o que mantém a dublagem em uma linha e o que se testa segue sendo a lógica do
provider. Fica registrado que o \cd{dispatcher} é a razão de o \cd{fetch} não
poder ser o global: quem apagar a importação ``porque o Node já tem \cd{fetch}''
reabre exatamente este defeito.

\textbf{Gatilho de reavaliação.} O Vertex passar a ser consumido por
\cd{streamGenerateContent?alt=sse}: com \emph{streaming} os \emph{headers} chegam
de imediato, o teto de 300\,s deixa de incomodar, e o \emph{dispatcher} ---
junto com todo este acoplamento --- pode simplesmente sair.

% ---------------------------------------------------------------------------
\section{ADR-46 --- O \cd{?hubOrigin=} é entrada não confiável: a experiência confirma quem a embedou}
\label{adr:hub-origin}

\begin{especificacao}{Contexto}
A regra 4 diz que o \emph{token} passa do hub para a experiência por
\cd{postMessage} com \cd{targetOrigin} explícito, nunca \cd{"*"}. A bridge
cumpria a letra da regra e não o espírito: o valor do \cd{targetOrigin} ---
e, o que é pior, o valor comparado em \cd{ev.origin !== hubOrigin} para decidir
de quem se aceita o \cd{hub:init} --- vinha CRU da \emph{query string}
(\cd{new URLSearchParams(location.search).get("hubOrigin")}).

O \emph{bundle} da experiência é público por construção (ADR-01: ele mora num
\emph{bucket} servido por CDN). Quem abrisse
\arq{https://cdn/exp-01/1.0.0/index.html?hubOrigin=https://evil.tld} dentro de
um \emph{iframe} próprio escrevia, com aquele parâmetro, a política de confiança
da experiência: passava a poder entregar um \cd{hub:init} forjado --- com
\cd{estudante}, \cd{apiUrl} e \cd{token} escolhidos por ele --- e a receber
\textbf{tudo} o que a experiência emitisse depois (\cd{exp:concluir},
\cd{exp:estado}). O Snyk classificou como \emph{code injection} no ponto em que
o parâmetro entra; a leitura é discutível, o caminho não.
\end{especificacao}

\textbf{Decisão.} Todo \cd{hubOrigin} passa por \cd{normalizarHubOrigin()}
(\arq{packages/sdk/src/bridge.ts}), chamado tanto onde o parâmetro é lido quanto
dentro do \cd{connectToHub} --- a função é idempotente, e a segunda chamada
existe porque é ali que o valor vira \cd{targetOrigin}. São três travas:

\begin{enumerate}
  \item \cd{new URL(...).origin}: o valor vira uma ORIGEM canônica, nunca uma
    \emph{string} com caminho, credencial ou \emph{query}. Não é higiene:
    \cd{postMessage} compara só a origem, mas a igualdade do \emph{handler}
    nunca casaria, e o sintoma seria \emph{timeout} de \emph{handshake} em vez
    de ``o parâmetro está errado''. Esquema fora de http(s) e origem opaca
    (\cd{"null"}) são recusados.
  \item Experiência servida por https exige hub em https --- a mesma regra do
    JWKS da portabilidade. Em desenvolvimento tudo é http e a trava não pesa.
  \item \cd{location.ancestorOrigins[0]}, quando o navegador publica: é a origem
    REAL de quem embedou, e ela VENCE a \emph{query string}. É a trava que
    fecha o caminho de fato --- o atacante controla o parâmetro, não o endereço
    do documento que abriu o \emph{iframe}.
\end{enumerate}

\textbf{Alternativas descartadas.}
\begin{itemize}
  \item \emph{Lista de origens permitidas embutida no \emph{bundle}}: a
    experiência é a mesma em todos os contratos, e cada contrato tem o seu
    domínio de hub. A lista viraria \emph{env} de \emph{build}, ou seja, um
    \emph{bundle} por contrato --- exatamente o que a ADR-01 evita.
  \item \emph{Deduzir a origem de \cd{document.referrer}}: depende da
    \emph{referrer policy} do hub e some com \cd{no-referrer}. O
    \cd{ancestorOrigins} não é forjável pela página.
  \item \emph{Deixar como estava, porque o \emph{token} forjado é do próprio
    atacante}: verdadeiro para o roubo direto e irrelevante para o resto ---
    o que o atacante ganha é uma experiência que fala com a API dele e devolve
    a ele o que a criança produziu.
\end{itemize}

\textbf{Consequências.} O Firefox não implementa \cd{ancestorOrigins}; lá
sobram as duas primeiras travas. \cd{hubOrigin} malformado deixa de morrer por
\emph{timeout} de \emph{handshake} e passa a lançar com a razão escrita ---
diagnóstico melhor, mas é uma falha NOVA para quem embedar a experiência fora
do hub. \cd{HUB\_ORIGIN} mal configurado no hub, que antes só travava o
\emph{handshake}, agora recusa explicitamente.

\begin{nota}
\textbf{O apontamento do Snyk que originou esta ADR NÃO some com ela, e está
certo assim.} A regra \emph{Code Injection} (CWE-94) desenha o fluxo
\cd{location.search} \seta{} \cd{setInterval} e afirma que o valor ``é
executado como JavaScript''. O \emph{sink} apontado é o
\cd{setInterval(anunciar, READY\_RETRY\_MS)} do \emph{retry} do
\cd{exp:ready}, e \cd{setInterval} só AVALIA código quando o primeiro
argumento é uma \emph{string}: ali ele recebe uma referência de FUNÇÃO, e o
\cd{hubOrigin} entra por captura de \emph{closure} para terminar como
\cd{targetOrigin} de um \cd{postMessage}, que não executa nada. O valor tem
três usos no pacote --- dois \cd{postMessage} e uma comparação com
\cd{ev.origin} ---, e nenhum é \emph{sink} de execução.

O que aquele fluxo tinha de verdadeiro era a origem de confiança vindo da URL, e
é isso que esta ADR fecha. Fica registrado para que ninguém CONTORÇA o código
para agradar o \emph{scanner}: esconder a captura do \emph{closure} deixaria o
\emph{retry} ilegível sem tirar risco nenhum. A nota vive junto da chamada, em
\arq{packages/sdk/src/index.ts}.
\end{nota}

\textbf{Gatilho de reavaliação.} O dia em que uma experiência precisar ser
embedada por um hub de terceiros (LTI, por exemplo): a trava 3 passa a exigir
uma lista de embedadores conhecidos, e ela terá de vir do manifesto --- não do
parâmetro.

% ---------------------------------------------------------------------------
\section{ADR-47 --- Nem senha de demonstração nem verificação de TLS moram no fonte}
\label{adr:segredo-fora-do-fonte}

\begin{especificacao}{Contexto}
Três apontamentos do Snyk com a mesma raiz: credencial e política de segurança
escritas como constante no código, onde não há como diferenciar
desenvolvimento de produção.

\begin{itemize}
  \item \arq{apps/api/prisma/seed.ts} trazia \cd{SENHA\_PADRAO = "unooria123"}.
    O seed é idempotente e roda contra o \cd{DATABASE\_URL} que estiver no
    ambiente: apontado por engano para o banco de um contrato, ele criaria 61
    contas --- inclusive o \cd{admin-contrato}, o único papel que atravessa
    escolas --- com uma senha publicada no \arq{README}.
  \item \arq{wso2/bootstrap/bootstrap.mjs} abria com
    \cd{process.env.NODE\_TLS\_REJECT\_UNAUTHORIZED = "0"}, porque o IS local usa
    o certificado \emph{self-signed} da distribuição (CN=\cd{localhost}, host
    \cd{identity-server}). Aquilo desliga a verificação do PROCESSO INTEIRO,
    para toda origem: o \emph{script} que carrega a senha do \emph{admin} do
    \emph{tenant} passava a aceitar qualquer certificado de qualquer servidor.
  \item \arq{apps/api/src/auth/auth.controller.ts} calculava o \emph{hash}
    ``dummy'' anti-\emph{timing} sobre uma \emph{string} literal.
\end{itemize}
\end{especificacao}

\textbf{Decisão.} O princípio é um só: o que é credencial vem do AMBIENTE, e o
que é confiança se PROVA em vez de ser dispensada.

\begin{itemize}
  \item \cd{SEED\_SENHA} é obrigatória e não tem \emph{default} no fonte ---
    ausente, o seed falha alto. O \emph{default} de desenvolvimento vive nos
    dois \arq{docker-compose}, exatamente como o \cd{docs}/\cd{unooria-docs-local}
    do Swagger (ADR-36). O seed também deixou de imprimir a senha.
  \item O \emph{bootstrap} do WSO2 PINA o certificado: lê o certificado do IS no
    primeiro contato (o único \cd{rejectUnauthorized: false} que resta, e ele
    não autentica nada --- serve para DESCOBRIR a chave), grava o PEM que api e
    hub já consumiam por \cd{NODE\_EXTRA\_CA\_CERTS} e o usa como a ÚNICA CA
    aceita nas chamadas seguintes. Só a conferência de NOME do host é
    dispensada, com \cd{checkServerIdentity}, porque o CN da distribuição não
    bate com o nome do serviço --- a identidade da chave continua obrigatória.
  \item O \emph{hash} dummy passa a ser calculado sobre
    \cd{randomBytes(32)}, sorteado a cada subida e descartado no mesmo instante.
\end{itemize}

\textbf{Alternativas descartadas.}
\begin{itemize}
  \item \emph{Manter o literal do seed e ignorar o apontamento no
    \arq{.snyk}}: o \arq{.snyk} deste repo é para vulnerabilidade cujo
    \emph{fix} exige \emph{major} de \emph{framework}, com prazo cobrado pelo
    CI. Aqui o \emph{fix} custa dez linhas.
  \item \emph{Sortear a senha do seed quando a \emph{env} falta}: nada ficaria
    no fonte, mas as suítes de ponta a ponta passariam a depender de ler
    \cd{stdout}, e o \arq{README} não teria o que documentar.
  \item \emph{Passar a CA do IS por \cd{NODE\_EXTRA\_CA\_CERTS} no
    \emph{bootstrap}}: a variável é lida na SUBIDA do processo, e o certificado
    só existe depois que o \emph{script} o busca. Exigiria dois processos.
  \item \emph{Aceitar também o nome do host, reescrevendo o certificado do IS}:
    trocar a distribuição por um certificado gerado no \emph{up} é mais infra
    local para operar, num serviço que já leva dois minutos para subir.
\end{itemize}

\textbf{Consequências.} \cd{pnpm --filter @edu/api seed} no \emph{host} exige
\cd{SEED\_SENHA} --- no \arq{apps/api/.env} (que o Prisma carrega) ou na linha
de comando. Nos dois \emph{compose} nada muda. O \emph{bootstrap} do WSO2 trocou
o \cd{fetch} global pelo \cd{node:https}, porque o \cd{fetch} do Node não aceita
agente e é por agente que se passa uma CA; são vinte linhas a mais, e é o preço
de não desligar o TLS do processo. O \emph{bootstrap} passou a FALHAR se não
conseguir ler o certificado, onde antes seguia com um aviso: sem CA não há como
falar com o IS sem voltar a aceitar qualquer certificado.

\textbf{Gatilho de reavaliação.} O WSO2 deixar de ser \emph{self-signed} em
desenvolvimento (certificado de uma CA local no \emph{up}): a pinagem e o
\cd{checkServerIdentity} saem juntos, e o \cd{chamar} pode voltar ao
\cd{fetch}.

\section{ADR-48 --- O teto do CSV é de BYTES e cabe debaixo do limite do
ambiente; o modelo do OneRoster vem em pares, e o exemplo leva a escola de
verdade}
\label{adr:modelos-oneroster}

\begin{especificacao}{Contexto}
A tela de importação tem duas abas. A da \emph{planilha da escola} oferece dois
arquivos para \emph{download} --- o modelo em branco e um exemplo preenchido ---
e os dois são GERADOS da mesma constante que o \emph{parser} consome. A do
\textbf{OneRoster 1.2} lista os sete arquivos do \emph{binding} CSV e não
oferecia nenhum: quem não tem o \emph{export} pronto do sistema de gestão --- e
é a maioria antes da primeira integração --- tinha de montar o cabeçalho lendo a
especificação da 1EdTech.

O teto do corpo era \SI{10}{\mega\byte}, dimensionado para a planilha de UMA
escola, e o \cd{MAX\_LINHAS} de 50 mil, que existe do lado da planilha, não vale
para o caminho OneRoster. Medido com o formato real do padrão
(\textasciitilde\SI{130}{\byte} por linha de \arq{enrollments.csv}), o pior
caso realista de UM contrato --- \textasciitilde10 mil estudantes, cinco
componentes cada --- dá \textasciitilde\SI{7}{\mega\byte}. Os
\SI{10}{\mega\byte} apertavam exatamente ali.

Acima da aplicação há um teto que não é nosso: o Cloud Run recusa requisição
maior que \SI{32}{\mebi\byte} em HTTP/1, no \emph{frontend} do Google e
ANTES do contêiner. Não é quota que se peça para aumentar --- é como a camada de
\emph{serving} funciona, porque ela bufferiza o corpo inteiro antes de
repassar.
\end{especificacao}

\textbf{Decisão.}

\begin{itemize}
  \item O teto de BYTES sobe para \SI{30}{\mega\byte} e passa a viver numa
    constante só, \cd{TETO\_DO\_ARQUIVO} em \arq{rostering.service.ts}, que o
    \emph{bodyParser} de \arq{main.ts} também lê. O número é escolhido para
    ficar DEBAIXO dos \SI{32}{\mebi\byte} do ambiente: a aplicação não promete
    o que o Cloud Run não entrega, e a mensagem que o gestor vê é sempre a nossa
    --- a que diz para dividir o arquivo. Ainda são
    \textasciitilde4$\times$ o pior caso medido.
  \item A ORDEM entre os dois tetos importa: com o \emph{parser} menor que a
    rota, o corpo morre no \emph{middleware} e o gestor recebe um
    \cd{PayloadTooLarge} do Express, sem o texto que diz o que fazer. O
    \cd{mb} do \emph{parser} é MEBIbyte (\cd{30mb} = \SI{31.4}{\mega\byte}),
    então ele fica \textasciitilde5\% ACIMA da rota --- e é essa folga que faz
    a nossa mensagem ser a que aparece.
  \item O \cd{MAX\_LINHAS} de 50 mil NÃO sobe: ele é da planilha de uma escola,
    e 50 mil pessoas numa unidade é um arquivo errado, não um arquivo grande.
  \item Cada arquivo do padrão ganha o mesmo par da planilha:
    \rota{GET /rostering/modelos/oneroster/:arquivo} (só o cabeçalho) e
    \rota{.../:arquivo/exemplo} (preenchido). As colunas do modelo são um
    SUBCONJUNTO declarado de \cd{COLUNAS\_POR\_ARQUIVO}, a constante que o leitor
    confere, e \arq{tests/oneroster-modelos.spec.ts} cobra as DUAS direções ---
    obrigatória que falte no modelo, e coluna do modelo que o \emph{parser} não
    conheça (o dado entraria e sumiria).
  \item O exemplo é gerado com a escola de VERDADE no \cd{sourcedId} do
    \arq{orgs.csv} e no \cd{schoolSourcedId} das turmas e matrículas, e com o ano
    corrente no calendário.
  \item Os sete exemplos formam um lote COERENTE --- uma escola, um ano letivo,
    duas turmas, quatro pessoas e o enturmamento delas. O teste aplica os sete
    contra o banco, na ordem da tela, e afirma turma, papéis e matrículas.
\end{itemize}

\textbf{Alternativas descartadas.}
\begin{itemize}
  \item \emph{Escrever os sete arquivos à mão, como texto fixo}: é a armadilha
    que o modelo da planilha já evitou uma vez. Modelo escrito à mão diverge do
    \emph{parser} no primeiro campo novo, e o sintoma chega como ``baixei o
    modelo de vocês e deu erro de cabeçalho''.
  \item \emph{Publicar o modelo com TODAS as colunas que a api sabe ler}:
    \arq{users.csv} tem 24. Modelo com 24 colunas não é modelo, é a
    especificação --- e a escola preencheria \cd{sms}, \cd{phone} e
    \cd{agentSourcedIds} sem que nada os leia.
  \item \emph{Um \cd{esc.exemplo} fixo no \arq{orgs.csv}}: para o gestor, o lote
    inteiro voltaria como ``escola não é a desta importação'', e o arquivo que
    existe para conferir o caminho provaria o contrário. Para o
    \cd{admin-contrato}, que PODE criar escola por CSV (ADR-34), o exemplo
    abriria um \emph{tenant} de mentira --- com RLS e dado de menor atrás dele.
  \item \emph{Um único arquivo por item, sem o modelo em branco}: é o exemplo que
    o gestor abre para entender, e o modelo que ele preenche. Linha de exemplo
    dentro do modelo vira pessoa fantasma importada no primeiro uso --- a mesma
    razão pela qual \arq{planilha.ts} mantém os dois separados.
  \item \emph{Empacotar os sete num \emph{zip}}: empacotamento é embalagem, não
    semântica. A rota de importação já é um arquivo por chamada, e um
    \emph{zip} para baixar exigiria dependência nova para produzir algo que o
    gestor descompactaria antes de usar.
  \item \emph{\SI{100}{\mega\byte}} --- foi o primeiro número escolhido, e ele
    é PIOR que \SI{30}{\mega\byte}: entre \SI{33.5}{\mega\byte} e
    \SI{100}{\mega\byte} a plataforma prometeria algo que o ambiente recusa, e
    o gestor receberia um 413 que não é nosso, sem mensagem útil e sem linha
    nenhuma no log da api. Teto que promete mais do que o ambiente entrega não é
    folga, é o erro que ninguém diagnostica.
  \item \emph{HTTP/2 ponta a ponta (h2c) para tirar o limite do Cloud Run}: é a
    saída documentada, porque em HTTP/2 a requisição é \emph{streamada} em vez
    de bufferizada. Custa trocar o servidor HTTP do Nest e fazer o hub repassar
    em HTTP/2 --- um modo de transporte novo para operar, por um problema que os
    números acima dizem que não existe.
  \item \emph{Streaming do corpo em vez de teto maior}: resolveria a memória, mas
    a leitura inteira do CSV é o que permite recusar o lote com o número da linha
    e o nome do campo. Trocar isso por um \emph{parser} incremental é reescrever
    as duas fronteiras para um caso que ainda não apareceu.
\end{itemize}

\textbf{Consequências.} O corpo do CSV é lido inteiro para a memória em três
lugares --- o \emph{route handler} do BFF do hub, o \emph{bodyParser} da api e o
\cd{parseCsv} ---, então \SI{30}{\mega\byte} de arquivo custam alguns múltiplos
disso em pico de \emph{heap} por importação simultânea. O \emph{rate limit} de
5 por minuto por usuário é o que segura isso hoje, e ele ficou
\textasciitilde3$\times$ mais barato do que ficaria com
\SI{100}{\mega\byte}.

A rede que precisar de mais que \SI{30}{\mega\byte} num arquivo só NÃO tem
saída pela tela: ela divide o \arq{enrollments.csv} por escola ou por diretoria,
que é uma operação que o formato já suporta (um arquivo por chamada). Os
sintomas dos três tetos --- o nosso, o do \emph{parser} e o do ambiente ---
estão no Apêndice~\ref{ap:operacao}.

\textbf{Gatilho de reavaliação.} Uma rede chegar com \emph{export} que não caiba
em \SI{30}{\mega\byte} depois de dividido: aí a conversa é sobre o CAMINHO ---
\emph{upload} direto para o \emph{bucket} com URL assinada e importação
assíncrona a partir dele ---, não sobre subir o número. Aquele caminho tira o
limite do ambiente e o pico de \emph{heap} de uma vez.

\section{ADR-49 --- A atividade do Instanooria declara a TURMA, e o painel dela
para de ler a escola inteira}
\label{adr:atividade-por-turma}

\begin{atencao}
\textbf{REVISTA em 2026-09-05 pela \adrr{ADR-69}: a turma virou
OBRIGATÓRIA.} A alternativa descartada aqui (\emph{uma atividade por turma,
obrigatória}) foi adotada --- \cd{POST /atividades} exige \cd{turmaId} e o
\enquote{Todas as turmas da escola} saiu do formulário. O que faltava em
2026-08-30 apareceu: a plataforma ganhou o modelo de \textbf{sala de aula}, em
que a turma é o contexto de tudo o que o professor faz, e a atividade sem turma
era a última peça que ainda falava de um lugar que não é uma sala --- era ela
que fazia 2EM-A e 2EM-B aparecerem como uma lista só no painel de perfis, no
monitor e na análise.

\textbf{O custo previsto aqui é real e foi aceito}: quem dá a mesma aula em
quatro salas cria quatro atividades. A tabela de junção continua sendo a
resposta certa se isso doer.

O resto desta ADR segue valendo, e a COLUNA continua anulável e sem
\emph{backfill}: as atividades criadas antes desta data seguem escolares, com a
faixa que já as explica. Fechou a porta de criar mais uma assim, não o passado.
\end{atencao}

\begin{especificacao}{Contexto}
A \cd{exp\_01\_activities} tinha \cd{escola\_id} e \cd{professor\_id}, e nenhum
vínculo com \cd{turmas} --- embora o próprio esquema já falasse como se tivesse:
o comentário de \cd{etapa} chama o campo de ``relógio COLETIVO da turma'', e a
função que monta o painel de perfis se chama \cd{distribuicaoDaTurma}.

Ao investigar por que a importação dos perfis da Vila do Lume
(\cd{POST /atividades/:id/perfis/da-vila}) varria a escola, apareceu uma coisa
pior: aquela função rodava \cd{user.findMany(\{ where: \{ role: 'estudante' \} \})}
--- \textbf{sem \cd{escola\_id}}. E \cd{users} está FORA da RLS: a lista de
tabelas com política em \arq{20260824145301\_init} não a inclui, e o
\cd{comTenancy} apenas arma \cd{app.escola\_id} para o Postgres, sem tocar no
\cd{where} do Prisma. O painel do professor listava nome e login de crianças de
TODAS as escolas do contrato, e \cd{contagem}, \cd{desequilibrio} e
\cd{precisaRedistribuir} --- que decidem quantos estudantes o botão
``Redistribuir a turma'' move --- eram calculados sobre esse universo.

O \cd{DaVilaSchema} já aceitava \cd{turmaId} opcional desde que a rota nasceu, e
a tela nunca o enviou: mandava \cd{\{\}}. O relatório dizia ``N ainda não
jogaram'' sobre crianças de salas que o professor não tem.
\end{especificacao}

\textbf{Decisão.} Duas coisas, e a primeira vale sozinha.

\begin{enumerate}
  \item \textbf{\cd{escola\_id} explícito} em \cd{distribuicaoDaTurma}. É a
    regra inviolável nº 1 --- toda \emph{query} filtra por escola ALÉM da RLS ---
    e aqui a RLS nem existia como rede.
  \item \textbf{\cd{exp\_01\_activities.turma\_id}}, ANULÁVEL e sem
    \emph{backfill} (\arq{20260830090000\_atividade\_por\_turma}). Com turma, o
    painel, a listagem de perfis, a importação da Vila e a entrada do estudante em
    \cd{POST /runs} falam todos da MESMA sala. Sem turma --- que é o caso de toda
    atividade anterior à migração --- vale a escola inteira, o comportamento de
    antes. A turma é conferida contra a MATRÍCULA do professor
    (\cd{professorLeciona}), nunca contra a escola do \emph{token}.
\end{enumerate}

\textbf{Alternativas descartadas.}
\begin{itemize}
  \item \emph{Só um seletor de turma na tela, sem coluna}: o professor
    escolheria a sala a cada visita e nada registraria para qual turma a
    atividade foi feita. Pior: \cd{POST /runs} não teria como saber, e o estudante do
    8B continuaria entrando na atividade do 8A.
  \item \emph{\cd{turma\_id} NOT NULL, com \emph{backfill}}: escolher uma turma
    para as atividades que já existem afirmaria uma decisão que nenhum professor
    tomou. É a mesma regra do \cd{inclinacao\_inicial} em
    \arq{20260829120000\_conclusao\_primeira\_e\_ultima}: NULL quer dizer
    ``anterior à mudança'', e isso é informação.
  \item \emph{Confiar só na RLS para o vazamento}: \cd{users} não tem política, e
    não deve ter --- ela é lida no \emph{login}, antes de se saber a escola. A
    tranca correta é o \cd{where}.
  \item \emph{Uma atividade por turma, obrigatória}: transformaria a criação de
    atividade num laço para o professor que dá a mesma aula em quatro salas, sem
    que nada no produto peça isso hoje.
\end{itemize}

\textbf{Consequências.} Atividades antigas seguem escolares, e o painel agora
DIZ isso em uma faixa --- antes ele mostrava a lista sem explicar de quem era. O
\cd{POST /runs} ganhou uma leitura a mais (\cd{users.turma\_id}, o cache da
matrícula principal) antes de escolher a atividade; a ordenação põe a atividade
DA SALA à frente da geral quando as duas estão ativas. O professor que criar
atividade sem escolher turma continua no comportamento antigo, sem aviso de
erro: é opção, não omissão.

\textbf{Gatilho de reavaliação.} Uma rede rodar o Instanooria em várias turmas
simultâneas com a MESMA pergunta: aí a pergunta é se a atividade deve ter várias
turmas (tabela de junção) em vez de uma, e a resposta muda o painel inteiro.

\textbf{Este gatilho disparou em 2026-09-05}, pelo outro lado: ver a advertência
no topo desta ADR.

\section{ADR-50 --- O perfil importado guarda a PARTIDA e a VERSÃO do
mapeamento, e não ganha chave estrangeira}
\label{adr:perfil-importado-versao}

\begin{especificacao}{Contexto}
\cd{exp\_70\_sessoes} (chamada \cd{vila\_sessoes} até a \adrr{ADR-20}) não tem
prazo de retenção por uma decisão registrada em
\arq{20260826170000\_vila\_sem\_prazo}: o bruto do jogo é a ÚNICA fonte para
recalcular o perfil quando \cd{VERSAO\_MAPEAMENTO} subir, e o mapeamento é
hipótese de design que vai mudar.

Só que a cópia que alimenta o Instanooria
(\cd{exp\_01\_exp\_70\_perfis\_importados}) guardava \cd{escores},
\cd{dominante}, \cd{origem}, \cd{lado\_forcado} e \cd{medido\_em} --- e nem a
sessão de onde veio, nem a versão sob a qual foi medida. Ao subir a constante,
os \cd{exp\_70\_perfis} seriam recalculados e as cópias ficariam na geração
anterior, decidindo o lado do feed, sem nada em lugar nenhum dizendo isso. É
exatamente a armadilha do \cd{expires\_at}: o registro fica, a regra sai, e quem
lê acredita.
\end{especificacao}

\textbf{Decisão.} Duas colunas anuláveis, sem \emph{backfill}
(\arq{20260830100000\_perfil\_importado\_sabe\_a\_partida}):
\cd{vila\_sessao\_id} e \cd{versao\_mapeamento}. Gravadas no \cd{create} E no
\cd{update} do \emph{upsert} --- se só o \cd{create} gravasse, reimportar não
tiraria a linha do estado ``desatualizada'' e o aviso da tela nunca sumiria. O
painel do professor passa a dizer quantos perfis vieram de uma versão anterior,
como AVISO: perfil de geração antiga ainda é melhor informação que o sorteio.

\textbf{SEM chave estrangeira para \cd{vila\_sessoes}}, e isso é o centro da
decisão. Zerar a turma na Vila apaga sessões e bolhas; com FK, isso levaria em
cascata --- ou travaria com \cd{RESTRICT} --- o perfil que está decidindo o lado
de uma aula em andamento. Hoje zerar a Vila NÃO toca \cd{perfis\_importados}, de
propósito, e a coluna não pode mudar isso: ela é RASTRO da medida, não vínculo
de integridade. Sessão apagada vira um id que não resolve, que é exatamente a
verdade do que aconteceu.

\textbf{Alternativas descartadas.}
\begin{itemize}
  \item \emph{Um valor novo em \cd{origem}} (por exemplo \cd{jogo-v3}):
    \cd{perfis\_importados\_origem\_check} e \cd{runs\_perfil\_origem\_check}
    congelam o vocabulário em dois lugares, e o contrato Zod num terceiro.
    \cd{vila\_sessao\_id IS NULL} já responde ``não veio da Vila desta
    instalação'' sem mexer em nenhum deles.
  \item \emph{Preencher a versão das linhas existentes com a atual}: afirmaria
    uma medida que ninguém fez. NULL é a leitura honesta de ``anterior a esta
    migração, ou preenchido à mão pelo professor''.
  \item \emph{Reimportar sozinho quando a versão sobe}: desfaria em silêncio a
    redistribuição que o professor fez antes da aula. O disparo continua sendo
    dele (ADR do Ato 1), e o que muda é que agora a tela AVISA que vale
    reimportar.
  \item \emph{Uma rotina de reprocessamento agora}: ela não existe nem para a
    própria Vila. Escrever a coluna sem o consumidor seria a armadilha de novo,
    ao contrário --- por isso a coluna nasce COM o aviso na tela que a lê.
\end{itemize}

\textbf{Consequências.} A capacidade de reprocessar continua sendo capacidade, e
não rotina: nada recalcula sozinho. O que existe agora é a informação para
decidir --- o painel sabe dizer que a medida envelheceu. Um perfil cuja sessão
foi apagada mantém o id órfão, e isso é aceito.

\textbf{Gatilho de reavaliação.} A primeira subida real de
\cd{VERSAO\_MAPEAMENTO} com dado de contrato em produção: aí se descobre se o
aviso basta ou se a reimportação precisa virar rotina assistida.

\section{ADR-51 --- A suíte de ponta a ponta é organizada por PAPEL, e a
descrição de cada tela mora no teste que a fotografa}
\label{adr:e2e-por-papel}

\begin{especificacao}{Contexto}
A validação visual eram doze arquivos soltos em \arq{scripts/}, nomeados pela
EXPERIÊNCIA que cobriam (\arq{walkthrough-estudante}, \arq{walkthrough-rede},
\arq{walkthrough-ato1-vila}). Três problemas mediram-se sozinhos.

O primeiro é de cobertura: telas inteiras não tinham dono. Criar uma atividade
do zero, os períodos letivos, a análise de produto, o binding do OneRoster, o
compositor de post do Instanooria --- nada disso era exercitado, porque nenhum
arquivo se chamava por elas.

O segundo é de divergência. Cada roteiro carregava sua própria cópia de
\cd{registrar}, de \cd{entrar}, do vigia de console e do relatório final, e as
doze cópias divergiram: uma tirava \cd{fullPage} e outra não, três reprovavam
por erro de console e nove não. Divergência em código de teste é pior que em
código de produção --- quem lê o resultado não sabe qual regra valeu.

O terceiro é o pior. O deck de evidências nascia de um segundo arquivo
(\arq{testes/.gerador-ppt/roteiro.mjs}, 66\,KB) mantido À MÃO, com a descrição
de cada slide escrita longe do teste que tirava o \emph{print}. Ele divergiu do
que os scripts faziam: descrevia telas que já não eram fotografadas e omitia
passos novos. É a mesma armadilha do \cd{expires\_at} e do
\cd{turmas.consentimento\_em} --- o registro fica, a regra sai, e quem lê
acredita nele.
\end{especificacao}

\textbf{Decisão.} A suíte passa a viver em \arq{scripts/e2e/<papel>/}, com uma
casca compartilhada em \arq{scripts/e2e/lib/roteiro.mjs} e um \emph{runner}
único (\arq{scripts/e2e/rodar.mjs}).

\textbf{Por PAPEL, não por experiência.} Quem usa o produto é uma pessoa com um
crachá, e é o crachá que decide o que ela alcança. Organizar por papel faz a
LACUNA aparecer: ``o gestor não tem suíte de períodos letivos'' é uma frase que
alguém diz em voz alta; ``não existe \arq{walkthrough-periodos}'' não é. A
contrapartida está registrada abaixo.

\textbf{O manifesto é a fonte.} Cada \cd{registrar()} grava, junto com o
\emph{print}, o TÍTULO da tela, o que ela MOSTRA e a regra que ela PROVA, num
\arq{passos.json} por suíte. O deck (\arq{gerar-deck.mjs}) e os vídeos
(\arq{gerar-videos.mjs}) só desenham --- não conhecem o produto. A descrição
passou a morar onde o passo mora, e sair de sincronia deixou de ser possível.

\textbf{Alternativas descartadas.}
\begin{itemize}
  \item \emph{Manter a organização por experiência}: é a que casa com o
    \arq{apps/experiencias/}, e para a exp-70 e a exp-03 seria mais direta. Mas
    ela esconde as telas da PLATAFORMA, que são a maioria e não pertencem a
    experiência nenhuma --- e foram justamente elas que ficaram sem cobertura
    por dois anos de projeto.
  \item \emph{Um processo só para todas as suítes}, mais rápido por reaproveitar
    o navegador: uma suíte que trava levaria as outras vinte junto, e o código
    de saída deixaria de ser por suíte. O \emph{runner} paga o custo do
    \cd{spawn} de propósito.
  \item \emph{Playwright Test} (\cd{@playwright/test}) no lugar da casca própria:
    traz \emph{runner}, \emph{fixtures}, \emph{retries} e relatório HTML de
    graça. O que ele NÃO traz é o manifesto descritivo, que é o objetivo aqui ---
    a saída dele é um relatório de teste, não um documento de evidência com uma
    frase de produto por tela. Adotá-lo significaria manter o roteiro do deck à
    mão de novo, que é o problema que esta ADR resolve.
  \item \emph{Manter o roteiro do deck à mão, só arrumando o que divergiu}:
    arruma a divergência de hoje e garante a de amanhã.
  \item \emph{Deck em PDF ou HTML}: o entregável pedido é \cd{.pptx} --- é o que
    a rede abre, edita e projeta.
\end{itemize}

\textbf{Consequências.} As suítes de dois contextos (o professor conduz o
relógio, o estudante joga) ficam no papel do PROTAGONISTA, e não em ambos: a
missão da exp-01 mora em \arq{estudante/}, ainda que metade dos \emph{prints}
sejam da tela do docente. O passo carrega de qual contexto veio, e o deck e o
vídeo dizem isso na tarja.

A ORDEM de execução passa a importar e está fixada no \emph{runner}: o estudante
produz o estado que o professor lê, e a administração vem por último porque
importa gente e troca senha. Rodar uma suíte isolada continua possível; rodar o
professor antes do estudante devolve telas de resultado vazias.

E há um custo aceito: uma experiência nova não tem mais um lugar óbvio. Ela
entra em \arq{estudante/} se for missão, em \arq{professor/} se for ferramenta
--- e a \arq{AGENTS.md} da pasta \arq{apps/experiencias/} passa a dizer isso.

\textbf{Gatilho de reavaliação.} Um papel novo no produto (o \emph{responsável},
por exemplo, se o acesso da família for adiante), ou a suíte passar de
40\,minutos de execução --- aí a paralelização por papel deixa de ser opcional e
o \emph{runner} precisa crescer.

\section{ADR-52 --- Falha do provedor de IA na destilação é 503 com a causa, não
500 mudo}
\label{adr:destilar-503}

\begin{especificacao}{Contexto}
Quatro lugares chamam \cd{generateJson}, e três decidem o que fazer quando ele
falha: as reações ao vivo caem em lista vazia (o feed fica quieto, como um app em
que ninguém respondeu), e a geração de posts e de personas vira \emph{job} falho
que o painel de progresso mostra. O quarto, \arq{geracao.controller.ts}
(\rota{POST /atividades/:id/destilar}), deixava a exceção subir.

O Nest a convertia em \cd{\{"statusCode":500,"message":"Internal server
error"\}}. E o \emph{log} não salvava: o \cd{ErroDeGeracao} guarda a causa real
em \cd{detalhe}, e o \cd{ExceptionsHandler} imprime apenas a \cd{message} --- que
no \cd{VertexProvider} é a genérica \enquote{Falha ao gerar conteúdo
estruturado.}.

Com \cd{AI\_PROVIDER=mock} isso deixa de ser um caso raro e vira certeza:
\cd{MockProvider.generateJson} lança SEMPRE, por decisão (o mock não sabe
satisfazer schema arbitrário). Ou seja, no modo que existe para subir sem GCP ---
o recomendado em desenvolvimento, e o \emph{default} dos dois \emph{compose} ---
\textbf{criar atividade era impossível, e nada em lugar nenhum dizia por quê}.
Custou uma tarde de investigação: a destilação reproduzida fora da api
funcionava, o \cd{UPDATE} do Prisma funcionava, e a resposta só apareceu num
\cd{ps eww} do processo, onde se via o \cd{AI\_PROVIDER=mock} que o \arq{.env}
não declarava.
\end{especificacao}

\textbf{Decisão.} A chamada entra num \cd{try/catch} e a falha vira
\textbf{503} com \cd{mensagem} legível, acrescida da causa quando existe. A causa
é lida como \cd{detalhe ?? message}: o \cd{VertexProvider} enche o
\cd{detalhe} e deixa a \cd{message} genérica; o \cd{MockProvider} faz o
contrário. Ler só um dos dois deixaria sem diagnóstico exatamente o modo em que a
falha é garantida.

\textbf{503 e não 500} porque a requisição está correta e repetir pode funcionar:
o que falhou é uma dependência. É a mesma distinção que a api já faz em
\rota{/ativar} com o 422.

\textbf{Alternativas descartadas.}
\begin{itemize}
  \item \emph{Um \cd{ExceptionFilter} global para \cd{ErroDeGeracao}}: pareceria
    mais econômico e seria pior --- transformaria em 503 também os três
    chamadores que hoje DEGRADAM de propósito, e a lista vazia das reações ao
    vivo viraria erro na cara da criança no meio da aula. A decisão de degradar é
    de cada chamador, e por isso o tratamento fica em cada um.
  \item \emph{Fazer o \cd{MockProvider} devolver um antagonismo plausível}: seria
    o mock passando a conhecer o domínio, e um dia divergiria dele em silêncio.
    Pior: esconderia o caminho de falha em vez de torná-lo legível --- em
    produção o provedor também cai, e é esse caminho que precisa estar exercido.
  \item \emph{Deixar o 500 e melhorar só o \emph{log}}: resolveria para quem tem
    acesso ao \emph{log} da api. O professor continuaria vendo
    \enquote{Internal server error}, e a suíte de ponta a ponta continuaria
    reportando um sintoma que não aponta para a causa.
\end{itemize}

\textbf{Consequências.} \arq{tests/destilar-ia-indisponivel.spec.ts} trava o
comportamento em cinco asserções, incluindo que a atividade NÃO é movida pela
metade (nem \cd{status}, nem \cd{pergunta\_chave}). O roteiro
\arq{scripts/e2e/professor/02-atividade-do-zero.mjs} passou a falhar em 9\,s com a
causa impressa, no lugar dos 180\,s esperando \enquote{Perguntas propostas pela
IA} --- é a mesma lição do \cd{esperarOuFalharCedo}, agora do lado do servidor.

\emph{Ruim}: a mensagem expõe texto do provedor ao cliente. Aceito porque a rota
é de professor autenticado e a causa é operacional (\enquote{sem
implementação}, \enquote{resposta vazia}), nunca dado de estudante --- e o corte
em 200 caracteres impede que um \emph{stack} inteiro vaze para a tela.

\textbf{Gatilho de reavaliação.} Um quarto consumidor de \cd{generateJson} sem
fallback, ou o dia em que a destilação ganhar um --- por exemplo, um antagonismo
escrito à mão pelo professor sem passar pela IA, que hoje ainda depende dela para
avaliar o equilíbrio.

\section{ADR-53 --- Um ambiente por \emph{branch}: \emph{environments} do GitHub
carregam as \cd{vars}, e a \emph{branch} escolhe a \emph{stack}}
\label{adr:ambiente-por-branch}

\begin{especificacao}{Contexto}
Até agosto de 2026 havia uma \emph{stack} GCP só (\cd{dev}) e o \emph{deploy}
disparava no \emph{push} para \cd{main}, lendo doze \cd{vars} do nível do
\textbf{repositório}. Com a \emph{stack} de \cd{hml} provisionada (projeto GCP
próprio, \emph{workspace} Terraform próprio, WIF \emph{pool} e \emph{deployer}
SA próprios), esse desenho não comporta dois ambientes: \cd{vars} de
repositório são um conjunto único, e \cd{main} não pode significar
\enquote{dev} e \enquote{hml} ao mesmo tempo.

O fluxo de trabalho também mudou: ninguém mais commita direto em \cd{main}.
\emph{Branch} de trabalho entra por PR em \cd{dev}; \cd{dev} promove por PR
para \cd{hml}; e, quando produção existir, \cd{hml} promoverá por PR para
\cd{main}.
\end{especificacao}

\textbf{Decisão.} Um \emph{environment} do GitHub por ambiente (\cd{dev},
\cd{hml}, futuramente \cd{prd}), cada um com o conjunto completo de \cd{vars}
e o \emph{secret} \cd{TOKEN\_SIGNING\_KEY\_PEM} escopados nele. O
\arq{deploy.yml} dispara em \cd{push} para \cd{dev} e \cd{hml}, e cada
\emph{job} de \emph{deploy} declara
\cd{environment: \$\{\{ github.ref\_name == 'main' \&\& 'prd' || github.ref\_name \}\}}
--- a \emph{branch} \textbf{é} o seletor de \emph{stack}, com \cd{main} como a
única exceção de nome. Cada \emph{environment} tem \emph{deployment branch
policy} restrita à própria \emph{branch}: o de \cd{hml} só aceita \emph{job}
vindo de \cd{hml}, então nem um \cd{workflow\_dispatch} de outra \emph{branch}
consegue assumir a identidade WIF do ambiente errado. O
\arq{scripts/configurar-github-ci.sh} deriva o \emph{environment} de
\cd{terraform workspace show} --- a mesma fonte dos valores --- e cria
\emph{environment}, \emph{policy}, \cd{vars} e \emph{secret} de uma vez.

\textbf{Alternativas descartadas.}
\begin{itemize}
  \item \emph{Prefixar as \cd{vars} de repositório} (\cd{DEV\_API\_SERVICE},
    \cd{HML\_API\_SERVICE}) e escolher no \emph{shell}: dobra os nomes, espalha
    \cd{if} por todos os \emph{jobs} e não dá nenhuma trava --- qualquer
    \emph{branch} leria as \cd{vars} de qualquer ambiente. O \emph{environment}
    dá o escopo E a \emph{branch policy} de graça.
  \item \emph{Um \arq{deploy-dev.yml} e um \arq{deploy-hml.yml}}: os dois
    divergiriam em silêncio --- é a mesma armadilha do roteiro de 66\,KB da
    suíte de ponta a ponta, mantido à mão até descrever telas que não existiam.
    Um \emph{workflow} só, parametrizado pelo \emph{environment}, não tem como
    divergir de si mesmo.
  \item \emph{Deploy por \emph{tag}} (\cd{v1.2.3-hml}): desacopla o histórico
    da promoção, mas cria um segundo vocabulário (qual \emph{commit} está em
    \cd{hml}? o da \emph{tag}, não o da ponta de nenhuma \emph{branch}) e
    dispensa o PR de promoção --- que é justamente onde o \emph{diff}
    \cd{dev..hml} fica visível e revisável.
\end{itemize}

\textbf{Consequências.} O bloco \cd{env:} global do \arq{deploy.yml} desceu
para dentro dos \emph{jobs}: \cd{vars} de \emph{environment} só resolvem em
\emph{job} que declara \cd{environment:}, e no nível do \emph{workflow}
chegariam vazias, sem erro. Referenciar um \emph{environment} inexistente o
cria em silêncio (com \cd{vars} vazias), então o \emph{job} \cd{mudancas}
ganhou um passo de guarda que derruba \emph{dispatch} de \emph{branch} que não
é de ambiente. O \emph{cache} do Docker ganhou o ambiente no escopo
(\cd{scope=api-dev}, \cd{scope=api-hml}): \cd{dev} e \cd{hml} não invalidam o
\emph{cache} um do outro a cada \emph{push} alternado. E o critério de
\emph{diff} por \cd{HEAD\^{}} passou a depender do método de \emph{merge}:
\emph{merge commit} e \emph{squash} preservam o primeiro pai como a ponta
antiga da \emph{branch} de ambiente; \emph{rebase} o quebraria, e fica
desabilitado nas opções do repositório.

\emph{Ruim}: são $12 \times N$ \cd{vars} para manter ($N$ = ambientes), e a UI
do GitHub não compara \emph{environments} lado a lado --- um valor desatualizado
só aparece como falha no meio do \emph{deploy}. Mitigado porque nenhum valor é
digitado: o \emph{script} os copia de \cd{terraform output}, e reexecutá-lo é
idempotente.

\textbf{Gatilho de reavaliação.} A chegada de \cd{main}/produção (que testará o
mapeamento \cd{main} \(\to\) \cd{prd} e pedirá \emph{required reviewers} no
\emph{environment}); ou o segundo contrato --- \emph{environments} escalam por
ambiente, não por contrato, e aí a conversa vira um repositório de
\emph{deploy} por contrato ou uma matriz.

\section{ADR-54 --- Pessoa nasce por IMPORTAÇÃO: não existe cadastro avulso de
usuário}
\label{adr:sem-cadastro-avulso}

\textbf{Estado: vigente, com o gatilho acionado} (2026-09-01). O cadastro
unitário passou a existir --- e existe do jeito que o gatilho abaixo prescreve:
um formulário que monta a MESMA linha de CSV e chama o MESMO importador, nunca
um segundo caminho de escrita em \cd{users}. Ver \adrr{ADR-59}, que também
abriu a rota separada de \cd{admin-contrato}. Tudo o mais nesta ADR continua
valendo, inclusive as três alternativas descartadas.

\begin{especificacao}{Contexto}
Quem pergunta \enquote{como eu crio um usuário?} procura, naturalmente, um botão
\enquote{Novo} na tela \rota{/admin/usuarios}. Ele não existe, e a ausência
nunca esteve escrita em lugar nenhum --- estava só implícita no fato de a API
não ter \rota{POST /usuarios}.

O provisionamento desta plataforma não é auto-serviço: a lista de pessoas vem da
REDE (a secretaria, o sistema de matrícula), e a identidade que importa é a
\textbf{matrícula da rede} --- com o CPF como chave que promove uma linha
repetida a transferência entre unidades (\adrr{ADR-28}, \adrr{ADR-32}). Uma
pessoa não é só um \cd{users}: é lotação, matrícula em turma, período letivo e
senha provisória, nascidos juntos na mesma transação.
\end{especificacao}

\textbf{Decisão.} As \textbf{únicas} duas chamadas que criam linha em
\cd{users} são \rota{POST /rostering/importar} (a planilha simples, uma linha
por pessoa) e \rota{POST /rostering/oneroster/users} (o \emph{binding} do
padrão, para a rede que já o exporta). \emph{Desde a \adrr{ADR-59}, some-se a
elas \rota{POST /usuarios/admin-contrato} --- o único papel que não cabe numa
planilha da escola; \rota{POST /usuarios} não é uma terceira, porque delega à
primeira.} As duas exigem \cd{gestor} --- na própria
escola --- ou \cd{admin-contrato}, com \cd{?escolaId=}. Criar UMA pessoa é
importar uma planilha de uma linha: a operação é idempotente pela matrícula (e
pelo CPF), então repetir não duplica nem reseta a senha de quem já entrou.

O que existe depois da criação são as operações de \textbf{balcão}, e essas sim
são unitárias: \rota{PATCH /usuarios/:id} corrige nome, CPF, e-mail, papel e
\cd{ativo}; \rota{POST /usuarios/:id/reset-senha} devolve uma senha temporária
uma única vez.

\textbf{Alternativas descartadas.}
\begin{itemize}
  \item \emph{Um \rota{POST /usuarios} com formulário na tela}: é uma segunda
    porta de escrita sobre o mesmo modelo, e portas de escrita divergem. Na
    primeira regra nova --- o CPF que vira transferência, a matrícula que
    precisa de período letivo aberto, o papel \cd{coordenacao} que normaliza
    para \cd{gestor} --- uma das duas a implementaria e a outra não, em
    silêncio. É exatamente o motivo pelo qual \rota{POST /usuarios/importar-csv}
    sobreviveu apenas como \emph{delegação} para o importador, e não como um
    segundo importador.
  \item \emph{Auto-cadastro pelo próprio usuário}: são menores. Não há prova de
    vínculo com a escola que uma pessoa possa apresentar sozinha, e a
    consequência de aceitar uma seria uma criança dentro do \emph{tenant}
    errado.
  \item \emph{Provisionamento \emph{just-in-time} no primeiro login} (criar a
    pessoa a partir do \emph{token} do IdP): tentador quando o WSO2 entrar, e
    recusado desde já --- o \emph{token} do IdP não carrega turma nem série, e a
    pessoa nasceria sem enturmamento, invisível para o professor e fora de
    qualquer conferência da rede. Quando o WSO2 chegar, ele autentica; quem
    provisiona continua sendo a importação.
\end{itemize}

\textbf{Consequências.} \emph{Ruim}: o gestor que percebe \enquote{esqueci de
incluir a Ana} precisa montar um arquivo para uma linha --- atrito real de
balcão, e a tela hoje só aceita \emph{upload}, sem campo de colar. Mitigado em
parte porque o balcão de verdade (nome errado, CPF trocado, senha perdida,
saída da rede) é editar, não criar, e essas rotas existem.

\textbf{Gatilho de reavaliação.} Se a operação pedir cadastro unitário com
frequência, a saída é um formulário que monta a MESMA linha de CSV e chama o
MESMO importador --- nunca um segundo caminho de escrita em \cd{users}.

\section{ADR-55 --- Instrução de agente: \arq{AGENTS.md} canônico por fronteira,
\arq{CLAUDE.md} como ponteiro de uma linha}
\label{adr:instrucao-de-agente}

\begin{especificacao}{Contexto}
Uma pessoa desenvolve esta plataforma, e boa parte do trabalho passa por agentes
de codificação. O que eles sabem sobre o projeto vem dos arquivos de instrução
que a ferramenta carrega sozinha ao tocar uma pasta --- e esses arquivos viraram
um padrão \emph{de fato}, nunca escrito: \arq{AGENTS.md} por domínio, com um
\arq{CLAUDE.md} de uma linha ao lado.

Padrão que não está escrito não se verifica, e os desvios apareceram. A doutrina
da Vila do Lume --- as regras que protegem a Camada~2 do instrumento --- viveu
meses em \arq{exp-70-vila-do-lume/docs/CLAUDE.md}: o carregamento é por caminho,
então quem editava o jogo em \arq{public/jogo/} nunca a recebia. Ela ainda
mandava rodar dois \arq{scripts/walkthrough-*.mjs} que deixaram de existir
quando a suíte passou a ser organizada por papel (\adrr{ADR-51}) --- instrução
desatualizada, que é pior que instrução nenhuma, porque o agente obedece.
\end{especificacao}

\textbf{Decisão.} Três arquivos, três papéis que não se misturam.
\arq{AGENTS.md} é \textbf{canônico}: o único lugar onde a doutrina de uma pasta é
escrita, e o nome que a maioria das ferramentas lê. \arq{CLAUDE.md} é
\textbf{uma linha}, \cd{@AGENTS.md} --- o Claude Code procura esse nome, o
ponteiro importa o canônico, e a regra existe uma vez só. \arq{README.md}
continua sendo para pessoas: o que a pasta é e como rodar, e ele \emph{não} entra
sozinho no contexto do agente.

Um \arq{AGENTS.md} nasce numa pasta onde se trabalha uma tarefa inteira sem sair
dela --- um app, um pacote, uma experiência, a suíte de ponta a ponta --- e
sempre na RAIZ dela, nunca numa subpasta. Entra nele o que, se você não souber,
você quebra: regra já violada e o que custou, contrato com quem consome, zona
congelada, armadilha que passa no \cd{typecheck}. Não entra o que o código já diz
nem o que o \arq{README.md} ao lado já explica. Onde o \arq{README.md} já é a
doutrina --- \arq{scripts/e2e}, \arq{documentacao-latex} --- o \arq{AGENTS.md} é
curto e serve de \emph{carregador}: as poucas regras invioláveis e a ordem de
abrir o README. Alterar a regra e o arquivo é o mesmo \emph{commit}, e isso entra
na \emph{Definition of Done}.

Decisão irmã, do mesmo movimento: \textbf{varredura de vulnerabilidade pertence à
esteira}, com o segredo e a política de prazos do repositório
(\adrr{ADR-39}), não à sessão do agente. O agente não dispara \cd{snyk} nem as
ferramentas do MCP correspondente; as regras estão em \arq{.claude/settings.json}
(escopo do projeto) e no equivalente do usuário. Varredura na sessão gasta cota
do plano, produz apontamento sem o portão que o cumpre, responde por uma árvore
de trabalho local que não é o que vai a produção --- e recicla falsos positivos
já analisados e registrados, como o \emph{code injection} do \cd{setInterval} da
\emph{bridge}. Quem precisar de uma varredura avulsa a roda na mão.

\textbf{Alternativas descartadas.}
\begin{itemize}
  \item \emph{Um arquivo só, na raiz}: já foi assim, e o arquivo cresceu até o
    ponto em que a regra da api ocupava contexto de quem mexia numa experiência.
    A divisão por pasta existe porque o carregamento é por caminho: cada sessão
    recebe o que é dela.
  \item \emph{\arq{CLAUDE.md} com conteúdo próprio ao lado do \arq{AGENTS.md}}:
    duas fontes de verdade sobre a mesma pasta. Elas divergem, sempre na direção
    de quem parou de ler a outra --- a mesma armadilha do \cd{expires\_at} que
    saiu do \emph{schema} e ficou no registro.
  \item \emph{Só \arq{README.md}, sem \arq{AGENTS.md}}: o README não entra no
    contexto sozinho. Foi o que fez a doutrina da Vila do Lume ficar invisível
    para quem ela tentava avisar.
  \item \emph{Gerar os \arq{AGENTS.md} a partir do código}: o que vale ali é
    justamente o que o código não diz --- por que a medida é 0{,}6 e não 0{,}5,
    o que custou descobrir. Gerador nenhum sabe disso.
\end{itemize}

\textbf{Consequências.} \emph{Ruim}: são vinte pares de arquivos a manter, e
cada um pode envelhecer sozinho --- o custo é real e recai sobre a mesma pessoa.
Mitigado por duas coisas: a proibição explícita de repetir conteúdo entre eles
(cópia é o que ninguém atualiza junto) e o item novo na \emph{Definition of
Done}. \emph{Ruim também}: bloquear o MCP de varredura no escopo do usuário vale
para todos os projetos da máquina, não só para este.

\textbf{Gatilho de reavaliação.} Um \arq{AGENTS.md} que ninguém consultou ao
mudar a pasta que ele governa --- sinal de que ele descreve o que o código já
diz, e deve encolher ou desaparecer. Do lado da varredura: a esteira deixar de
cobrir uma classe de vulnerabilidade que só a análise local pega.

\section{ADR-56 --- O fluxo por ambiente é imposto por \emph{hooks} versionados,
não por \emph{branch protection}}
\label{adr:fluxo-por-hooks}

\begin{especificacao}{Contexto}
A \adrr{ADR-53} definiu o fluxo --- \emph{branch} de trabalho \(\to\) \cd{dev}
\(\to\) \cd{hml} \(\to\) \cd{main} --- e o \emph{deploy} passou a depender dele:
o \emph{merge} do PR é o que sobe código para a escola. A regra existia em prosa
no \arq{AGENTS.md} e nada a sustentava: \cd{git push origin dev} funcionava, e
funcionaria de madrugada, no fim de uma sessão longa, sem portão e sem PR.

A resposta natural seria \cd{branch protection} ou \cd{rulesets}. As duas APIs
respondem \textbf{403 --- \enquote{Upgrade to GitHub Pro or make this repository
public}}: o repositório é privado numa conta de plano \emph{Free}. Tornar
público não é opção (é código de plataforma com dado de menor no desenho), e
pagar por assento para uma pessoa, hoje, não se justifica --- mas a regra
precisa valer \textbf{agora}, porque \cd{hml} já é um ambiente real.

Faltavam também as duas pontas que só aparecem quando produção existe: como se
corrige produção sem furar a fila, e como se diz \emph{qual} versão está lá.
Não havia uma \emph{tag} no repositório.
\end{especificacao}

\textbf{Decisão.} A barreira é local e versionada: \arq{.githooks/}, ativado por
\cd{core.hooksPath} no \cd{prepare} do \arq{package.json} --- ou seja,
\cd{pnpm install} instala. \arq{pre-push} recusa \emph{push} direto em
\cd{dev}, \cd{hml} e \cd{main} (inclusive deleção) e nome de \emph{branch} fora
do padrão; \arq{commit-msg} valida \emph{Conventional Commits} com os escopos
deste monorepo. Os escapes são \textbf{nomeados}
(\cd{PERMITIR\_PUSH\_DIRETO=1}, \cd{PERMITIR\_BRANCH\_FORA\_DO\_PADRAO=1}), e
não a \emph{flag} genérica \cd{-{}-no-verify}: emergência acontece, e é melhor
que ela deixe rastro numa variável que diz o que se está furando.

Junto com isso, duas regras que o fluxo não tinha. \textbf{Promoção para
\cd{main} leva \emph{tag} semver e \emph{release}}, criadas depois do
\emph{merge}: MAJOR só em quebra de contrato com quem está fora do repositório
(\cd{PROTOCOL\_VERSION} da ponte, manifesto, telemetria, contrato do Terraform),
MINOR no caso comum, PATCH em correção. \textbf{\emph{Hotfix} sai de \cd{main},
entra por PR em \cd{main} e volta para \cd{hml} e \cd{dev} por retro-merge} no
mesmo dia. O fluxo operacional inteiro vive em \arq{CONTRIBUTING.md} --- a raiz
do repositório, onde o GitHub o exibe ao abrir um PR.

\textbf{Alternativas descartadas.}
\begin{itemize}
  \item \emph{Deixar só a prosa do \arq{AGENTS.md}}: é o estado que produziu o
    risco. Regra que depende de lembrar dela, às 2h da manhã, não é regra --- e
    o custo de esquecer aqui é \emph{deploy} para escola sem portão.
  \item \emph{Tornar o repositório público para liberar \cd{rulesets}}: trocar
    a superfície de exposição do produto inteiro por um recurso de processo.
    Desproporcional.
  \item \emph{\emph{Husky} + \emph{commitlint} + \emph{lint-staged}}: três
    dependências, uma cadeia de configuração e um \emph{hook} gerado, para o que
    são dois arquivos de \emph{shell} que o repositório inteiro consegue ler. A
    regra da raiz vale aqui (\enquote{simplicidade > sofisticação}), e
    \cd{core.hooksPath} é nativo do git desde a 2.9.
  \item \emph{Um \emph{workflow} de \emph{release} que cria a \emph{tag}
    sozinho no \emph{push} para \cd{main}}: automação que decide o número da
    versão precisa de uma fonte para ele, e a única honesta seria um arquivo que
    alguém edita --- o mesmo trabalho manual, com uma engrenagem a mais para
    falhar. Enquanto a promoção for um ato deliberado de uma pessoa, o comando
    documentado basta.
  \item \emph{\emph{Trunk-based} puro, com \emph{feature flags} no lugar de
    \cd{hml}}: exige disciplina de \emph{flag} e um custo de teste que uma
    pessoa não paga; e \cd{hml} existe porque a rede cliente precisa de um lugar
    para homologar antes da escola.
\end{itemize}

\textbf{Consequências.} \emph{Ruim}: \emph{hook} é conselho, não muralha --- quem
clonar e não rodar \cd{pnpm install} não tem nenhum, e \cd{-{}-no-verify} fura.
A barreira protege contra descuido, não contra intenção, e isso é aceitável
enquanto o time é de uma pessoa; deixa de ser no primeiro colaborador externo.
\emph{Ruim também}: o \cd{prepare} roda em CI e em contêiner, onde o
\emph{hook} não serve para nada --- é inofensivo (\cd{git config} num
\emph{checkout} que existe), mas é ruído.

\textbf{Gatilho de reavaliação.} O primeiro colaborador com \emph{commit} no
repositório, ou a conta virar \emph{Pro}\,/\,organização: nos dois casos as
mesmas regras viram \cd{ruleset} no servidor, e os \emph{hooks} passam a ser
conveniência local (retorno rápido) em vez de barreira única.

% ---------------------------------------------------------------------------
\section{ADR-57 --- Observabilidade nativa do Google Cloud, com OpenTelemetry
como fronteira; nenhuma \emph{stack} auto-hospedada}
\label{adr:observabilidade-nativa}

\textbf{Estado: aceita} (2026-09-02; proposta em 2026-09-01 a partir do
estudo do Capítulo~\ref{cap:observabilidade}). Aceita com a primeira peça
no código: log JSON com \emph{redaction} na API e no hub (fase 1) e
\emph{traces} OpenTelemetry com o exportador direto do Cloud Trace nos dois
processos (fase 2, \emph{traces}); as métricas da aplicação e o erro de
cliente seguem o roteiro do capítulo.

\begin{especificacao}{Contexto}
A plataforma é diagnosticável e não é observável: log em texto sem
correlação, nenhuma métrica de aplicação, nenhum alerta, nenhuma sonda
configurada no Cloud Run, erro dentro do \cd{iframe} sem rastro. O risco
\cd{R-05} exige alerta antes de novembro de 2026. Uma pessoa opera; um
contrato é um projeto GCP; haverá vários contratos. A regra da raiz proíbe
serviço novo sem necessidade demonstrada, e a franquia gratuita do Google
Cloud Observability (50\,GiB de log, 150\,MiB de métrica, 2,5\,M de
\emph{spans}, 1\,M de \emph{uptime checks} por mês) cobre um contrato em
regime com folga, desde que o log de requisição seja amostrado na origem.
\end{especificacao}

\textbf{Decisão.} Os sinais operacionais vão para o Cloud Logging, Cloud
Monitoring, Cloud Trace e Error Reporting \textbf{do projeto do contrato}.
A API e o hub emitem log JSON pelo \emph{stdout} (\cd{pino}, com
\emph{redaction} mecânica e \cd{severity}/\cd{trace} nos campos que o Cloud
Logging reconhece) e \emph{traces} e métricas pelo SDK OpenTelemetry com os
exportadores do Google, \textbf{direto do processo, sem \emph{sidecar}}.
Alertas são poucos e por sintoma (SLO com \emph{burn rate}), em duas
severidades. Tudo o que é recurso --- canais, políticas, \emph{checks},
painel, exclusões, \emph{bucket} de auditoria, \emph{dead letter}, sondas do
serviço --- é Terraform, e passa a constar do contrato em
\arq{docs/INFRA.md}. A visão entre contratos é um \emph{metrics scope}
num projeto de observação sem recursos; log não atravessa projeto.

\textbf{Alternativas descartadas.}
\begin{itemize}
  \item \emph{Grafana + Loki + Tempo + Prometheus auto-hospedados}: uma VM
    central (ou uma por contrato), com retenção, \emph{backup},
    atualização e autenticação por conta própria, recebendo dado de N
    projetos por rede e IAM próprios --- e custando, no mínimo, mais que o
    nativo no pior caso. Para uma pessoa, é a peça que quebra às 2h da
    manhã. Volta se uma rede exigir que os sinais fiquem fora do Google
    Cloud, ou se uma consulta demonstrada não couber no Logs Explorer e no
    Log Analytics.
  \item \emph{SigNoz / OpenObserve / HyperDX} (tudo-em-um OpenTelemetry):
    mesma objeção operacional, com a vantagem de um só binário. Gatilho
    igual ao anterior.
  \item \emph{Grafana Cloud Free} como destino (10\,k séries, 50\,GB de log,
    14 dias): o dado sairia do projeto do contrato para uma conta de
    terceiro sem DPA da rede. Aceitável apenas como \textbf{painel} de leitura
    sobre a fonte de dados Cloud Monitoring, se a ergonomia do Grafana
    compensar --- é leitura, não coleta.
  \item \emph{Datadog / New Relic / Dynatrace}: centenas de dólares por mês
    por contrato, para um time de operação que não existe. A instrumentação
    OpenTelemetry deixa a porta aberta sem pagar por ela agora.
  \item \emph{OpenTelemetry Collector como \emph{sidecar}} (a recomendação do
    Google): justifica-se com vários destinos ou \emph{tail sampling}; hoje
    há um destino e amostragem de cabeça basta. É um contêiner a mais em
    cada serviço de cada contrato, com versão e falha próprias.
  \item \emph{Log de requisição na aplicação} (\cd{morgan}, interceptor):
    duplicaria o log que o Cloud Run já grava, dobrando o sinal mais caro
    para acrescentar nada.
  \item \emph{Continuar como está, com o \emph{checklist} semanal do
    Apêndice~F}: é o estado que produziu o \cd{R-05}. Funciona em \cd{dev}
    com uma pessoa olhando; não funciona com cem escolas às 8h.
\end{itemize}

\textbf{Consequências.} \emph{Boa}: zero peças novas para operar; custo
estimado em cerca de US\$~1 por contrato por mês hoje e US\$~7 quando a
cobrança de política de alerta entrar; um contrato novo herda tudo do
Terraform. \emph{Ruim}: dependência do fornecedor nos \textbf{destinos}
(painel, política, consulta são em formato do Google); mitigada porque a
instrumentação é OpenTelemetry e o log é JSON --- trocar de destino é trocar
exportador e reescrever painéis, não reinstrumentar. \emph{Ruim também}: dois
custos crescem com o tráfego e precisam de teto na origem --- log de
requisição (por isso a exclusão e a amostragem são recursos do Terraform) e
cardinalidade de métrica (por isso \cd{escola\_id} não é rótulo). E a
cobrança anunciada de política de alerta muda a economia de \enquote{uma
política por curiosidade}: a lista fica curta também por dinheiro.

\textbf{Gatilho de reavaliação.} Uma rede exigir sinais fora do Google
Cloud; um segundo operador (plantão, rodízio); um segundo destino
(SIEM, Grafana) ou a necessidade de \emph{tail sampling} --- aí o
\emph{sidecar} entra; ou a franquia estourar em dois meses seguidos apesar da
exclusão e da amostragem.

\textbf{Adendo (2026-09-02): métricas por OTLP, não pelo exportador do
Google.} O texto acima diz \enquote{exportadores do Google}. Vale para os
\emph{traces} (\cd{@google-cloud/opentelemetry-cloud-trace-exporter}, gRPC,
leve). Para as métricas, o \cd{-cloud-monitoring-exporter} depende de
\cd{googleapis} --- \textbf{202\,MB} medidos em \cd{node\_modules} --- para
chamar um método da API de monitoramento; entraria inteiro na imagem da API.
A decisão é exportar métricas em OTLP/JSON direto para
\cd{telemetry.googleapis.com/v1/metrics}, com um cliente de 60 linhas que
pede o \emph{token} ao \cd{google-auth-library} a cada envio
(\arq{observabilidade/otlp-google.ts}) --- o mesmo cliente por onde a API
reexporta os \emph{spans} do navegador. Alternativa descartada: o exportador
OTLP/HTTP do próprio OpenTelemetry, que recebe cabeçalhos estáticos e não
acompanha a expiração do \emph{token} OAuth.

\textbf{O risco foi exercido em \cd{dev} no mesmo dia, e a decisão se
sustentou} --- ao custo de duas descobertas que valem mais que a decisão:

\begin{enumerate}
  \item \textbf{O \emph{resource} precisa de \cd{gcp.project\_id}}, e o
    \cd{gcpDetector} \emph{não} o preenche (ele traz \cd{cloud.account.id},
    que a Telemetry API não aceita no lugar). Sem o atributo, TODA exportação
    volta \cd{HTTP 400 Resource is missing required attribute}, a cada 60\,s
    --- com log e \emph{trace} funcionando, o que torna a falha invisível para
    quem olha o painel e visível só no \cd{diag} do SDK. Corrigido em
    \arq{observabilidade/recurso.ts}, com teste de regressão.
  \item \textbf{A ingestão OTLP grava no domínio
    \cd{prometheus.googleapis.com}}, não em \cd{workload.googleapis.com} ---
    onde o exportador legado gravava, e onde este documento mandava procurar.
    A série sai como
    \cd{prometheus.googleapis.com/\B unooria.telemetria.descartados/\B counter},
    no recurso monitorado \cd{prometheus\_target} (\cd{job} =
    \cd{service.name}, \cd{instance} = a instância do Cloud Run,
    \cd{location} = a região), com \cd{otel\_scope\_name} nos rótulos. Por
    algumas horas isso fez parecer que \emph{nenhuma} métrica estava sendo
    exportada, quando estavam todas. É consequência declarada da migração, não
    defeito: a nota oficial avisa que a mudança de domínio \enquote{quebra
    painéis existentes}. Aqui saiu de graça, porque painel ainda não existe ---
    e é o tipo de custo que se paga barato \emph{antes} da produção.
\end{enumerate}

O caminho de volta (exportador do Google, 202\,MB) segue registrado e não foi
preciso.

\textbf{Segundo adendo (2026-09-02): os exportadores do Google para Node estão
DEPRECIADOS --- os PACOTES, não os produtos.} O SDK avisa no \emph{stdout}:
\cd{DEP\_GCP\_OTEL\_TRACE\_EXPORTER --- deprecated and will be archived after
October 30th, 2026}. A nota oficial
(\cd{GoogleCloudPlatform/\B opentelemetry-operations-js}, \arq{MIGRATION.md})
delimita o escopo: \enquote{\emph{all packages in this repository are
deprecated}} e \enquote{\emph{the underlying Google Cloud services (Cloud
Trace, Cloud Monitoring) remain operational}}. \textbf{O Cloud Trace não vai a
lugar nenhum}: muda o caminhão de entrega (OTLP pela Telemetry API), não o
destino.

\textbf{A migração foi feita no mesmo dia} (\cd{feat/obs-traces-otlp}), e o
prazo --- 30 de outubro de 2026, \emph{anterior} à entrada em produção ---
deixou de ser um risco carregado. Nenhum pacote \cd{@google-cloud/opentelemetry-*}
sobrou na api nem no hub.

O desenho mudou numa coisa em relação ao primeiro adendo, e para melhor: os
sinais do PRÓPRIO processo passam pelos exportadores \textbf{oficiais} do
OpenTelemetry (\cd{@opentelemetry/\B exporter-trace-otlp-http} e
\cd{-metrics-otlp-http}), não por cliente escrito à mão. O que tornou isso
possível foi a opção \cd{headers} aceitar uma \emph{função assíncrona} --- é
ela que resolve o problema que motivou o cliente próprio (o \emph{token} da
ADC expira em uma hora, e cabeçalho estático pararia de exportar sozinho). A
fábrica está em \arq{observabilidade/otlp-google.ts}
(\cd{cabecalhosParaTelemetryApi}); o exportador de métricas escrito à mão foi
apagado, e com ele um defeito que ele tinha e a biblioteca não tem: ele
tratava qualquer 2xx como sucesso, ignorando o \cd{partialSuccess} que o OTLP
devolve quando aceita o lote mas descarta pontos.

O \cd{ClienteOtlpGoogle} continua existindo para o que a biblioteca não sabe
fazer: reencaminhar um corpo OTLP \emph{já pronto}, vindo do navegador
(\rota{POST /observabilidade/v1/traces}).

Saldo de dependências: saíram \cd{@google-cloud/opentelemetry-cloud-trace-exporter}
e o \cd{@grpc/grpc-js} que ela arrastava; entraram os dois exportadores
oficiais, que compartilham \cd{otlp-transformer} e \cd{otlp-exporter-base} com
o que já estava instalado.

% ---------------------------------------------------------------------------
\section{ADR-58 --- Erro de cliente entra pela API; nenhum SDK de terceiro
no \emph{bundle} da experiência}
\label{adr:erro-cliente-pela-api}

\textbf{Estado: aceita} (2026-09-02; proposta em 2026-09-01). O lado do
servidor está no código: \rota{POST /observabilidade/erros-cliente},
\rota{/vitals} e \rota{/v1/traces} (este último não estava no texto original:
com o OpenTelemetry Web dentro da SDK, os \emph{spans} do navegador também
entram pela API, que os valida, reescreve o \cd{resource} e reexporta com a
própria credencial --- o mesmo argumento, estendido). O lado do cliente
também está no código (\cd{feat/obs-sdk}, 2026-09-02): captura na SDK com
dedupe e teto de 5 por sessão, \cd{exp:erro} na ponte, \emph{sourcemaps}
\cd{hidden} publicados em \emph{bucket} privado à parte --- e não na pasta
do \emph{bundle} com ACL, como o texto original supunha: o \emph{bucket}
público tem acesso uniforme, e ACL por objeto não existe lá.

\begin{especificacao}{Contexto}
A classe de falha mais cega da plataforma é a que acontece no navegador do
estudante, dentro de um \cd{iframe} de origem diferente: não há
\cd{window.onerror}, \cd{unhandledrejection} nem mensagem \cd{exp:erro} na
ponte, e os \emph{bundles} não têm \emph{sourcemap}. O caminho padrão da
indústria é um SDK de rastreio de erro (Sentry e semelhantes) dentro do
\emph{bundle}, apontando para um \emph{DSN}. A regra 2 da raiz proíbe segredo
e URL interna em \emph{bundle}; a regra 6 proíbe PII fora de controle. Um
\emph{DSN} não é segredo, mas é um destino: o SDK manda \emph{user agent},
URL, \emph{stack}, \emph{breadcrumbs} e o que mais o erro arrastar para um
processador fora do contrato --- para um erro que aconteceu na tela de uma
criança de uma rede pública, coberta por um DPA que não menciona esse
processador.
\end{especificacao}

\textbf{Decisão.} A \cd{@edu/sdk} captura \cd{window.onerror} e
\cd{unhandledrejection} e envia para \rota{POST
/observabilidade/erros-cliente} da \textbf{própria API}, com o \emph{token}
que já tem: \emph{schema} Zod fechado (\cd{expId}, \cd{versao},
\cd{mensagem}, \cd{stack}, \cd{url}, \cd{userAgent}), 8\,KB no máximo,
\emph{dedupe} por \emph{hash} e no máximo cinco por sessão no cliente, cota
por \cd{(escola, usuario)} na API. A API redige, acrescenta \cd{escola\_id}
do \emph{token} e emite \textbf{uma} linha \cd{severity: ERROR} com
\cd{@type} de \cd{ReportedErrorEvent}; o Error Reporting agrupa por
\cd{(expId, versao)}. O \emph{build} das experiências gera \emph{sourcemaps}
ocultos, publicados com ACL privada ao lado do \emph{bundle}, na mesma pasta
versionada. A ponte ganha \cd{exp:erro}, opcional e compatível com hub
antigo, só para o hub trocar a tela. O hub faz o equivalente com
\arq{error.tsx} e \arq{global-error.tsx}. Web Vitals seguem o mesmo caminho,
amostrados e sem identificador de pessoa.

\textbf{Alternativas descartadas.}
\begin{itemize}
  \item \emph{SDK do Sentry (SaaS) no \emph{bundle}}: melhor produto da
    categoria (\emph{session replay}, \emph{breadcrumbs}, \emph{release
    health}); plano Developer grátis para uma pessoa. Descartado pelo
    destino, não pela ferramenta: dado de menor saindo do contrato sem DPA.
  \item \emph{GlitchTip ou Sentry auto-hospedados no projeto do contrato}:
    resolve o destino, mas é uma peça a mais (Postgres, Redis, workers) por
    contrato --- ou uma central com o mesmo problema de N projetos. Fica
    como saída se o agrupamento do Error Reporting se mostrar insuficiente
    para 48 experiências.
  \item \emph{Só \cd{exp:erro} na ponte, e o hub registra}: cria dependência
    do hub para um sinal que a experiência já consegue mandar sozinha (tem
    \emph{token} e base da API), e perde o erro que acontece antes do
    \emph{handshake} --- justamente o mais comum. A ponte fica com o papel de
    interface, não de transporte.
  \item \emph{\emph{Sourcemaps} públicos no \emph{bundle}}: expõem o fonte
    da experiência --- e, na exp-01, a lógica da revelação que o estudante
    não deve ver antes da hora.
  \item \emph{Não capturar erro de cliente}: é o estado atual, e é a maior
    cegueira da Figura~\ref{fig:cadeia-visibilidade}.
\end{itemize}

\textbf{Consequências.} \emph{Boa}: nenhum terceiro no \emph{bundle}; a regra
6 vale por construção, porque a \emph{redaction} é a da API; a versão da
experiência já está no caminho, então a \emph{stack} é mapeável sem
\emph{release} à parte. \emph{Ruim}: sem \emph{session replay} nem
\emph{breadcrumbs} --- o que se tem é mensagem, \emph{stack} e contexto; um
laço de erro numa experiência é contido pela cota, mas o que passa da cota se
perde de propósito. \emph{Ruim também}: mais uma rota pública autenticada na
API, com Zod e cota, para manter.

\textbf{Gatilho de reavaliação.} O Error Reporting agrupar mal os erros das
experiências (assinaturas instáveis por minificação), ou o diagnóstico exigir
\emph{replay}: nos dois casos a saída é um Sentry-compatível
\textbf{dentro do projeto}, nunca o SaaS.

\section{ADR-59 --- Cadastro unitário existe, e ele DELEGA ao importador;
\cd{admin-contrato} tem rota própria}
\label{adr:cadastro-unitario-delega}

\textbf{Estado: aceita} (2026-09-01). Aciona o gatilho de reavaliação da
\adrr{ADR-54}, que segue válida em tudo o mais.

\begin{especificacao}{Contexto}
A \adrr{ADR-54} previu este dia por escrito: \enquote{se a operação pedir
cadastro unitário com frequência, a saída é um formulário que monta a MESMA
linha de CSV e chama o MESMO importador}. A operação pediu. O atrito que a
própria ADR registrou como consequência ruim --- \enquote{o gestor que percebe
que esqueceu de incluir a Ana precisa montar um arquivo para uma linha} --- é
diário no balcão da escola, e a tela só aceitava \emph{upload}.

O que NÃO mudou é o motivo da decisão original: uma pessoa não é uma linha em
\cd{users}. É lotação, matrícula em turma, período letivo aberto e senha
provisória, nascidos na mesma transação, com o CPF decidindo se aquilo é gente
nova ou transferência (\adrr{ADR-28}, \adrr{ADR-32}). Duas portas de escrita
sobre esse conjunto divergem na primeira regra nova.
\end{especificacao}

\textbf{Decisão.} Existe \rota{POST /usuarios}, e ele \textbf{não escreve em
\cd{users}}: valida a forma do corpo com Zod, pergunta se a identidade já
existe e, não existindo, serializa os campos como UMA linha da planilha
(\cd{planilhaDeUmaPessoa}, na mesma constante \cd{COLUNAS\_MODELO} que o leitor
consome) e chama \cd{importarPlanilha}. Toda a regra de domínio --- CPF com
dígito verificador, série e turno normalizados, login gerado do nome, turma
criada quando falta, lotação, matrícula --- continua existindo em UM lugar só.

Duas diferenças em relação a importar um arquivo de uma linha, e as duas são
deliberadas:

\begin{itemize}
  \item \textbf{Criar é criar.} A importação é idempotente de propósito:
    reimportar a planilha do ano não pode duplicar ninguém. Num formulário
    chamado \enquote{Novo usuário} a mesma virtude vira armadilha --- o gestor
    digita a Ana de novo, a linha antiga é atualizada e a tela diz
    \enquote{criado}. Por isso o cadastro unitário \textbf{pergunta antes}
    (CPF na rede; matrícula e login na escola) e responde \textbf{409 sem
    escrever}, apontando o caminho certo: editar, transferir ou importar.
  \item \textbf{CPF de outra escola não revela nada.} O 409 desse caso sai sem
    id, sem nome e sem dizer qual unidade --- é dado de outro \emph{tenant}. A
    transferência continua sendo pela importação, que reconhece o CPF e move a
    matrícula junto.
\end{itemize}

\textbf{E \cd{admin-contrato} ganha rota separada}
(\rota{POST /usuarios/admin-contrato}, \cd{@Roles("admin-contrato")}): ele não
passa pelo importador porque não cabe numa planilha DA ESCOLA --- \cd{escola\_id}
nulo, sem lotação, sem turma, e \cd{PAPEIS\_IMPORTAVEIS} sempre o excluiu.
O login dele é único na REDE inteira, não na escola: o índice
\cd{(escola\_id, username)} não cobre linha com \cd{escola\_id} nulo, e a
autenticação procura o \emph{username} na rede toda --- dois donos do mesmo
login viram recusa no login dos dois.

\textbf{Alternativas descartadas.}
\begin{itemize}
  \item \emph{Escrever direto em \cd{users} a partir do formulário}: é
    exatamente a segunda porta que a \adrr{ADR-54} recusou. Na primeira regra
    nova --- um papel que normaliza, um período letivo que precisa estar aberto
    --- uma das duas aprenderia e a outra não, em silêncio.
  \item \emph{Extrair do importador um método que receba objetos em vez de
    CSV}: mais limpo à primeira vista, e pior no que importa. A validação por
    campo vive no leitor da planilha; um caminho que a contorne precisa
    reescrevê-la em Zod, e aí são duas validações para manter --- a divergência
    de novo, um nível abaixo. O custo aceito é a conversão
    JSON~\seta{}~CSV~\seta{}~pessoa, invisível em latência para uma linha.
  \item \emph{Montar o CSV no NAVEGADOR e postar em \rota{/rostering/importar}}:
    põe a regra de escape RFC 4180 no cliente e faz a api receber um arquivo
    que ela não sabe que veio de um formulário --- sem chance de dar 409 no
    lugar de atualizar.
  \item \emph{\cd{admin-contrato} pela MESMA rota, com \cd{papel} no corpo}:
    metade dos campos ficaria inerte, a autorização viraria um \cd{if} escondido
    no \emph{service} em vez do \cd{@Roles} do método, e a matriz papel x rota
    perderia a linha que hoje se lê sozinha.
\end{itemize}

\textbf{Consequências.} \emph{Ruim}: são duas rotas para \enquote{criar
pessoa}, e quem lê o apêndice C precisa entender por que a segunda existe. A
conversão em três saltos é indireta o bastante para assustar quem depura pela
primeira vez --- daí o comentário no serializador. E o 409 tem uma janela de
corrida de milissegundos: se alguém criar a mesma pessoa entre a pergunta e a
escrita, o importador atualiza (corretamente) e a rota responde 409
\cd{JA\_EXISTIA} em vez de \enquote{criado}, que é a resposta honesta.
\emph{Bom}: a regra de domínio segue em um lugar, o formulário ganha de graça
tudo o que a importação aprendeu, e o gestor deixa de abrir o Excel para
incluir uma criança.

\textbf{Gatilho de reavaliação.} Se o cadastro unitário precisar de um campo que
a planilha não tem --- e que portanto não caiba numa coluna de
\cd{COLUNAS\_MODELO} ---, a delegação deixou de servir: a saída é acrescentar a
coluna ao modelo (e ao leitor), não abrir o caminho direto de escrita.
\section{ADR-60 --- Microsoft Clarity carrega no hub \textbf{e} dentro do
\emph{bundle} da experiência}
\label{adr:clarity}

\begin{especificacao}{Contexto}
A plataforma media comportamento por um caminho só: a telemetria própria da
\arq{@edu/sdk} (\rota{POST /telemetry/batch}), que registra os \emph{verbos}
que cada experiência declara no manifesto --- \cd{attempted}, \cd{interacted},
\cd{progressed}, \cd{scored}, \cd{completed}. Ela responde \enquote{o que a
pessoa fez}, e responde bem, porque é desenhada com o domínio dentro. Ela não
responde \enquote{o que a pessoa \emph{tentou} fazer e não conseguiu}: o clique
repetido num elemento que não é botão, o \emph{scroll} que sobe e desce numa
tela que não explica, a tela abandonada antes do primeiro evento. Nenhum
\emph{verbo} nasce de uma frustração --- só de uma ação que o código previu.

Há uma decisão anterior em sentido contrário, e ela é explícita: a
\adrr{ADR-58} tirou o SDK de erro de terceiro do \emph{bundle} da experiência,
com o argumento de que um \emph{bundle} público manda dado \enquote{para um
processador fora do contrato e do DPA da rede --- para um erro que aconteceu na
tela de uma criança}. Esta ADR \textbf{abre exceção} a ela, para analytics de
produto, e o faz de olhos abertos: o argumento da \adrr{ADR-58} continua
verdadeiro, e as consequências abaixo o repetem.

O pedido foi cobertura total: hub e missões. A restrição técnica que decide o
desenho é que o \cd{iframe} da experiência é \textbf{cross-origin} (o hub sai
do Cloud Run, os \emph{bundles} do \emph{bucket}). Um \emph{script} de gravação
carregado no hub não enxerga nada dentro dele: grava um retângulo vazio
exatamente onde o estudante passa a maior parte da sessão. Cobrir a missão
exige carregar o \emph{script} \textbf{dentro} do \emph{bundle}.
\end{especificacao}

\textbf{Decisão.} O \emph{snippet} oficial do Microsoft Clarity carrega nas duas
superfícies, a partir de \textbf{uma} fonte:
\arq{packages/sdk/src/clarity.ts}, que exporta o corpo do \emph{snippet} e as
chamadas tipadas (\cd{identify}, \cd{set}, \cd{event}, \cd{upgrade}), todas
\emph{no-op} quando \cd{window.clarity} não existe --- bloqueador de anúncio e
rede de escola barrando domínio são o caso comum, não a exceção.

\begin{tabularx}{\linewidth}{@{}L{30mm}X@{}}
\toprule
\headrow \thd{Superfície} & \thd{Como carrega}\\
\midrule
Hub & \arq{components/clarity-hub.tsx} monta o \emph{snippet} com
  \cd{next/script} (\cd{afterInteractive}) no \arq{app/layout.tsx}
  \textbf{raiz} --- o único que cobre a \emph{landing} pública, o login, as
  quatro áreas de papel, \rota{/exp/**} e \rota{/dev-embed}.\\
\zebra Identificação & \arq{components/clarity-identidade.tsx}, montado na
  \arq{casca-plataforma.tsx} --- o único componente por onde
  \rota{(plataforma)/**} e \rota{/exp/**} passam os dois. Re-dispara a cada
  mudança de rota: o App Router navega sem recarregar a página, e o Clarity
  pede \cd{identify} por \textbf{página}.\\
Experiência & \cd{createExperiencia()} chama \cd{carregarClarity} depois do
  \emph{handshake}, com o id que chega no \cd{config} do \cd{hub:init} ---
  campo aditivo, sem mexer no \cd{PROTOCOL\_VERSION}.\\
\bottomrule
\end{tabularx}

O id vai por \textbf{uma} variável, \cd{NEXT\_PUBLIC\_CLARITY\_PROJECT\_ID}, e
\textbf{ausente significa desligado} nas duas superfícies. Ela é inlinada no
\emph{bundle} do \emph{client} durante o \cd{next build}, então entra por
\cd{-{}-build-arg} no \arq{apps/hub/Dockerfile}, alimentado por
\cd{vars.CLARITY\_PROJECT\_ID} do \emph{environment} no \arq{deploy.yml}.

Duas escolhas de conteúdo fecham a decisão. O \cd{identify} manda o
\textbf{UUID opaco} do estudante e nada mais --- o \cd{friendlyName}, que
levaria o nome, fica de fora, e escola e turma vão por \cd{id}. E o
\cd{track()} da \arq{@edu/sdk} passa a alimentar os dois destinos: a telemetria
própria com a carga útil inteira, o Clarity só com \cd{verbo:objeto}. A
\emph{resposta} do estudante não sai do contrato.

\textbf{Alternativas descartadas.}
\begin{itemize}
  \item \emph{Só nas superfícies adultas} (professor, gestor,
    \cd{admin-contrato}): respeita o aviso do próprio Clarity sobre menores e
    não abre exceção nenhuma. Descartada porque deixa de fora o público
    principal --- é no estudante que a dúvida de interface custa aprendizagem.
  \item \emph{Só no hub, missões pela telemetria própria}: manteria terceiro
    fora do \emph{bundle} público. Descartada pelo motivo técnico do contexto:
    o \cd{iframe} é cross-origin, e \enquote{cobrir o hub} deixaria a maior
    parte da sessão do estudante como um retângulo vazio na gravação.
  \item \emph{Pacote \arq{@microsoft/clarity} (npm)}: 7,6\,KB, zero
    dependências, com tipos. É CJS puro --- sem campo \cd{module} --- e o Vite
    das experiências já custou tempo neste repositório com ESM\,/\,CJS.
    Internamente ele só injeta o mesmo \emph{script} que está no
    \arq{clarity.ts}; pagar uma dependência por isso não se justifica.
  \item \emph{\cd{VITE\_CLARITY\_PROJECT\_ID} no \emph{bundle} de cada
    experiência}: seria o precedente do \cd{VITE\_VIDEOS\_BASE\_URL}.
    Descartada porque \cd{import.meta.env} dentro da \arq{@edu/sdk} quebra o
    \emph{build} do hub (que transpila o mesmo arquivo), e porque desligar
    passaria a exigir republicar as seis experiências.
  \item \emph{Um projeto Clarity separado para o \cd{iframe}}: dois painéis
    para uma sessão só. O mesmo id com a \emph{tag} \cd{superficie} separa o
    que precisa ser separado e mantém a costura pelo \cd{identify}.
\end{itemize}

\textbf{Consequências.} \emph{Ruim, e o principal}: a documentação do Clarity
declara que ele \textbf{não deveria ser usado em sites ou aplicações destinadas
a menores de 18 anos}, e a plataforma vai do 1º ano do Fundamental ao 3º do
Médio. Isto é decisão de produto assumida, não descuido, e precisa constar de
qualquer conversa de DPA com uma rede --- ver \cd{R-16} no
Capítulo~\ref{cap:riscos}.

\emph{Ruim}: a gravação de sessão sai do contrato e do projeto GCP da rede. Ao
contrário do log e da métrica, que ficam no projeto do contrato, o
\emph{replay} vai para um processador de terceiro, sob os termos dele.

\emph{Ruim}: as experiências disfarçadas (\cd{exp-01}, \cd{exp-02},
\cd{exp-03}) carregam no \arq{index.html} o aviso de que nada ali pode
denunciar o exercício. Requisições para \cd{clarity.ms} aparecem no DevTools, e
um estudante que abra o inspetor vê um domínio que o clone imitado não teria.

\emph{Ruim}: mexer na \arq{@edu/sdk} republica as seis experiências (o filtro de
mudanças do \arq{deploy.yml} inclui o pacote), então toda alteração aqui custa
um \emph{bump} de versão em seis manifestos.

\emph{Bom}: o \cd{track()} existente virou instrumentação de funil sem tocar em
nenhuma experiência, e o kill-switch é a ausência de uma variável --- desligar
não exige alterar código, e a máquina de quem desenvolve e a suíte e2e já
nascem fora do painel.

\textbf{Gatilho de reavaliação.} Uma rede pública recusar o processamento em
contrato, um parecer jurídico sobre o aviso de menores, ou a gravação deixar de
responder pergunta que a telemetria própria não responda --- nos três casos o
caminho de volta é apagar uma variável de \emph{environment}, e o de frente é
um \emph{replay} auto-hospedado dentro do projeto do contrato.

\section{ADR-61 --- Trabalho paralelo acontece em \emph{worktree}, com a base
escrita à mão}
\label{adr:worktree}

\textbf{Complementada pela \adrr{ADR-71}} (2026-09-05), que responde o que esta
não respondeu: \emph{worktree} separa a árvore de trabalho, não o TRABALHO ---
e \arq{apps/api} e \arq{apps/hub} são comuns a todas as frentes. Lá está a
regra da fatia vertical, a lista de costuras e o \cd{pnpm worktree criar}, que
transforma o passo a passo abaixo em um comando (e escreve o \cd{origin/dev}
que este parágrafo pede que se escreva à mão).

\begin{especificacao}{Contexto}
O fluxo da \adrr{ADR-53} e da \adrr{ADR-56} descreve \emph{branches}, não
diretórios: pressupõe uma pessoa, uma cópia do repositório, uma tarefa por vez.
Passou a haver mais de um agente trabalhando ao mesmo tempo no mesmo
repositório, e alternar de \emph{branch} na mesma cópia não serve para isso ---
há uma árvore de trabalho só, e ela carrega as alterações não commitadas de
quem estava antes. Duas tarefas em paralelo na mesma cópia produzem um estado
misturado, e é dele que sai o \emph{commit}.
\end{especificacao}

\textbf{Decisão.} Cada frente de trabalho paralela ganha um \emph{worktree} em
\arq{.claude/worktrees/<nome>}, criado com a base \textbf{escrita
explicitamente} --- e a base é \cd{origin/dev}.

Um clone antigo deste repositório carrega \cd{origin/HEAD} apontando para
\cd{origin/main} (a \emph{branch} padrão no GitHub já é \cd{dev}; o ponteiro
local não acompanha sozinho), então qualquer ferramenta que \enquote{parta da
\emph{branch} padrão} abre a \emph{branch} sobre produção, e a \adrr{ADR-53}
diz que trabalho nasce de \cd{dev}. O erro é silencioso: a \emph{branch}
compila, os testes passam, e o PR mostra um diff cheio de \emph{commits} que já
estavam em \cd{dev}. A raiz se conserta uma vez por clone
(\cd{git remote set-head origin -{}-auto}); escrever a base no comando é o que
evita depender do \emph{default}, consertado ou não.

\textbf{O paralelismo é de edição e teste unitário; o \emph{runtime} é
\emph{singleton}.} Portas fixas com \cd{strictPort}, \emph{compose} de nome
fixo e um banco com um \emph{schema}: dois \emph{worktrees} não rodam a
\emph{stack} ao mesmo tempo, e isso é deliberado --- é o que faz \cd{pnpm dev}
falhar alto em vez de servir a experiência errada. No \emph{worktree}: editar,
\cd{typecheck}, \cd{lint}, os testes unitários e de contrato, o \emph{build},
e o Vite de \textbf{uma} experiência em porta alternativa, embedado pelo hub
que já está no ar (\rota{/dev-embed?src=}). Na cópia principal, uma frente por
vez: a \emph{stack} inteira, a suíte de ponta a ponta e as migrações. A
\emph{Definition of Done} exige a suíte do papel afetado para toda tela ---
numa frente em \emph{worktree} isso significa trazer a \emph{branch} para a
cópia principal e rodá-la de lá, um agente por vez.

\arq{.gitignore} já cobre o caminho (\cd{.claude/*}), então o \emph{worktree}
não aparece no estado da cópia principal. O ciclo de \emph{commit}, PR e
promoção não muda em nada --- o \emph{worktree} é onde o trabalho acontece, não
um fluxo novo. O passo a passo está em \arq{CONTRIBUTING.md}, e desde a
\adrr{ADR-71} ele é um comando: \cd{pnpm worktree criar}.

\textbf{Alternativas descartadas.}
\begin{itemize}
  \item \emph{Alternar de \emph{branch} na mesma cópia}: é o estado que
    produziu o problema. Serializa o trabalho e mistura alteração não
    commitada de tarefas diferentes.
  \item \emph{Clones separados do repositório}: funciona, e custa uma cópia
    inteira do histórico e um \cd{pnpm install} do zero por frente. O
    \emph{worktree} compartilha o diretório de metadados e só duplica o que
    precisa ser duplicado.
  \item \emph{\emph{Worktree} fora de \arq{.claude/}}: exigiria uma entrada
    nova no \arq{.gitignore}, e um esquecimento ali significa \emph{worktree}
    inteiro entrando num \cd{add -A} --- o mesmo acidente que já custou 250\,MB
    de \emph{cache} de \emph{webpack} no histórico.
\end{itemize}

\textbf{Consequências.} \emph{Ruim}: \cd{node\_modules} e o \emph{client} do
Prisma são por \emph{worktree}, então cada frente paga um \cd{pnpm install} (e
um \cd{prisma generate}) na criação. \emph{Ruim}: \cd{pnpm docker:dev} de
dentro de um \emph{worktree} recria os contêineres da \emph{stack} principal
--- o nome do projeto do compose vem do diretório, e os dois disputam as mesmas
portas. Quem precisa da \emph{stack} inteira usa a cópia principal.
\emph{Ruim}: quatro arquivos que toda frente toca brigam no \emph{merge}
--- a numeração de ADR (dois agentes já escreveram a mesma \enquote{ADR-N}
ao mesmo tempo: reserve o número conferindo \cd{origin/dev}), o
\arq{main.pdf} binário (resolve-se regerando, nunca escolhendo um lado), o
\arq{AGENTS.md} da raiz e o \arq{pnpm-lock.yaml}. \emph{Ruim}: cada
\emph{merge} em \cd{dev} dispara um \emph{deploy} que não cancela o anterior;
PRs em rajada viram \emph{deploys} em fila. \emph{Bom}: a árvore principal
fica intocada, com o que estiver em andamento nela preservado.

\textbf{Gatilho de reavaliação.} O trabalho paralelo deixar de acontecer, ou o
repositório ganhar CI que rode as frentes em contêiner --- nos dois casos o
\emph{worktree} vira conveniência local em vez de regra. E o inverso: a fila
da suíte de ponta a ponta na cópia principal virar gargalo \textbf{medido} (duas
frentes esperando uma pela outra, mais de uma vez) é o que justifica
parametrizar porta e nome do \emph{compose} por \emph{worktree} e um banco por
\emph{branch} --- não antes.

\section{ADR-62 --- O WAF não inspeciona o corpo de rota autenticada nem de
texto livre de estudante}
\label{adr:waf-excecoes}

\begin{especificacao}{Contexto}
O Cloud Armor da \adrr{ADR-57} entrou com sete regras pré-configuradas do
OWASP CRS 3.3 em \cd{sensitivity 1}, e em \cd{dev} já saiu de \emph{preview}
--- bloqueando de verdade. Nos primeiros dias, \textbf{242 requisições}
\rota{POST /telemetry/ingest} responderam 403 com
\cd{body\_denied\_by\_security\_policy}, todas do agente \cd{APIs-Google}: é o
\emph{push} do Pub/Sub, negado pela assinatura
\cd{owasp-crs-v030301-id921150-protocolattack} --- \enquote{HTTP Header
Injection via payload (CR/LF detectado)}.

O envelope do \emph{push} é um JSON com a mensagem em \emph{base64}, e a 921150
casa qualquer corpo que carregue \cd{\textbackslash r\textbackslash n}. O
prejuízo é silencioso duas vezes: a \emph{subscription} reentrega para sempre
(não há DLQ), e a tela de resultados do professor fica vazia sem erro em lugar
nenhum --- a requisição morre no \emph{load balancer}, antes da api, então não
há log de aplicação a consultar.

A mesma assinatura ameaça a exp-03: a pergunta do estudante é texto livre, e
\emph{Shift+Enter} ou colar de outra aba produzem quebra de linha legítima. O
403 do WAF não carrega cabeçalho de CORS, então o navegador reporta falha de
rede: o estudante lê \enquote{A conexão falhou} e não há nada no servidor
apontando para o WAF.
\end{especificacao}

\textbf{Decisão.} Duas exceções cirúrgicas na política, e nenhuma regra
desligada.

\begin{enumerate}
  \item Uma regra \cd{allow} em \textbf{priority 900} --- antes das WAF, que
    começam em 1000 --- casando
    \cd{request.path.startsWith('/telemetry/ingest')}. Abrir o caminho não abre
    a rota: ela exige o \emph{token} OIDC do Pub/Sub
    (\env{TELEMETRIA\_PUSH\_SA} e \env{TELEMETRIA\_PUSH\_AUDIENCE}) e responde
    401 sem ele. O WAF ali era redundância que só custava dado.
  \item Um \cd{preconfigured\_waf\_config} na regra \cd{protocolattack} (1006)
    excluindo \textbf{apenas} o \cd{target\_rule\_id 921150} nos prefixos
    \rota{/unoogpt/} e \rota{/ai/} --- as rotas que recebem texto livre de
    estudante. As outras seis assinaturas de protocolo continuam valendo
    nesses caminhos, e a 921150 continua valendo em todo o resto da api.
\end{enumerate}

\textbf{Alternativas descartadas.}
\begin{itemize}
  \item \emph{Desligar a regra \cd{protocolattack} inteira}: resolveria os dois
    sintomas com uma linha a menos, e removeria \emph{smuggling} e
    \emph{splitting} da superfície inteira da api por causa de duas rotas.
  \item \emph{Voltar a política para \emph{preview}}: transforma um bloqueio
    visível em bloqueio nenhum, e junto some a proteção contra os
    varredores que já batem no IP do \emph{load balancer} todos os dias.
  \item \emph{Fazer a api aceitar o \emph{push} por outro caminho} (endpoint
    interno, fora do LB): exigiria expor a api sem o LB ou montar rede nova
    para um único endpoint --- muito mais superfície do que a exceção.
  \item \emph{Sanitizar o CR/LF antes de enviar} (no SDK e no publicador da
    telemetria): trata o sintoma no emissor errado. A quebra de linha na
    pergunta do estudante é conteúdo legítimo, e removê-la mudaria o dado que a
    missão mede.
\end{itemize}

\textbf{Consequências ruins.} \rota{/telemetry/ingest} passa a não ter WAF
nenhum --- se o \emph{guard} de autenticação daquela rota regredir, não há
segunda camada. E as exclusões são por prefixo de caminho: uma rota nova sob
\rota{/ai/} herda a exclusão sem que ninguém decida isso.

\textbf{Gatilho de reavaliação.} A telemetria deixar de entrar por \emph{push}
no LB (um \emph{worker} separado, como o desenho original previa) devolve o
caso 1 ao WAF. E o Cloud Armor publicar exclusão por \emph{método} além de
caminho permitiria estreitar o caso 2 para os \cd{POST} que de fato carregam
texto do estudante.

\section{ADR-63 --- O CI não publica a mídia pesada, mas se recusa a publicar
sem ela}
\label{adr:midia-externa}

\begin{especificacao}{Contexto}
Mídia acima de 2\,MB não entra no git nem no \emph{bundle}: os 21 \emph{mp4}
da exp-02 (52\,MB) vivem no \emph{bucket} de \emph{assets} do contrato, e o
\emph{build} recebe a base pública em \env{VITE\_VIDEOS\_BASE\_URL}. O
\emph{upload} sempre foi manual, uma vez por contrato, e isso estava registrado
como lacuna conhecida em \arq{docs/INFRA.md}.

Em 2026-09-01 a lacuna cobrou. A exp-02 foi publicada no contrato
\cd{entre-redes-dev} com o \emph{bucket} de \emph{assets} vazio: os 21 arquivos
responderam 404, e como o \cd{ShortsPlayer} do \emph{design system} deriva
\cd{waiting} só de \cd{waiting}/\cd{playing} --- sem escutar \cd{error} --- o
evento \cd{playing} nunca chegou e o feed ficou em \enquote{Carregando o
vídeo\ldots} para sempre. Salvar e Curtir são desabilitados enquanto o vídeo
carrega (para não salvar sem assistir), então a missão ficou \textbf{impossível
de concluir}. Nada falhou: nem o \emph{deploy}, nem o log da api. Quem
descobriu foi um estudante, na frente da tela.
\end{especificacao}

\textbf{Decisão.} O CI continua \textbf{não} publicando a mídia --- o
\emph{runner} faz \emph{checkout} do git, e os arquivos não estão lá --- e
passa a \textbf{conferi-la}. Uma experiência com mídia externa declara um
\arq{midia-externa.json} versionado ao lado do \arq{package.json}, com o
prefixo no \emph{bucket} e a lista de arquivos esperados. Antes de subir o
\emph{bundle} daquela experiência, o \emph{job} \cd{publicar-exps} lista o
prefixo no \emph{bucket} de \emph{assets} e \textbf{falha o deploy} se faltar
algum, imprimindo o que falta e o comando de \emph{upload}.

Isso exige a var \env{ASSETS\_BUCKET} no \emph{environment} do GitHub, escrita
por \arq{scripts/configurar-github-ci.sh} a partir do \emph{output}
\cd{bucket\_assets} do Terraform. Sem a var, o passo apenas avisa --- um
\emph{environment} ainda não reconfigurado não deve travar \emph{deploy} por
uma verificação que ele não tem como fazer.

Em paralelo, e independente disto, a experiência passou a tratar o \cd{error}
do \cd{<video>} por conta própria: vídeo que falha sai do estado de
carregamento, as ações voltam e a missão continua concluível. As duas coisas
respondem à mesma falha em camadas diferentes --- uma impede publicar quebrado,
a outra impede que o estudante fique preso quando algo quebrar mesmo assim.

\textbf{Alternativas descartadas.}
\begin{itemize}
  \item \emph{Versionar os \emph{mp4} no git}: 52\,MB por experiência, num
    repositório que já teve 250\,MB de \emph{cache} de \emph{webpack} entrando
    por um \cd{add -A} distraído. Para 48 experiências, o histórico ficaria
    impraticável.
  \item \emph{Git LFS}: resolveria o peso e acrescentaria um serviço com cota,
    um passo de \emph{checkout} novo e uma forma nova de o \emph{deploy}
    falhar --- por mídia que muda uma vez por contrato.
  \item \emph{O CI baixar a mídia de uma origem canônica e republicar}: exige
    escolher e operar essa origem (outro \emph{bucket}, \emph{release assets}),
    e transfere o problema em vez de resolvê-lo --- alguém continua tendo que
    colocar o arquivo lá na primeira vez.
  \item \emph{Só documentar melhor o passo manual}: já estava documentado. A
    lacuna estava escrita em \arq{docs/INFRA.md} e mesmo assim o contrato subiu
    sem a mídia.
\end{itemize}

\textbf{Consequências ruins.} O \emph{deploy} de uma experiência com mídia
passa a depender de um recurso que o \emph{job} não controla: \emph{bucket}
renomeado ou permissão de leitura retirada da \emph{service account} do CI
quebra a publicação de um \emph{bundle} que estava correto. E a lista de
arquivos no \arq{midia-externa.json} é mantida à mão --- acrescentar um vídeo
sem acrescentar a linha faz a conferência passar sem conferir.

\textbf{Gatilho de reavaliação.} A terceira experiência com mídia externa: com
uma, a lista à mão é honesta; com três, o \arq{midia-externa.json} deve passar
a ser gerado do catálogo da api, que já sabe o nome de cada arquivo.

\section{ADR-64 --- O marco da exp-03 carrega o relógio do instante; a faixa
da decisão é gravada, não recomputada}
\label{adr:marco-com-relogio}

\begin{especificacao}{Contexto}
O pedido do produto para o relatório do professor da Missão Tempestade foi
sobre \emph{tempo}: quanto demorou até a primeira fala, com quanto cronômetro
cada fonte foi carregada, quantas vezes mudou o destino antes de confirmar, com
quanto relógio confirmou (ou se estourou), e se apertou Recomeçar --- e se o
professor \enquote{tinha liberado}. Ao levantar o que já existia, duas
imprecisões de gravação apareceram. A primeira: \cd{sessoes.minutos\_gastos}
só era atualizado pelo relógio dentro do turno do assistente; as rotas diretas
(\cd{alocar}, \cd{carregarConjuntos}, \cd{simular}) e as ferramentas chamadas
de dentro do turno gravavam o marco com o minuto do \textbf{último turno} ---
e um turno leva de 10 a 30\,s, dois a seis minutos da cidade. Pior: a hora de
\emph{saída} das equipes vinha do mesmo retrato, então quem alocava pelo mapa e
enviava sem falar de novo saía no minuto de um turno atrás, e chegava a bairros
que já não alcançaria. A segunda: a fala do estudante era gravada depois do
\cd{await} do provedor, datada pelo fim do turno, não pelo envio.

Havia ainda o que o relatório precisaria e não tinha como reconstruir: o
cronômetro de parede que a criança via num dado marco. \cd{criado\_em} existe,
mas \enquote{+2 min} altera \cd{duracao\_seg} sem histórico e \enquote{zerar a
turma} apaga o relógio --- a conta feita depois mentiria em toda aula que teve
extensão. E a faixa da decisão (200 corridas de \cd{esperadoDe}) era
recomputada a cada leitura: determinística hoje, mas uma mudança em
\cd{PARAMETROS} reescreveria em silêncio o relatório de toda partida antiga.
\end{especificacao}

\textbf{Decisão.} Cinco movimentos, uma migração só de \cd{ALTER}, sem tocar
RLS nem \cd{escola\_id}:

\begin{enumerate}
  \item \cd{gravarMarco} recebe o \cd{EstadoDoRelogio} \textbf{recalculado no
    instante} (\cd{relogioAgora}), e grava \cd{minuto},
    \cd{segundos\_restantes} e \cd{motivo\_do\_fim} dele. As rotas diretas
    passam a atualizar \cd{minutos\_gastos} pelo mesmo estado, e
    \cd{aoConfirmar} usa esse minuto para a saída das equipes.
  \item A fala do estudante leva \cd{criadoEm: agora} --- a hora em que apertou
    enviar.
  \item Dois tipos de marco: \cd{inicio} (entrou na Investigação, uma vez por
    tentativa --- é de onde se mede o tempo até a primeira fala) e
    \cd{recomeco}, gravado na tentativa que \emph{fecha}, com \cd{origem}
    (qual botão: a lixeira da Investigação ou o \enquote{Tentar de novo} do
    Resultado --- a única parte que só o cliente sabe, e vai no corpo do
    \rota{POST /unoogpt/reiniciar}), \cd{fase} e \cd{modo} do relógio da turma.
  \item \cd{exp\_03\_simulacoes.esperado}: a faixa e o percentil do resgate
    real, gravados no confirmar. Precedente: \cd{exp\_70\_sessoes.leitura\_mensagens}.
    O \cd{envio} passa a carregar também as justificativas como o professor lê,
    porque \cd{reiniciar} zera as da sessão.
  \item O resto é \textbf{derivado na leitura}, sem coluna: mudanças de destino
    (pares \cd{(alfa, bravo)} consecutivos distintos antes do envio --- contar
    linhas mentiria, porque o cartão de duas opções manda um POST por equipe e
    justificar passa pela mesma rota), tempo até a primeira fala (da largada e
    do Começar, os dois, porque antes da largada o campo está bloqueado e
    contar do Começar mediria o professor), \% do risco coberto e posição na
    faixa. Mesmo princípio de \cd{leitura.ts}: o servidor recalcula do bruto.
\end{enumerate}

\enquote{Se o professor tinha liberado} recebe a resposta honesta: não existe
liberação de recomeço --- o que o professor controla é o relógio, e é o
relógio que fica gravado ao lado do clique. Criar a liberação seria decisão de
produto contra o desenho registrado (\enquote{o professor controla quando dá
para falar, não o que a criança decide}) e fica para uma ADR própria, se o dado
gravado mostrar que faz falta.

\textbf{Alternativas descartadas.}
\begin{itemize}
  \item \emph{Telemetria como fonte}: os eventos já existem
    (\cd{interacted:alocacao}, \cd{attempted:recomeco}). Mas o canal é
    preenchido pelo cliente, assíncrono (Pub/Sub, desligado em \cd{dev}),
    perde lote no fechar da aba, e nenhum código do repositório lê
    \cd{telemetria\_eventos} --- além da regra escrita de que o que o professor
    lê nunca vem do canal que o cliente preenche.
  \item \emph{Reconstruir o cronômetro de \cd{criado\_em} + relógio da turma}:
    falha em toda aula com \enquote{+2 min} (a duração muda sem histórico) e
    em toda turma zerada (o relógio some). Gravar no instante custa duas
    colunas e não mente.
  \item \emph{Sessão por tentativa} (uma linha de \cd{exp\_03\_sessoes} por
    rodada, preservando justificativas e conjuntos): resolveria a perda no
    \cd{reiniciar} de forma geral, mas mudaria a chave primária de quatro
    tabelas e toda consulta do módulo. A cópia das justificativas no marco de
    \cd{envio} resolve o que o relatório precisa, por uma fração do custo.
  \item \emph{Persistir a contagem de mudanças e a cobertura}: contrariaria o
    princípio do módulo (o servidor recalcula do bruto, e a versão da leitura
    sobe quando a hipótese muda). Só se grava o que \textbf{não} dá para
    recalcular --- o cronômetro e a faixa.
\end{itemize}

\textbf{Consequências.} \emph{Ruim}: a correção do minuto muda o desfecho de
quem confirma sem turno novo --- é correção, mas partidas antigas e novas não
são comparáveis nesse detalhe, e o produto foi avisado antes de qualquer teste
controlado. \emph{Ruim}: marcos e simulações anteriores à migração têm as
colunas novas nulas; a tela diz \enquote{sem registro} e a leitura recomputa a
faixa. \emph{Ruim}: o cliente da exp-03 sobe de versão pela \cd{origem} no
corpo --- cliente antigo continua aceito (corpo vazio é válido). \emph{Bom}:
nenhum evento de telemetria novo, nenhuma tabela nova, RLS intocada; o
relatório do professor sai inteiro de \cd{GET /unoogpt/professor/estudantes/:id},
e o gestor continua sem rota que aceite \cd{estudanteId}.

\textbf{Gatilho de reavaliação.} O produto pedir a comparação de
\emph{justificativas} ou \emph{conjuntos} entre tentativas além do que o marco
de envio guarda --- aí a sessão por tentativa volta à mesa. Ou o dado de
\cd{recomeco} mostrar que crianças recomeçam depois de o professor encerrar,
em número que justifique a liberação explícita.

\section{ADR-65 --- O guia da aula é dado renderizado por um componente único,
com blocos vivos para o que tem fonte de verdade em outro lugar}
\label{adr:guia-como-dado}

\begin{especificacao}{Contexto}
Só o Instanooria tinha guia da aula. Ele era uma página
(\arq{.../[ano]/[unidade]/guia/page.tsx}), 447 linhas de JSX com \cd{notFound()}
para qualquer outro \cd{expId}: conteúdo, ordem e marcação na mesma camada. As
outras quatro missões publicadas --- UnooTube, Missão Tempestade, O Eco no Grupo
e a Vila do Lume --- chegavam ao professor com a ficha do catálogo (competência,
atos, preparo) e nada sobre como conduzir a aula: nenhuma pergunta de debate,
nenhuma mediação, nenhuma rubrica. A Vila tinha um guia \emph{do debate},
excelente e parcial: cobria o pós-jogo e não o preparo nem o fechamento.

E o problema não era só de cobertura. Duas partes do guia do Instanooria já
vinham de fonte viva por decisão anterior --- a tabela dos momentos de
\cd{CONTROLE\_DO\_PROFESSOR} e as telas da revelação de
\cd{@edu/conteudo-exp-01} --- porque a versão escrita à mão tinha divergido: o
guia falava em \enquote{os seis botões} e o painel tinha sete. A Vila serve os
cartões da api por turma pela mesma razão. Qualquer solução para as outras
missões precisava carregar esses casos junto.
\end{especificacao}

\textbf{Decisão.} O guia vira \textbf{dado}
(\arq{apps/hub/lib/entre-redes/guias/<exp>.ts}, tipo \cd{GuiaDaAula}) renderizado
por um \textbf{componente único} (\arq{apps/hub/components/guia/}), com a mesma
ordem de seções em toda missão: objetivos e BNCC Computação, antes da aula
(técnico, pedagógico, o que não dizer), roteiro cronometrado com a coluna
\enquote{o que você faz e fala}, banco de perguntas em quatro níveis, folha de
observação, mediação, depois da aula (síntese, registro, rubrica,
desdobramentos), cuidados, adaptações e dúvidas frequentes.

O que tem fonte de verdade em outro lugar entra por \textbf{bloco vivo}: um
\cd{ReactNode} que a página da missão ancora ao fim de uma seção. É por aí que
entram os momentos e a revelação do Instanooria, os cartões da Vila (por turma)
e a tabela dos oito vieses da Missão Tempestade --- esta última por uma rota
nova, \rota{GET /unoogpt/professor/fontes}, com o viés \textbf{derivado} da
instrução do agente em vez de escrito duas vezes.

Uma folha de impressão acompanha o renderizador: o professor conduz a aula com o
guia na mão enquanto o computador projeta a tela da turma.

\textbf{Alternativas descartadas.}
\begin{itemize}
  \item \emph{Uma página por missão, como a do Instanooria}: cinco arquivos de
    cerca de 450 linhas, sem nada que os obrigue a ter as mesmas seções. A seção
    que alguém esquecesse de copiar não existiria para aquela turma, sem sinal
    nenhum --- e a segunda missão que o professor rodasse teria a informação em
    outro lugar da página.
  \item \emph{Guia em Markdown ou MDX}: resolveria a escrita e perderia as duas
    coisas que importam --- o tipo que cobra seção faltando e os blocos vivos,
    que precisam de componente.
  \item \emph{Guia no catálogo} (\arq{lib/entre-redes/catalogo.ts}): o catálogo
    é a grade, lido por meia dúzia de telas; o guia é material de aula, pesado, e
    só a rota do guia precisa dele.
  \item \emph{Copiar a tabela dos vieses para o hub}, evitando a rota nova: era
    a opção mais barata e é exatamente a que a doutrina do repositório proíbe.
    Recalibrado o cenário, o professor contaria à turma um viés que o jogo já não
    produz --- o mesmo erro que os cartões da Vila vindos da api existem para
    evitar.
  \item \emph{Guia por turma} (com \cd{?turmaId=}): só a Vila precisa disso, e
    ela já resolve com seletor próprio dentro da revelação. Genéricos, todos os
    guias passariam a exigir turma para abrir --- e o professor lê o guia na
    véspera, quando ainda não escolheu a sala.
\end{itemize}

\textbf{Consequências.} \emph{Bom}: as cinco missões publicadas passam a ter
guia com a mesma estrutura; \arq{guias/index.test.ts} recusa missão publicada sem
guia, com uma lista de exceções que exige razão escrita. \emph{Bom}:
\arq{lib/rotas-da-missao.ts}, que já era o único tradutor \cd{expId} \seta{}
rota, passa a responder também se a missão tem guia --- então a sala de aula
mostra o \emph{link} em toda turma × toda missão sem inventar URL.
\emph{Ruim}: o conteúdo pedagógico das quatro missões novas é \textbf{autoral},
derivado do código e da \emph{copy} congelada, e não passou por playtest em sala
--- está escrito para ser corrigido por quem der a primeira aula. \emph{Ruim}: o
alinhamento com a BNCC Computação é proposta da plataforma, não classificação
oficial, e a tela precisa dizer isso. \emph{Ruim}: uma rota a mais na api para
uma tela de leitura, com a documentação e o teste que ela obriga.

\textbf{Gatilho de reavaliação.} Se o guia passar a variar por turma em mais de
uma missão --- aí o parâmetro deixa de ser exceção da Vila e vira propriedade do
renderizador. Ou se a rede pedir o guia como PDF/impresso oficial, com marca e
paginação de documento: aí a folha de impressão deixa de bastar e o caminho é
gerar o documento, não estilizar a tela.

\section{ADR-66 --- A exp-07 escreve com IA na tela do estudante, e o gabarito
continua fora do modelo}
\label{adr:exp07-ia-em-runtime}

\begin{especificacao}{Contexto}
A missão \cd{exp-07-o-eco-no-grupo} chegou dos gates de produto com uma decisão
explícita: \enquote{NÃO HÁ USO DE IA GENERATIVA nesta missão. Todas as
mensagens de chat, reações da cascata, áudios e falas de personagens são 100\%
roteirizados e deterministicamente controlados.} As cinco bancas aprovaram
sobre essa premissa, e a banca de dados chegou a registrar como ponto forte a
ausência de campo de texto livre de menor (\cd{DAD-08}).

O produto reabriu a decisão: quer interação por IA na missão. Isso não é um
detalhe de implementação --- muda o que o instrumento mede, o que a criança
escreve e quanto a aula custa.
\end{especificacao}

\textbf{Decisão.} A IA entra em \textbf{três superfícies}, e em nenhuma delas
ela decide resultado.

\begin{enumerate}
  \item \emph{As reações do bairro} (Etapa 4). O desfecho --- se o grupo se
    acalma, entra em pânico ou azeda --- sai de \cd{Opcao.desfecho} no catálogo
    do servidor, antes de qualquer modelo ser chamado. O que a IA escreve é
    \textbf{como isso soa}, reagindo ao que a criança de fato postou.
  \item \emph{Amara no privado}, e só DEPOIS da decisão no grupo. O silêncio
    dela durante a investigação é não-negociável de ofício: se ela apontasse a
    mensagem falsa, o reconhecimento de padrão que a missão mede deixaria de
    acontecer. O portão é do servidor (\cd{amara-em-silencio}), não uma condição
    de tela.
  \item \emph{A devolutiva do \enquote{Meu percurso}}, um parágrafo sobre as
    escolhas. Não é nota: o \cd{score} sai de \arq{regras/pontuacao.ts} sobre o
    que aconteceu no grupo, e o formulário reflexivo não entra nele.
\end{enumerate}

As três falham para o lado seguro. Sem saldo de IA, sem provider ou com a saída
reprovada pelo sanitizador, a missão fecha inteira: as reações caem no roteiro
autoral do catálogo, Amara \enquote{fica sem sinal} e a devolutiva não aparece.
Nenhuma aula depende de o Vertex estar de pé --- e é isso que
\arq{tests/o-eco-no-grupo-postgres.spec.ts} exercita, rodando a suíte com o
\cd{MockProvider}, que falha fechado.

\textbf{O texto livre de menor volta, com cinco travas.} A conversa com Amara é
o único campo aberto da missão: 600 caracteres, teto de 3 turnos, gravado em
\cd{exp\_07\_mensagens} sob RLS, fora de log e fora da telemetria (só a
contagem de turnos sai), e sem nenhuma rota de gestor que o alcance. É a mesma
postura do \cd{exp\_03\_mensagens}, com um teto menor.

\textbf{Alternativas descartadas.}
\begin{itemize}
  \item \emph{Manter tudo roteirizado}: é o que as bancas aprovaram, e continua
    sendo o piso do produto (o roteiro existe e é exercido). O que ele não dá é
    a reação ao que \emph{aquela} criança escreveu --- e o pedido do produto era
    exatamente esse.
  \item \emph{IA decidindo o desfecho}: um modelo com um dia ruim mudaria a
    lição da aula. O desfecho é gabarito, e gabarito é servidor (regra 3).
  \item \emph{Amara conversando durante a investigação}: destruiria o
    não-negociável do silêncio cognitivo. A conversa foi movida para depois da
    decisão, onde consolida em vez de entregar.
  \item \emph{Prompts no bundle}: proibido pela regra 2, e pior que publicar o
    gabarito --- o prompt descreve o mecanismo inteiro da missão em português.
\end{itemize}

\textbf{Consequências.} \emph{Ruim}: o pacote de insumos aprovado passa a
divergir do que está em produção em dois pontos (IA em runtime e texto livre),
e as bancas C e D precisam revisitar \cd{PED-14} e \cd{DAD-08}. \emph{Ruim}: a
missão passa a consumir quota de IA de estudante --- estimados 30 mil tokens
por criança contra os 300 mil diários. \emph{Ruim}: um sanitizador a mais para
manter (\arq{regras/sanitizador.ts}), e ele já barrou uma saída real: o mock
ecoava o próprio prompt para dentro da tela. \emph{Bom}: o desfecho, a
pontuação e a fixação pelo administrador continuam determinísticos e testáveis
sem rede.

\textbf{Gatilho de reavaliação.} Um playtest mostrar que a reação escrita ao
vivo confunde mais que o roteiro autoral --- ou que crianças usam a conversa com
Amara para escrever o que não deveriam. Nos dois casos, o roteiro do catálogo
já está no código e volta a ser o caminho único trocando uma linha.


\section{ADR-67 --- A mídia da exp-07 é gerada por IA, preparada pelo professor
e cacheada por contrato}
\label{adr:exp07-midia-gerada}

\begin{especificacao}{Contexto}
O brief de assets da exp-07 pede oito peças de conteúdo autoral: duas gravações
de voz, um podcast masterizado de 3m30 com dois locutores e três fotografias
estilizadas de um portal de notícias fictício. Nenhuma delas existe, e nenhuma
sai de código --- são voz humana e design gráfico.

Sem elas, três telas da missão ficam mudas: o áudio de Amara, o áudio da
moradora que espalha o boato e o episódio que destrava a Aula 2.
\end{especificacao}

\textbf{Decisão.} A mídia é \textbf{gerada pelo Vertex}, publicada no bucket de
assets e servida como arquivo estático. Cinco movimentos:

\begin{enumerate}
  \item \cd{AiProvider} ganha \cd{generateSpeech} (MP3, via Cloud
    Text-to-Speech) e \cd{generateImage} ganha proporção opcional. Os dois são
    aditivos: nada fora de \arq{src/ai} conhece o vendor, e \cd{AiTimbre} é
    neutro (\cd{jovem-feminino}, \cd{idoso-feminino}) para o roteiro descrever
    \emph{quem} fala, não o nome de catálogo de uma voz do Google.
  \item Quem prepara é o \textbf{professor}, no painel da turma. A quota diária
    dele é 2 milhões de tokens contra 300 mil do estudante, gerar leva dezenas
    de segundos, e o preparo é o momento em que ele confere se a missão está
    pronta.
  \item A mídia é do \textbf{contrato}, não da escola:
    \cd{experiencias/\{expId\}/\{arquivo\}}, sem \cd{escola\_id} no caminho.
    Toda a rede ouve a mesma Dona Lúcia --- duas turmas discutindo vozes
    diferentes do mesmo boato não estariam falando da mesma coisa.
  \item \textbf{O print viral é um recorte da foto da reportagem}, não uma
    segunda geração. Mesma cena, mesma luz, mesmo instante: o corte visual
    espelha o corte do texto que o estudante vai descobrir. Duas gerações
    independentes dariam duas fotos parecidas, e o par casado que o brief exige
    seria mentira.
  \item \textbf{A página do portal é HTML, não imagem.} Modelo de imagem erra
    texto longo, e o parágrafo suprimido precisa ser legível, selecionável por
    leitor de tela e com contraste medido. A IA faz a foto; o layout é código.
\end{enumerate}

\textbf{Alternativas descartadas.}
\begin{itemize}
  \item \emph{Gerar sob demanda para cada estudante}: cada criança ouviria uma
    Dona Lúcia diferente, o \cd{intervention\_latency} mediria a latência do
    modelo, e o custo se multiplicaria por dez mil.
  \item \emph{Script one-shot com os arquivos no git}: funciona, mas cada
    contrato novo dependeria de alguém rodar o script --- e o precedente é caro
    (a exp-02 subiu com o bucket vazio e a missão ficou impossível de concluir,
    sem nada falhar no deploy nem no log). O script existe
    (\arq{apps/api/scripts/smoke-midia-exp-07.ts}), como ferramenta de
    conferência.
  \item \emph{Print inteiro gerado por IA}: manchete e parágrafos distorcidos, o
    trecho omitido ilegível, e o texto invisível para leitor de tela.
  \item \emph{WAV ou PCM}: não há ffmpeg neste repositório, e um WAV de 3m30
    passaria de 10\,MB. O \cd{audioEncoding: MP3} resolve na origem.
\end{itemize}

\textbf{Consequências.} \emph{Ruim}: a duração do podcast não bate com o brief
--- o roteiro aprovado rende 1\,min\,58\,s, não 3\,min\,30\,s, e chegar lá
exigiria uma cadência de 11 caracteres por segundo, que não é ritmo de podcast.
O cronograma da Aula 2 ganha folga; a copy é que teria de mudar. \emph{Ruim}:
com ADC de usuário a Cloud TTS exige \cd{x-goog-user-project}, e o
\cd{aiplatform} não --- no mesmo preparo, imagem funciona e voz responde 403 com
uma mensagem que não fala de voz. \emph{Ruim}: uma dependência a mais para
habilitar por projeto (\cd{texttospeech.googleapis.com}). \emph{Bom}: sem mídia
preparada a missão roda igual --- o áudio cai na transcrição, que já está sempre
na tela por acessibilidade, e a foto na ilustração procedural.

\textbf{Gatilho de reavaliação.} A rede contratar locução humana para a trilha,
ou um playtest mostrar que a voz sintética atrapalha a compreensão de quem tem
12 anos. Nos dois casos, o caminho é substituir os arquivos no bucket --- o
código da tela não muda.


\section{ADR-68 --- O clone de WhatsApp da exp-07 é escuro, e as quatro peças
que faltam são compensação marcada}
\label{adr:exp07-clone-escuro}

\begin{especificacao}{Contexto}
Os wireframes e o sheet de copy da exp-07 descrevem a missão em \textbf{modo
claro}: balões verdes e brancos sobre fundo bege, o WhatsApp que a maioria das
pessoas usa. A seção \cd{whatsappmobile} do \cd{@unooria/design-system} 0.1.17
publica \textbf{exclusivamente o tema escuro} --- \cd{WaTheme = 'dark'}, e o CSS
traz \cd{:root, [data-wa-theme=dark]} e nada mais.

Além do tema, o pacote não publica quatro componentes que a missão pede, e três
deles sustentam não-negociáveis de ofício: o selo de encaminhada, o banner do
alfinete, o balão de voz com transcrição e o menu de resposta citada.
\end{especificacao}

\textbf{Decisão.} \textbf{O clone é escuro}, fiel ao que o pacote publica.
Derivar uma paleta clara aqui é proibido pela mesma regra do \cd{unitube}: a
fidelidade é ao que o design system entrega, e uma paleta inventada seria um
WhatsApp que não existe em lugar nenhum. O contraste do escuro
(\cd{\#D3DAE0} sobre \cd{\#101D25}, cerca de 11:1) passa o AA que o
\cd{RNF-ACES-001} exige com folga.

Das quatro peças, \textbf{três saem de composição} com o que o pacote já tem: o
selo é \cd{Icon name="forward"} com \cd{Badge}; o banner do alfinete usa a
superfície de appbar; o menu é \cd{Radio} com o texto no \cd{label}. A quarta
--- o balão de voz --- é a única compensação de verdade, marcada
\cd{(b) COMPENSAÇÃO} no CSS, com pedido registrado ao \cd{storybook-unooria}.

Duas notas de mecânica que custaram tempo: o \cd{PhoneFrame} do pacote
\textbf{não entra} na experiência (quem desenha o aparelho é o hub, pelo
\cd{viewport} do manifesto --- empilhar os dois daria um celular dentro de
outro); e o \cd{ListItem} \textbf{não serve} para as três opções de resposta,
porque o título dele tem \cd{white-space: nowrap} --- truncar o texto faria a
criança escolher sem ler, que é exatamente o comportamento que a missão existe
para desfazer.

\textbf{Alternativas descartadas.}
\begin{itemize}
  \item \emph{Esperar o tema claro no \cd{storybook-unooria}}: bloquearia o
    front inteiro da missão por uma diferença que o RNF de acessibilidade não
    cobra --- o escuro passa o AA melhor que o claro.
  \item \emph{Derivar a paleta clara localmente}: proibido, e é a porta para
    cada experiência ter a sua versão do mesmo clone.
  \item \emph{Publicar os quatro componentes antes de começar}: bloquearia as
    etapas 2 a 4 da missão por trabalho em outro repositório, e três deles nem
    seriam componentes novos --- são composições do que já existe.
  \item \emph{Desfazer o \cd{nowrap} do \cd{ListItem} com dois seletores}: é o
    caminho que a doutrina permite, mas usar o componente errado e depois
    desfazer o CSS dele é pior que usar o certo.
\end{itemize}

\textbf{Consequências.} \emph{Ruim}: o pacote de design aprovado descreve telas
claras, e a missão em produção é escura --- os wireframes precisam de uma nota,
e o QA de design revisitaria \cd{DES-06}. \emph{Ruim}: uma compensação local
aberta, que sai quando o pacote publicar o balão de voz. \emph{Bom}: a ponte
\cd{@edu/ui/ds-whatsapp} existe e a seção deixa de ser \enquote{sem consumidor};
a próxima experiência de mensageria começa do que esta deixou pronto.

\textbf{Gatilho de reavaliação.} O \cd{storybook-unooria} publicar o tema claro
da seção \cd{whatsappmobile} --- aí a decisão volta à mesa, e a troca é de uma
folha de estilo, não de código de tela.

\section{ADR-69 --- A missão vive DENTRO da turma: a sala de aula por liberação}
\label{adr:sala-de-aula}

\begin{especificacao}{Contexto}
O enturmamento estava modelado desde a \adrr{ADR-28} --- \cd{periodos\_letivos},
\cd{turmas} com série, \cd{matriculas} datadas com papel --- e o escopo do
professor já era a matrícula. Faltava a camada de cima: \textbf{nada ligava
\cd{turmas} a \cd{experiencias}}.

O efeito era visível nas três telas que o produto tem:

\begin{itemize}
  \item o hub do estudante listava as missões de uma \textbf{constante compilada
    no \emph{bundle}} (\arq{apps/hub/lib/missoes-do-estudante.ts}), igual para
    toda criança de toda escola --- sem olhar série, turma nem escola. E
    \rota{/exp/:expId} não conferia nada: qualquer estudante logado abria
    qualquer missão digitando a URL;
  \item o professor navegava por ANO. Existe uma página por
    \enquote{ano/unidade}, então as duas turmas da 2ª série do EM dividiam a
    mesma ficha, o mesmo guia e a mesma tela de resultados --- a turma só
    reaparecia num \cd{<select>} no fim;
  \item \enquote{liberar} existia só DENTRO de três das cinco experiências, com
    um mecanismo diferente em cada uma: o \cd{status} da atividade da exp-01, o
    relógio da turma da exp-03 (\adrr{ADR-25}), os \emph{gates} da exp-07. Todos
    respondem \enquote{a aula começou}; nenhum responde \enquote{esta missão
    pertence a esta sala}.
\end{itemize}

É a camada que todo LMS tem entre a turma e o conteúdo --- \emph{resources}
ligados a uma \emph{class} no OneRoster, \emph{deep linking} num contexto no
LTI. A \adrr{ADR-28} descartou o LTI porque esta plataforma não vive dentro de
um LMS; a consequência que faltava tirar é a inversa: então ela precisa ter a
própria camada de atribuição, porque é ela quem faz esse papel.
\end{especificacao}

\begin{decisao}
\textbf{Uma tabela de plataforma, \cd{liberacoes}: turma $\times$ experiência
$\times$ estado.} Sem linha, a missão não abre. \cd{liberada} abre;
\cd{encerrada} fecha a porta e mantém tudo o que a criança fez.

Quem libera é \textbf{quem tem matrícula de professor na turma} --- gestor lê a
sala e não decide a aula. Não há prefixo \cd{exp\_NN\_} porque a tabela serve a
todas as experiências: ela é contrato entre a plataforma e cada uma.

\textbf{A porta é do SERVIDOR}, em \arq{comum/liberacao.ts}: \cd{exigirMissao\-Liberada}
é chamada na entrada de estudante de cada experiência --- \rota{POST /runs} e
\rota{GET /runs/boot}, e o funil de sessão do UnooTube, do unooGPT, da Vila e do
Eco. Esconder o \emph{card} no hub é UX; a regra que vale é esta, pela mesma
razão da regra inviolável nº 3 (correção no servidor).

\textbf{Ela é PRÉ-CONDIÇÃO dos relógios e \emph{gates}, não substituta.} Eles
continuam dizendo \enquote{a aula começou} e continuam morando na experiência.
Sem liberação, a pergunta do relógio nem chega a ser feita.
\end{decisao}

\textbf{Três escolhas de tela, e nenhuma é neutra.}

\begin{itemize}
  \item \textbf{Missão bloqueada aparece com CADEADO}, não some. Sumir era mais
    simples e foi recusado: a grade do ano é a promessa do que vem depois, e uma
    lista que muda de tamanho sozinha não deixa a criança entender que a próxima
    unidade existe. \cd{bloqueada} não é estado do banco --- é a ausência de
    linha.
  \item \textbf{Missão encerrada continua VISÍVEL}, com o rótulo. O que a
    criança fez não pode evaporar da tela quando o professor fecha a aula.
  \item \textbf{Professor com uma turma só entra nela clicando}, como quem tem
    quatro. Havia a opção de redirecionar direto, e ela foi recusada: dois
    fluxos diferentes para o mesmo papel é o tipo de atalho que faz a tela de
    lista parecer quebrada para metade da rede.
\end{itemize}

\textbf{Alternativas descartadas.}

\begin{itemize}
  \item \emph{Filtrar só por série} (\cd{experiencias.series}, que já existe e
    ninguém lia): resolve a criança do 4º ano ver a missão do 8º, e não resolve
    nada do resto --- as duas turmas do mesmo ano continuariam idênticas, e o
    professor continuaria sem dizer quando a missão entra.
  \item \emph{Continuar liberando dentro de cada experiência}: é o que existe, e
    o custo está medido --- três mecanismos, e a criança só descobre que a missão
    não abriu DEPOIS de entrar nela, num 409 ou numa tela de espera.
  \item \emph{Esconder no hub e não checar na API}: o \emph{bundle} é público e
    \rota{/exp/:expId} é endereço adivinhável.
  \item \emph{Guardar \enquote{bloqueada} como linha}: obrigaria a criar uma por
    turma $\times$ experiência para o mundo nascer fechado. Ele nasce fechado de
    graça.
  \item \emph{FK de \cd{exp\_id} para \cd{experiencias}}: aquele catálogo é
    escrito pelo CI ao publicar, e em desenvolvimento está vazio --- a FK faria
    liberar depender de um \emph{deploy} que não aconteceu. É a mesma escolha de
    \cd{progresso.exp\_id}.
\end{itemize}

\textbf{Consequências, inclusive as ruins.}

\begin{itemize}
  \item \textbf{Estudante sem turma não vê missão nenhuma}, e isto é mudança de
    comportamento. A exp-03 tinha decisão registrada de que a criança solta via o
    relógio \enquote{aguardando} em vez de um cronômetro só dela; agora ela nem
    entra. Sem sala não há missão --- e quem explica isso é a tela do hub, porque
    um erro dentro do aparelho não teria como dizer \enquote{peça para a
    coordenação te matricular} sem sair do disfarce.
  \item \textbf{O professor ganhou um passo antes da aula.} Missão que ninguém
    liberou não abre, e a sala precisa dizer isso com todas as letras --- é o
    risco novo do capítulo 19.
  \item \textbf{O \emph{seed} passou a escrever \cd{experiencias}.} Em produção
    quem escreve é o CI ao publicar; sem catálogo não há grade, e o professor
    abriria a sala sem nada para liberar.
  \item A checagem custa três consultas indexadas na entrada de cada
    experiência. Ficam na CRIAÇÃO da sessão, não em toda rota: quem já entrou
    fica, e encerrar fecha a entrada em vez de expulsar a criança no meio da
    aula.
\end{itemize}

\textbf{Gatilho de reavaliação.} Uma rede querer liberar em lote pela
coordenação (hoje gestor não libera: é uma linha no \cd{@Roles} e uma frase na
matriz do capítulo 7), ou pedir \textbf{janela de datas} --- \enquote{abre
segunda, fecha sexta} ---, que hoje é um clique manual e viraria duas colunas e
uma rotina que as cumpra. Enquanto não houver quem cumpra, a coluna não nasce:
é a armadilha do \cd{expires\_at}.

\section{ADR-70 --- As páginas legais descrevem o sistema que existe, e a rede é
a controladora}
\label{adr:paginas-legais}

\begin{especificacao}{Contexto}
O rodapé da landing oferecia \enquote{Privacidade} e \enquote{Termos} com
\cd{href="\#"} --- dois links que não levavam a lugar nenhum. Escrever as duas
páginas expôs uma pergunta que não é de redação: \textbf{o que elas podem
afirmar?}

Uma auditoria do repositório inventariou o tratamento de dados a partir do
código, não do desejo. Ela encontrou vinte afirmações que uma política de
privacidade padrão faria e que \textbf{seriam falsas aqui}: não existe entidade,
coluna ou tela de responsável legal; não existe rota de exclusão de pessoa
(desligar alguém é \cd{ativo = false}); não existe \textbf{nenhuma} rotina de
retenção no repositório; a stack está em \cd{us-east1} e a IA no endpoint
\cd{global} do Vertex; \cd{users.cpf} é \cd{TEXT} em claro (ADR-33); a
telemetria é pseudonimizada, não anônima; e não há log de acesso individual ---
ninguém sabe hoje qual gestor abriu o CPF de qual criança.

A pesquisa jurídica acrescentou o ponto que decide a forma do documento: o
cliente é uma \textbf{rede pública de ensino}, e o Guia do Poder Público da ANPD
usa como exemplo literal a matrícula de estudante para concluir que o
consentimento obtido em relação compulsória é \textbf{nulo} (Enunciado CD/ANPD
nº 1/2023).
\end{especificacao}

\textbf{Decisão.} \textbf{As páginas descrevem o sistema que existe.} Toda frase
tem lastro em tabela, rota ou envio verificado, e o que hoje não existe é dito
como não existindo --- inclusive quando isso é desconfortável. Onde faltava
número (prazo de retenção), o texto atribui a definição à rede em vez de
inventar meses; onde faltava processo (revisão humana de decisão automatizada),
o texto não promete o processo.

\textbf{Não há pedido de consentimento.} A base é o art. 7º, III da LGPD ---
política pública ---, e a página tem uma seção inteira explicando \emph{por que}
não existe botão de aceitar, para que a ausência não seja lida como descuido. O
único consentimento legítimo do produto é o do professor no DASS-21
(art. 11, I), e ele está isolado ali. \textbf{A rede é a controladora; a Unooria
é operadora} --- com uma exceção declarada: nos registros de acesso do art. 15
do Marco Civil, a controladora é a Unooria.

\textbf{Os Termos não elegem foro} (Lei 14.133/2021, art. 92, parágrafo 1º manda o foro
da sede da Administração no contrato, e cláusula de foro contra titular menor é
abusiva) \textbf{e não limitam responsabilidade de forma ampla} (art. 51 do CDC).

Duas decisões de forma, que são de conteúdo disfarçado. \textbf{Zero
\cd{<table>}}: texto jurídico numa tabela de quatro colunas produz uma coluna de
três palavras e outra de quarenta, célula altíssima e rolagem lateral no
celular --- que é onde a família abre o link. O mesmo dado vira \cd{legal-card},
que cresce com o conteúdo e reflui sozinho. E \textbf{nenhum \cd{[A DEFINIR]}
espalhado pelo texto}: o que depende do contrato mora num bloco
\cd{IDENTIFICACAO} no topo de cada arquivo, e o que ainda não foi decidido tem
uma seção final própria. Razão social e CNPJ ficam \textbf{vazios} até virem do
cadastro --- um CNPJ de exemplo numa página pública é declaração falsa que passa
na revisão justamente por parecer preenchida.

\textbf{Alternativas descartadas.}
\begin{itemize}
  \item \emph{Adaptar um modelo de política de mercado}: todo modelo pressupõe
    consentimento do titular e autocadastro. Aqui os dois são falsos, e o
    consentimento seria nulo --- o modelo erra justamente no que é específico.
  \item \emph{Prometer o estado desejado e implementar depois}: promessa vincula
    (LGPD art. 6º, X). Seria a mesma armadilha do \cd{expires\_at} que saiu do
    schema e ficou no registro: a regra sai do código, o texto fica, e quem lê
    acredita nele.
  \item \emph{Publicar só a Política e deixar os Termos para depois}: quem lê uma
    procura a outra, e metade dos pontos (simulações, IA, conduta) não cabe na
    Política.
  \item \emph{Marcar \cd{[A DEFINIR]} em cada lacuna do texto}: foi tentado e
    rejeitado --- o documento fica ilegível e a lacuna deixa de ser vista de
    tanto se repetir.
  \item \emph{Uma casca visual própria para as páginas legais}: elas são a mesma
    superfície pública da landing, alcançadas pelo rodapé dela. Casca própria
    faria a página parecer de outro produto na hora em que a confiança está
    sendo decidida.
\end{itemize}

\textbf{Consequências.} \emph{Ruim}: as páginas declaram em público limitações
reais do produto (sem exclusão por titular, sem retenção, sem log de acesso,
processamento de IA fora do Brasil) --- é material que uma rede pode usar em
negociação, e o incentivo a fechar essas lacunas passa a ser externo.
\emph{Bom, por tabela}: escrever a página encontrou duas contradições que o
texto não podia encobrir. A do \arq{docs/INFRA.md} \textbf{foi corrigida
aqui} --- ele afirmava que o Vertex fica em \cd{southamerica-east1} e, mais
abaixo, que a localização é \cd{global}; o código concorda com a segunda em
quatro lugares (\arq{config/env.ts}, os dois \emph{compose} e o
\cd{LOCALIZACAO\_PADRAO} do provider), e o \cd{3.6-flash} responde 404 na
região brasileira. Era a linha que uma política copiaria para afirmar que os
dados ficam no Brasil. \emph{Ruim}: a outra é de produto e segue aberta --- a
exp-91 exibe \enquote{Só você lê o que escreve aqui} sob a caixa cujo conteúdo
vai ao Vertex para gerar o acolhimento. \emph{Bom}: o inventário existe, e com ele o RIPD que a Resolução
CD/ANPD nº 2/2022 torna exigível deixa de ser trabalho do zero.

\textbf{Gatilho de reavaliação.} Qualquer uma das lacunas ser fechada no código
--- rotina de retenção, rota de exclusão por titular, log de acesso individual,
região brasileira para o Vertex ---, o que muda o que a página pode afirmar; ou
a primeira rede assinar, quando a controladora passa a ter nome e CNPJ.

\section{ADR-71 --- O trabalho paralelo se reparte por fatia vertical, e o
runtime continua singleton}
\label{adr:fatia-vertical}

\begin{especificacao}{Contexto}
A \adrr{ADR-61} resolveu a \emph{árvore de trabalho}: cada frente paralela ganha
um \cd{git worktree}, e o \cd{git status} de uma não contamina o commit da
outra. Ela não resolveu a pergunta seguinte, que aparece no minuto em que existe
mais de um terminal aberto: \textbf{\arq{apps/api} e \arq{apps/hub} são comuns
às 48 experiências}. Worktree separa arquivos não commitados; não separa duas
frentes que precisam editar o mesmo arquivo.

A divisão que o instinto sugere é por camada --- uma frente no \emph{backend},
outra no \emph{frontend}. Ela é a pior possível aqui: as duas passam a morar nos
mesmos arquivos por motivos diferentes, e o conflito não é textual (duas linhas
vizinhas), é semântico (dois desenhos do mesmo \emph{service}). Nenhuma
ferramenta de merge ajuda com isso.

O segundo eixo é o tempo. O runtime deste repositório é \textbf{singleton} por
decisão: portas com \cd{strictPort}, um \cd{docker-compose} de nome fixo
(\cd{unooria-plataforma}), um Postgres com um schema. Isso não é descuido --- é
o que faz \cd{pnpm dev} falhar alto em vez de servir a experiência errada dentro
do \emph{iframe}, que é uma falha silenciosa e cara de diagnosticar.
\end{especificacao}

\textbf{Decisão.} \textbf{Reparta por fatia vertical.} Cada frente leva o módulo
da api, a rota do hub e a experiência \emph{juntos} --- \arq{apps/api/src/unootube/}
com a rota do hub e \arq{apps/experiencias/exp-02-unootube/}, e não "a api" de um
lado e "o hub" do outro. Com isso api e hub deixam de ser recurso disputado: a
fatia de cada frente dentro deles é uma \textbf{pasta própria}. Não é sorte --- é
o desenho de um módulo por domínio (\cd{catalogo.ts} + \cd{regras.ts} + service +
controllers por papel) sendo usado para o que ele serve. \textbf{Duas frentes no
mesmo módulo não se paralelizam: serializam.}

\textbf{O que sobra compartilhado é linha de REGISTRO, não lógica}, e vai em UM
commit no fim da frente, depois de \cd{git merge origin/dev}: um \cd{import} e
uma linha no array de \arq{app.module.ts}; \cd{model} novo \textbf{no fim} do
\arq{schema.prisma} (duas inserções no fim conflitam limpo, duas no meio não); a
migração; a entrada em \arq{middleware.ts} / \arq{experiencias-dev.ts}. A porta
de dev se reserva como o número da ADR: conferindo \cd{origin/dev}.

\textbf{O runtime continua singleton}, e o trabalho se divide em duas fases. A
\emph{paralela}, N frentes: editar, \cd{typecheck lint test}, \cd{build --filter},
e o Vite de UMA experiência em porta \cd{52xx} embedada por
\cd{/dev-embed?src=}. A \emph{serial}, uma por vez na cópia principal:
\cd{pnpm dev}/\cd{docker:dev}, migração, e a suíte de ponta a ponta.

\textbf{E a fatia é BARREIRA, não conselho.} \arq{scripts/worktree.mjs}
(\cd{pnpm worktree criar/listar/remover}) faz o caminho inteiro --- worktree
sobre \cd{origin/dev}, os seis \cd{.env} copiados (todos ignorados pelo git, nenhum
viaja), \cd{pnpm install} e \cd{prisma generate} (\arq{node\_modules} e o client
do Prisma são por worktree), porta derivada de um \emph{hash} do slug --- e
\textbf{resolve o \cd{--fatia} contra a árvore} (\cd{api/unootube} vira
\arq{apps/api/src/unootube}; \cd{hub/professor} acha o grupo de rota do App
Router) antes de gravar a lista em \cd{<git-dir>/fatia}. Num worktree esse
caminho é \arq{.git/worktrees/<slug>/}: privado daquela frente, fora do
versionamento, e portanto sem conflito possível entre frentes.

O \arq{.githooks/pre-commit} lê essa lista e \textbf{recusa} o que estiver fora
dela, com o escape nomeado \cd{PERMITIR\_FORA\_DA\_FATIA=1} --- o mesmo desenho
dos outros hooks. Sem o arquivo, o hook é inerte, e a cópia principal segue
livre. A costura passa, mas passa \emph{avisando}: o lugar dela é um commit
separado no fim. O efeito colateral que vale mais que o hook: a lista de
costuras deixa de existir só em prosa e passa a existir como código, num lugar
onde ficar desatualizada quebra alguém em vez de enganar em silêncio.

\textbf{Alternativas descartadas.}
\begin{itemize}
  \item \emph{Repartir por camada (uma frente na api, outra no hub)}: é a origem
    do problema, não a solução. Some com a fronteira que os módulos por domínio
    construíram e transforma todo merge em releitura de desenho.
  \item \emph{Stack por worktree --- offset de portas, \cd{name} do compose por
    slot, um banco por frente}: tornaria a fase serial paralela também, e foi a
    alternativa mais séria. Descartada por três razões. O \cd{strictPort} existe
    para \textbf{falhar alto}: com portas móveis, o \emph{iframe} volta a poder
    apontar para a experiência errada em silêncio, que é exatamente o bug que a
    porta fixa previne. N bancos divergem --- a suíte e2e é a etapa que mais
    depende de um estado conhecido, e depurar "passou no meu banco" custa mais
    que a fila. E o custo de manutenção (compose parametrizado, envs por slot,
    \arq{experiencias-dev.ts} dinâmico, a suíte e2e ciente do slot) cai sobre
    \textbf{uma pessoa}, para destravar uma fila que hoje tem dois ou três itens.
  \item \emph{Banco por frente, mantendo as portas fixas}: metade do custo acima
    e menos da metade do ganho --- continua sem poder rodar duas stacks, que é o
    que a fila realmente pede.
  \item \emph{Deixar a mecânica só como prosa no \arq{CONTRIBUTING.md}}: era o
    estado anterior. O passo que morde é o \textbf{primeiro}, e é o que a prosa
    não consegue impor --- \cd{git worktree add} sem \cd{origin/dev} escrito abre
    a branch sobre produção em silêncio (\adrr{ADR-61}).
  \item \emph{Barrar na ESCRITA, pelas permissões do agente} (um
    \arq{.claude/settings.local.json} por worktree negando \cd{Edit}/\cd{Write}
    fora da fatia): pegaria o erro antes do trabalho ser feito, e não no commit.
    Descartada por cobertura e por forma. Cobre \textbf{um} agente --- outra
    ferramenta, ou a edição à mão, passa direto ---, enquanto o hook cobre
    qualquer coisa que produza um commit. E \emph{glob} não expressa "tudo menos
    X": negar o vizinho exigiria listar os trinta e um módulos da api um a um,
    regerando a lista toda vez que nasce um módulo. Fica como camada extra
    possível no dia em que o custo do trabalho jogado fora superar o do commit
    recusado.
  \item \emph{Registro de portas num arquivo versionado}: estado a mais para
    manter sincronizado, e conflito num arquivo que TODA frente tocaria. O
    \emph{hash} do slug dá uma porta estável sem guardar nada --- e estável
    importa: o \cd{src=} do \cd{/dev-embed} fica na barra do navegador de quem
    está trabalhando, e não pode mudar porque um vizinho abriu uma frente.
\end{itemize}

\textbf{Consequências.} \emph{Ruim}: o teto útil é de duas a três frentes ---
toda frente que mexe em tela precisa da suíte e2e, que é serial, e além disso a
fila come o ganho. \emph{Ruim}: o \cd{pre-commit} pega no \textbf{commit}, não na
escrita --- o agente ainda gasta trabalho num arquivo que não devia tocar, e
descobre ao commitar --- e a alternativa que pegaria antes, descartada acima,
cobre um agente só. \emph{Ruim}: hook é conselho para quem clona e não roda
\cd{pnpm install} --- sem \cd{core.hooksPath}, não há barreira nenhuma, e isso
vale para os três hooks. \emph{Ruim}: a fatia não reprova caminho inexistente,
porque a frente pode estar criando o módulo; um erro de digitação vira uma fatia
que protege menos do que parece --- por isso o script \emph{diz} qual caminho
não existe. \emph{Bom}: a classe de erro da \adrr{ADR-61} morre na origem,
porque o \cd{origin/dev} passa a estar escrito no script em vez de na memória de
quem digita. \emph{Bom}: o script recusa nome de branch que o \cd{pre-push}
recusaria, e a recusa acontece no minuto zero em vez de depois de horas de
trabalho numa branch mal nomeada.

\textbf{Gatilho de reavaliação.} Passar a existir mais de uma pessoa
desenvolvendo ao mesmo tempo (hoje é uma), quando a fila da fase serial deixa de
ser aceitável; ou a suíte de ponta a ponta virar o gargalo medido --- aí a stack
por worktree volta à mesa com o custo justificado por número, não por intuição.

% ---------------------------------------------------------------------------
\section{ADR-72 --- A fronteira do disfarce é a revelação, e a Etapa 4 da
exp-06 fica fora dele}
\label{adr:72}

\textbf{Data.} 2026-09-06 \quad \textbf{Status.} Aceita

\textbf{Contexto.} A exp-06 termina com duas telas analíticas: um gráfico
comparando a velocidade do boato com a da nota oficial, e um mapa da topologia
de rede que desenha o alcance da decisão do estudante. Os wireframes e o
\emph{handoff} da missão colocam as duas \textbf{dentro do clone}, descritas
como \enquote{painel de métricas do grupo} --- diegéticas, faladas pela seção
\cd{whatsappmobile}.

O problema é que o aplicativo real não tem nada parecido. Ele mostra
\enquote{Encaminhada com frequência} num balão; não mostra grafo de propagação,
não conta conexões e não desenha nós. Uma tela dessas dentro do disfarce não é
verossimilhança seletiva --- que admite \textbf{ausências} (sem teclado, sem
chamadas, sem Stories) --- é \textbf{invenção}: uma funcionalidade que o app
imitado não possui, apresentada como se possuísse.

\textbf{Decisão.} O disfarce vale \textbf{até o fim da Etapa~3}. Da pausa entre
as aulas em diante quem fala é a escola: o card de espera, o gráfico, o mapa, o
construtor do ciclo e o \enquote{Meu percurso} usam \cd{@edu/ui/ds} e os tokens
\cd{--unooria-*}.

O critério que decide, e que já valia no repositório, é
\textbf{\enquote{quem está falando?}} --- não \enquote{é depois no tempo?}. A
exp-02 já o havia aplicado: feed, cromo e playlists falam de dentro da ficção; o
ranking e os diálogos da missão falam pela escola, ainda que apareçam no meio
dela.

\textbf{Alternativas descartadas.}
\begin{itemize}
  \item \emph{Seguir os wireframes e desenhar o painel dentro do clone}: era o
    caminho especificado, e o QA de design o aprovou. Descartada porque quebra o
    próprio teste de fidelidade que aquele QA aplicou (\enquote{o frame passaria
    por uma captura do app real?}): um grafo de topologia não passa, porque o
    app não tem grafo de topologia. E a fidelidade do clone é o instrumento ---
    é ela que faz o comportamento medido ser espontâneo.
  \item \emph{Levar o pedido ao \arq{storybook-unooria}}, publicando um
    componente de visualização de rede na seção \cd{whatsappmobile}: pioraria o
    problema em vez de resolvê-lo. O pacote passaria a publicar, como parte do
    clone, uma tela que o original não tem --- e a próxima missão que usasse a
    seção herdaria a invenção.
  \item \emph{Manter as duas telas no clone e apenas trocar os tokens}: é o pior
    dos dois mundos. A tela continuaria dentro do aparelho, com o cromo do
    mensageiro em volta, e a marca institucional apareceria no meio do disfarce
    --- a mistura de seções que a regra de ouro proíbe.
  \item \emph{Tirar as telas analíticas da missão} e deixar a revelação apenas
    para o debate mediado: mataria o beat memorável. A descoberta central da
    missão é a diferença entre 1.200 e 850 conexões, e ela precisa ser
    \textbf{vista}, não contada --- é isso que a criança leva.
\end{itemize}

\textbf{Consequências.} \emph{Ruim}: a missão troca de identidade visual no
meio, e a transição precisa ser justificada ao estudante pelo próprio roteiro
--- o card de pausa é o que faz esse trabalho (\enquote{sua investigação está
salva}), e sem ele a troca pareceria um erro de carregamento. \emph{Ruim}: o
insumo aprovado e o código divergem neste ponto, e a divergência tem de viver
escrita em três lugares (aqui, no \arq{AGENTS.md} da experiência e no
\arq{catalogo.ts} do hub) para não ser \enquote{corrigida} de volta por quem
chegar depois lendo só os wireframes. \emph{Bom}: o teste do app sem legenda
continua valendo para todas as telas que se apresentam como o mensageiro.
\emph{Bom}: os tokens da Etapa~4 passam a ser os da plataforma, e isso corrigiu
um defeito real --- com os tokens do clone, calibrados para fundo escuro, os nós
do mapa saíam quase invisíveis sobre o branco e o beat memorável não chegava.

\textbf{Gatilho de reavaliação.} O dia em que o aplicativo imitado publicar de
fato uma visualização de propagação (o \enquote{encaminhada com frequência} já
foi nessa direção). Aí a tela deixa de ser invenção, e a fronteira pode voltar
para depois dela.


% --------------------------------------------------------------------------

\section{ADR-73 --- O planejamento de ensino é dirigido pela BNCC extraída do
PDF homologado, e o componente da rede não é o componente da Base}
\label{adr:plano-de-ensino-bncc}

\begin{especificacao}{Contexto}
A exp-90 gera plano anual a partir do \textbf{currículo da rede importado como
texto} (\adrr{ADR-47}): um documento em prosa entra no \emph{prompt}, e o modelo
infere dele a progressão do ano. O arranjo serve a quem tem documento curricular
próprio e material didático adotado --- e é inútil para quem não tem, que é a
maioria das redes pequenas. Ele também não desce da intenção anual para a aula
de quinta-feira, que é onde o trabalho do professor realmente acontece.

Havia um segundo limite, mais estreito: o produto nasceu para a
\textbf{BNCC Computação}. Mas a Resolução CNE/CEB nº 1/2022 não cria componente
isolado --- ela \textbf{delega às redes} definir \enquote{parâmetros e abordagens
pedagógicas} (art. 2º) ---, e a Lei 14.533/2023 fez da educação digital
componente do Fundamental e do Médio pela LDB (art. 26, parágrafo 11). Computação
é camada transversal, não disciplina; um planejador que só a atenda serve a uma
fração da carga horária de um professor.

Três achados do levantamento decidiram a forma da solução, e nenhum deles é
óbvio:

\textbf{Primeiro: no Ensino Médio, componente da matriz e componente da Base
NÃO coincidem.} A matriz de qualquer rede tem Biologia, Física e Química; a BNCC
do Médio desce apenas até a \textbf{área}, e as três compartilham as 26
habilidades de Ciências da Natureza e suas Tecnologias (\cd{EM13CNT}). O mesmo
vale para Filosofia, Geografia, História e Sociologia, que dividem as de Ciências
Humanas. As duas exceções são Língua Portuguesa e Matemática, os únicos
componentes que a Base do Médio nomeia.

\textbf{Segundo: a Base não é uniforme no organizador acima do objeto de
conhecimento.} Sete componentes usam unidade temática; Língua Portuguesa usa
prática de linguagem dentro de campo de atuação; Língua Inglesa usa eixo; o
Médio agrupa por competência específica de área.

\textbf{Terceiro: o PDF que o MEC serve na página do Ensino Médio é o RASCUNHO}
de 11/05/2018, com 179 habilidades contra as 183 do homologado e
\textbf{textos divergentes} --- a \cd{EM13MAT101} do rascunho perdeu
\enquote{criticamente} e \enquote{fatos relativos}. O documento consolidado de
dezembro de 2018, distribuído sob o nome \cd{BNCC\_EI\_EF\_110518}, é o que traz
as três etapas.
\end{especificacao}

\textbf{Decisão. Nasce a exp-92, ao lado da exp-90 e sem substituí-la}, com o
currículo como DADO ESTRUTURADO extraído do PDF homologado, cobrindo todos os
componentes do 1º ano do Fundamental à 3ª série do Médio, em dois níveis: plano
\textbf{anual} (macro, por período de avaliação) e plano de \textbf{período}
(granular, aula a aula, semanal ou quinzenal à escolha da rede).

\textbf{A fonte é o PDF homologado, extraído por coordenadas.} As tabelas da
Base são um \emph{spread} de duas páginas --- unidade temática e objeto de
conhecimento na par, habilidades na ímpar --- e no texto corrido do
\cd{pdftotext} as três colunas viram três blocos sequenciais: a ligação
habilidade\,$\leftrightarrow$\,objeto de conhecimento se perde, e atribuí-la
seria chute. O que salva é o alinhamento do InDesign: o objeto
\enquote{Estratégia de leitura} sai em \cd{yMin=260,5} e a \cd{EF15LP02} em
\cd{yMin=260,4}, na página seguinte. A associação é geométrica e determinística
(\cd{scripts/extrair-bncc.mjs}), e rende \textbf{1.304} habilidades do
Fundamental --- o número oficial ---, \textbf{183} do Médio e \textbf{129} de
Computação, das quais 98\,\% com objeto de conhecimento.

\textbf{O componente da rede é modelado separado do componente da Base}
(\cd{componenteBncc} em \cd{curriculo/componentes.ts}), e a tela AVISA quando as
habilidades são da área. Sem isso, o professor de Biologia do 2º ano --- o caso
mais comum do Médio --- receberia \enquote{nenhuma habilidade encontrada}.

\textbf{A distribuição de habilidades pelos períodos é DECISÃO HUMANA.} Ela
expressa a leitura que o professor faz da própria turma, e é a única coisa aqui
que nenhum modelo deveria escolher por ele; a IA entra depois, para desenvolver
pedagogicamente o que ele decidiu. Pelo mesmo motivo, o \textbf{esqueleto das
aulas não vem do modelo}: número, data e horário saem do calendário e da grade
real. Deixar o modelo decidir quantas aulas cabem numa quinzena produzia dez ---
o número redondo --- para um professor de Arte que tem duas.

\textbf{Há três barreiras contra código de habilidade inventado}, e nenhuma
sozinha basta: o \emph{prompt} proíbe, a rota de distribuição recusa código fora
do recorte, e a costura filtra o que o modelo devolveu
(\cd{planejamento/costura.ts}). Um plano que cite \cd{EF67LP42} chega à
coordenação como erro do professor, e destrói a confiança no produto inteiro.

\textbf{A articulação com Computação é rotulada como PROPOSTA.} Não existe
\emph{crosswalk} normativo entre habilidade de Computação e habilidade de
componente --- a norma delega a abordagem à rede ---, e vender inferência como
norma é risco de auditoria numa venda a poder público. Em contrapartida, a
extração aproveita as colunas \enquote{Explicação da habilidade} e
\enquote{Exemplos} do anexo do Parecer CNE/CEB nº 2/2022: texto pedagógico
oficial que nenhum \emph{dataset} público da Base empacota.

\textbf{Alternativas descartadas.}

\textbf{Consumir um \emph{dataset} público de terceiro} (o \cd{bncc-dev/bncc-dados},
CC BY 4.0, com DOI e proveniência auditável) foi seriamente considerado e é
tecnicamente bom. Recusado por dois motivos: num contrato B2G a pergunta
\enquote{de onde veio este texto?} precisa terminar no documento oficial, não
num repositório de terceiro; e o mantenedor é uma \emph{edtech}, o que faz da
base curricular do produto uma dependência de um concorrente potencial. Ele
continua útil como conferência de completude no CI.

\textbf{Ingerir o Mapa de Foco do Instituto Reúna} (priorização de habilidades
por ano, pedagogicamente excelente) foi \textbf{vedado por licença}: os Termos
de Uso proíbem alteração e comercialização sem autorização expressa. Se
interessar, entra por contrato, não por clone de repositório.

\textbf{Um sistema de RAG sobre os PDFs} foi descartado por ser a ferramenta
errada, não por custo. A pergunta que o produto faz é sempre \enquote{quais
habilidades de \emph{componente} no \emph{ano}} --- consulta exata sobre dois
campos estruturados. Busca vetorial responderia com aproximação onde existe
resposta certa, e um plano anual que \enquote{esqueceu} a \cd{EF67LP08} porque
ela ficou fora do \emph{top-k} é um plano que descumpre o currículo.

\textbf{Substituir a exp-90} foi descartado: são ferramentas para redes
diferentes, e migrar os planos já salvos custaria mais do que manter as duas.

\textbf{Consequências ruins, declaradas.}

\textbf{Ficam duas telas de planejamento no menu do professor}, e a diferença
entre elas não é auto-evidente para quem chega. A mitigação é a descrição de
cada card no hub; se a confusão aparecer em campo, uma delas some.

\textbf{A matriz de carga horária é mediana de dez redes, não norma --- e a
mediana esconde uma bimodalidade.} A BNCC não define grade horária (LDB art. 23
e 26; Parecer CNE/CEB nº 8/2004), e o total semanal do Fundamental tem duas
famílias sem nada no meio: seis redes em 25 aulas (833\,h) e quatro em 28--30
(890--1.000\,h, entre elas as duas maiores do país). Uma média entre as duas
descreveria rede nenhuma, então o produto pré-preenche o subtotal da base comum
e deixa o professor ajustar. Nos anos iniciais há um agravante: em várias redes o
modelo nem é disciplinar --- São Paulo e o município do Rio agregam cinco
componentes num bloco de professor polivalente, e para elas a matriz sugerida
não descreve a realidade.

\textbf{Vinte e seis habilidades do Fundamental (2\,\%) ficaram sem objeto de
conhecimento} na extração, e aparecem num grupo \enquote{Sem classificação}.
Preferimos mostrá-las a escondê-las: habilidade escondida é currículo não
cumprido.

\textbf{A extração depende de \cd{pdftotext} e de um PDF que o MEC pode
substituir.} Se o layout mudar, o extrator quebra --- e quebra em silêncio, o
que é pior. Contra isso há testes que travam o DADO, não só o código: os totais
oficiais de cada documento, e o texto de \cd{EM13MAT101} com as palavras que só
existem na versão homologada.

\textbf{Gatilho de reavaliação.} Uma revisão da BNCC (há discussão pública sobre
a do Ensino Médio desde a Lei 14.945/2024); a publicação de uma exportação
oficial máquina-legível estável pelo MEC; ou o primeiro contrato cujo currículo
próprio precise ser exibido AO LADO da Base --- caso em que a fronteira entre
exp-90 e exp-92 deixa de fazer sentido e as duas convergem.